% NEG_NIGHT GCT track. Authoring run /root/gct.
% Proof checks and source-reading scope are recorded separately, not inferred
% from the presence of a theorem environment.
\section{What the GCT negations say, and what arithmetic realizations retain}
\label{sec:nn-gct}

Throughout this section, E and A have the status assigned in the prompt.
No theorem below uses M except where this is expressly indicated. A finding
is new to this run; no historical-priority claim is intended. The statements
P-F31, P-F56, P-F57, PLE-F5 and PLE-F8 are inputs from the supplied
ALPOGE\_GCT return. The rank-three multiplication table is the supplied
integral table with the corrected sign of $b_2$, not a newly computed
rank-four canonical basis.

\subsection{NN-G001: the arithmetic zeta of the actual $D_5$ plethysm algebra}
\paragraph{Statement.}
Let
\[
 X=\operatorname{Spec}\mathbb Z[U,V,Z]/(Z^2-UV^2-4U^4),
 \qquad \deg(U,V,Z)=(2,3,4).
\]
For every finite field $\mathbb F_q$, including characteristic two,
\[
 |X(\mathbb F_q)|=q^2,\qquad
 Z(X_{\mathbb F_q},T)=(1-q^2T)^{-1},\qquad
 \zeta_X(s)=\zeta(s-2).
 \tag{G001.1}
\]
The grading nevertheless has Hilbert series
\[
 H_X(t)=\frac{1+t^4}{(1-t^2)(1-t^3)},\qquad
 \dim X_s=\lfloor s/3\rfloor+\mathbf1_{2\mid s},
 \tag{G001.2}
\]
where $X_s$ in this formula means the degree-$s$ coordinate space. Its
weight support is $S=\{0,2,3,4,\ldots\}$, its saturation in its group
$\mathbb Z$ is $\mathbb N$, and its conductor is $2$. Thus the ungraded
arithmetic zeta agrees with that of $\mathbb A^2_{\mathbb Z}$ while the
actual plethysm grading has a hole and has no homogeneous rational
Ehrhart model. Under A the nontrivial zeros in (G001.1) lie on
$\operatorname{Re}s=5/2$; this supplies no information recovering that
hole or an Ehrhart model.

\paragraph{Inputs and locators.}
ALPOGE\_GCT, PLE-F5 and PLE-F8, respectively
\texttt{work/plethysm.tex}, the subsection beginning ``The actual
plethysm ray'', and \texttt{work/positivity\_actual\_plethysm.tex}:
the complex graded algebra is the highest-weight multiplicity algebra
for $S^{(7s,5s)}\mathbb C^2$ in
$\operatorname{Sym}^{4s}(\operatorname{Sym}^3\mathbb C^2)$.
Kahle--Micha\l{}ek, \emph{Obstructions to combinatorial formulas for
plethysm}, Theorem 5, supplies the non-Ehrhart assertion. A is used
only for the final zero-location assertion.

\paragraph{Proof.}
Suppose first that $q$ is odd and let $\chi$ be its quadratic character,
extended by $\chi(0)=0$. For $U=0$ there are $q$ points: $V$ is arbitrary
and $Z=0$. For $U=u\ne0$ the number is
\[
 q+\sum_{v\in\mathbb F_q}\chi(uv^2+4u^4)
 =q+\chi(u)\sum_v\chi(v^2+4u^3)=q-\chi(u).
\]
Here, for every $a\ne0$, the equation $y^2-x^2=a$ has $q-1$
solutions: its two factors $y-x,y+x$ are a freely chosen nonzero
element and its prescribed reciprocal multiple. It also has
$q+\sum_x\chi(x^2+a)$ solutions. Therefore that character sum is
$-1$. Summing $q-\chi(u)$ over $u\ne0$ and using
$\sum_{u\ne0}\chi(u)=0$ proves $q^2$.
If $q$ is even, squaring is a bijection of $\mathbb F_q$, so each
pair $(U,V)$ admits exactly one $Z$. This proves the count for all
finite fields, hence for every extension of each residue field.
The defining exponential formula for local zeta functions gives
$\exp(\sum_{r\ge1}q^{2r}T^r/r)=(1-q^2T)^{-1}$, and the product
over rational primes gives (G001.1).
The relation is monic in $Z$, so $1,Z$ is a free basis over
$\mathbb Z[U,V]$. This proves flatness over $\mathbb Z$ and the
displayed Hilbert series. Its coefficient formula is the supplied
PLE-F5 calculation; directly it counts $U^iV^jZ^\epsilon$ with
$\epsilon=0,1$. Degrees $2,3$ generate precisely $S$; $1$ is absent,
$2\cdot1\in S$, and every integer at least two is present. The
remaining conclusions now follow from the stated inputs.\hfill$\square$

\paragraph{Algebraic geometry and Grothendieck dictionary.}
This is an actual flat affine scheme over $\mathbb Z$, together with a
$\mathbb G_m$ action. The action, equivalently its grading, is the
extra datum in SGA 3, I, Proposition 4.7.3 and Corollary 4.7.3.1.
Forgetting it loses the Hilbert function. SGA 5, III, 4.7 is the
Lefschetz--Verdier trace framework for the local zeta; equality of
zeta functions alone is not an identification of cohomological
degrees, motives, or canonical bases. EGA II, Proposition 2.4.7(i),
additionally shows that taking a Veronese subring preserves
$\operatorname{Proj}$ as a scheme while changing the recorded grading.
There is no implication from the normality of this complex $D_5$
surface to saturation of the support of an arbitrary torus grading:
the example itself refutes that inference.

\paragraph{Complexity impact.}
An arithmetic realization with the same zeta as affine space cannot
by itself establish PH1, PH2, PH3, or a positivity algorithm. This is
an explicit obstruction to the proposed recovery from the classical
zeta sheet, rather than a complexity-class separation.

\subsection{NN-G002: the exact output obstruction from the exponential poles}
\paragraph{Statement and proof.}
On the P-F31 cubic family, put $M=2^{n-1}$ and let $f_n$ be its
stretching quasi-polynomial. A full constituent table for $f_n$ has
at least $M$ entries. Every denominator presentation
\[
 H_n(t)=Q(t)\prod_j(1-t^{a_j})^{-b_j},\qquad a_j,b_j>0,
 \tag{G002.1}
\]
has some $a_j$ divisible by $M$; its expanded denominator has degree
at least $M$. Indeed P-F31 gives a genuine pole at every primitive
$M$th root. At one such root a denominator factor in (G002.1) must
vanish, which is exactly $M\mid a_j$. A quasi-polynomial of period
$L$ has a generating function whose poles are $L$th roots of unity:
sum its polynomial constituents over arithmetic progressions and
differentiate geometric series. Thus $M\mid L$. This proves both
output assertions.\hfill$\square$

\paragraph{Scope and dictionary.}
The input is P-F31, \texttt{work/positivity\_actual\_cubic\_pole.tex},
the theorem ``actual cubic pole'' and its denominator consequence;
P-F56 makes the two-adic valuation of a quasi-period exact. This is
an obstruction to an \emph{expanded} Hilbert-series representation.
The integer $M$ itself has only $n$ binary digits, so this proves no
exponential lower bound for arithmetic circuits, succinct rational
expressions, evaluation, or positivity. In geometric language these
are Hilbert characters of the graded invariant coordinate rings;
the grading/action dictionary is SGA 3, I, 4.7.3.1 and the projective
construction is EGA II, 2.4.1--2.4.7. No statement about the size of
a defining system or a Quot/Hilbert parameter space follows from
the length of this particular output format.

\subsection{NN-G003: arbitrarily large strict indices with no support holes}
\paragraph{Statement.}
For every $L\ge2$ define
\[
 f_L(n)=L\lfloor n/L\rfloor+1\quad(n\ge0).
\]
Its support is all of $\mathbb N$, while the positivity and strict
saturation indices in GCT VI, Definitions 1.2.2 and 1.2.4, are both
$L-2$. It has a positive witness description of size polynomial in
$\log L+\log(n+1)$, and
\[
 \sum_{n\ge0}f_L(n)t^n=
 \frac{1+(L-1)t^L}{(1-t)(1-t^L)}.
 \tag{G003.1}
\]

\paragraph{Proof.}
The constituent on $n\equiv r\pmod L$ is $X-r+1$.
For $f_L(n+p)$ the constituent on $n\equiv r$ is
$X+p+1-((r+p)\bmod L)$. Its least constant coefficient is
$p-L+2$. All coefficients are nonnegative exactly when $p\ge L-2$.
All these constituents have slope one; each is positive at every
integer $X\ge1$ exactly when its value at one is positive. The least
such value is $p-L+3$, and strict positivity again means
$p\ge L-2$. No constituent is identically zero. The count is always
at least one. It counts a distinguished object and the pairs
$(i,j)$ with $1\le i\le\lfloor n/L\rfloor$ and $1\le j\le L$.
Binary comparisons and division verify this set with polynomial
resources. Finally
$f_L(n)=1+L\sum_{k\ge1}\mathbf1_{n\ge kL}$ gives (G003.1).
\hfill$\square$

\paragraph{What this resolves.}
The inference ``exponential strict saturation index implies an
exponentially distant realized hole'' is false even for compactly
described positive counts. This is a logical countermodel, not a
new plethysm family. P-F57 supplies actual plethysm examples with
large strict indices; its off-residue evaluation must not silently
be replaced by an actual vanishing coefficient. The supplied
positive-start assertion is retained at statement level only.

\paragraph{Geometric translation.}
The series is the Hilbert series of the graded
$k[x,y]$-module
$k[x,y]\oplus k[x,y](-L)^{\oplus(L-1)}$, with
$\deg x=1$, $\deg y=L$. Thus the same failure of inference occurs
for effective coherent graded modules, equivalently equivariant
modules under SGA 3, I, 4.7.3. The associated sheaf construction is
EGA II, Proposition 2.5.2. This uses a module, not a normal domain
or an actual GCT multiplicity algebra; no stronger geometric
realization is asserted.

\subsection{NN-G004: cancellation and precision on the cubic pole formula}
\paragraph{Statement.}
Fix $0<c<1/4$. At the witness $x_n$ in P-F57, write the specific
pole expansion as $f_n(x_n)=\sum_\rho a_{\rho,n}$, with the dyadic
wave $W_n(x_n)$ among its grouped summands. If $f_n(x_n)>0$, its
relative summation condition number satisfies
\[
 \kappa_n:=\frac{\sum_\rho|a_{\rho,n}|}{f_n(x_n)}
 \ \ge\ 2^{(1/24-c/6)n^4-O(n^3\log n)}.
 \tag{G004.1}
\]
If $f_n(x_n)=0$ and the nonzero wave is retained, the relative
condition number is infinite. The lower bound is for this formula;
it gives a loss of $\Omega(n^4)$ bits of relative accuracy, not an
exponential number of required bits.

\paragraph{Proof.}
P-F57 gives
$2W_n(x_n)\ge2^{n^4/24-2n^3-1}$ and
$f_n(x_n)\le2^{(c/6)n^4+O(n^3\log n)}$.
The sum of the absolute values of the pole summands is at least
the absolute value of their dyadic subsum, hence at least
$|W_n(x_n)|$. Division gives (G004.1).
For a finite sum $F=\sum a_j\ne0$, relative perturbations
$|\Delta a_j|\le\eta|a_j|$ change it by at most
$\eta\sum|a_j|$; equality in the bound is attained by aligning
the perturbations with one complex direction. This proves the
stated relative condition number. A relative error of size
$2^{-b}$ in the terms can therefore cause relative error as large
as $2^{-b}\kappa_n$ in the sum. Taking logarithms proves the bit
statement.\hfill$\square$

\paragraph{Inputs, geometry and complexity.}
The inputs are ALPOGE\_GCT P-F56d and P-F57, with the exact wave
and dimension estimates in the supplied intake part 4 and the
paired-wave/strict-index proof files. A cyclotomic partial-fraction
expansion is a character computation in a localized representation
ring; the tensor and restriction operations themselves appear in
cote 2, scan 17, and symmetric-group operations in cote 133,
scan 28. A virtual-character cancellation is not a proof that
no different effective realization exists. The inequality gives
no Boolean or algebraic circuit lower bound and no obstruction
to exact integer arithmetic or symbolic cancellation. It makes
the ``near-supported zero'' precise only relative to a specified
expansion and norm.

\subsection{NN-G005: the universal support valuation and its exact linearization}
\paragraph{Statement.}
Let $L$ be a finite distributive lattice with least element $0$,
and let $P$ be its poset of join-irreducibles. The universal abelian
group for normalized inclusion--exclusion valuations is
\[
 V(L):=\mathbb Z[L]/\langle[0],\ [x]+[y]-[x\vee y]-[x\wedge y]\rangle
 \simeq\mathbb Z^P.                                      \tag{G005.1}
\]
Under the identification $L\simeq\mathcal I(P)$ with order ideals,
the universal valuation is $I\mapsto\sum_{p\in I}e_p$.
In contrast, the abelian group completion of the join monoid
$(L,\vee,0)$ is zero. For a Boolean lattice on $r$ atoms,
the kernel of the augmentation $V(L)\to\mathbb Z$ is $A_{r-1}$
with its usual integral coordinates; this is an extra valuation
construction, not the group completion of idempotent addition.

\paragraph{Proof.}
For completeness, every element of a finite lattice is the join
of the join-irreducibles below it, by induction on the interval
below it. In a distributive lattice a join-irreducible $p$ is
join-prime: if $p\le x\vee y$, then
$p=(p\wedge x)\vee(p\wedge y)$, so $p\le x$ or $p\le y$.
Consequently $x\mapsto\{p:p\le x\}$ identifies $L$ with the
order ideals in $P$; the inverse takes joins. Characteristic
functions of ideals satisfy inclusion--exclusion, giving a map
$V(L)\to\mathbb Z^P$. Define the inverse on a basis vector by
\[
 e_p\longmapsto[\mathord\downarrow p]
                    -[\mathord\downarrow p\setminus\{p\}].
\]
If $p$ is maximal in an ideal $I$, use inclusion--exclusion for
$I\setminus\{p\}$ and $\mathord\downarrow p$ to obtain
\[
 [I]=[I\setminus\{p\}]+[\mathord\downarrow p]
              -[\mathord\downarrow p\setminus\{p\}].
\]
Induction on $|I|$ proves the inverse assertion and the universal
property. A homomorphism from the join monoid to an abelian group
sends $x\vee x=x$ to $2a=a$, hence every $x$ to zero.
For an antichain $P$ with $r$ elements, augmentation is
$(n_1,\ldots,n_r)\mapsto\sum_i n_i$; its kernel is precisely the
integral root lattice with basis $e_i-e_r$, $i<r$.
\hfill$\square$

\paragraph{Scope and geometric translation.}
This corrects the meaning of ``support linearization'' in
\texttt{03\_split\_zero\_defect\_preprint.tex}, its support section,
and \texttt{04\_primitive\_support\_defects\_preprint.tex}, the
primitive-support construction. A finite set of locally closed
strata gives an inclusion--exclusion valuation to its constructible
functions; assigning an Euler characteristic, a compactly supported
cohomology class, or a rank to each stratum is a further choice of
the weights $e_p\mapsto a_p$. For finite disjoint rational points
this is literal $K_0$ of their vector spaces. The external/direct
sum and tensor operations are the framework of cote 2, scan 17;
SGA 5, III, 4.7 concerns the subsequent trace realization, not a
canonical trace on the bare lattice. This theorem creates neither
witness multiplicities nor an effective counting algorithm from
Boolean support.

\subsection{NN-G006: faithful Artin--Tate zeta encoding without RH}
\paragraph{Statement.}
Let
$\mathcal A=\mathbb Z[q^{\pm1},\varepsilon]/(\varepsilon^2-1)$ and
$\chi=\chi_{-4}$. The assignment
\[
 a=\sum_j(a_j+b_j\varepsilon)q^j\longmapsto
 Z_a^{(2)}(s)=\prod_j\zeta^{(2)}(s-j)^{a_j}
                         L(s-j,\chi)^{b_j}                 \tag{G006.1}
\]
is injective, even as an identity of meromorphic functions. Here
$\zeta^{(2)}=(1-2^{-s})\zeta$ and all products are finite. It
identifies multiplication of coefficients with the Witt tensor
product of local Euler factors. No RH is needed for injectivity.
For the full completed products the correct factors are
\[
 \Lambda_\zeta(s)=\pi^{-s/2}\Gamma(s/2)\zeta(s),\qquad
 \Lambda_\chi(s)=(4/\pi)^{(s+1)/2}\Gamma((s+1)/2)L(s,\chi).
 \tag{G006.2}
\]
The signed pole order of
$\prod_j\Lambda_\zeta(s-j)^{a_j}\Lambda_\chi(s-j)^{b_j}$ at an
integer $m$ is $a_m+a_{m-1}$. Consequently these signed orders
recover the finitely supported sequence $(a_j)$. The conventional
entire $\xi(s)=\tfrac12s(s-1)\Lambda_\zeta(s)$ cannot be used
in this pole-recovery formula.

\paragraph{Proof.}
All factors in (G006.1) have absolutely convergent Dirichlet
series in a common sufficiently far right half-plane, even with
integer negative exponents. If their product is one, its logarithm
there has zero coefficients by uniqueness of absolutely convergent
Dirichlet series. A quick proof of that uniqueness takes the least
index with nonzero coefficient, multiplies by its $s$th power,
and lets real $s\to+\infty$; absolute convergence at a fixed
abscissa dominates the tail. At the coefficient indexed by $p^n$
this gives
\[
 \sum_j\{a_j+\chi(p)^n b_j\}p^{nj}=0\quad(n\ge1).
\]
At $p=5$, the distinct nonzero numbers $5^j$ and the Vandermonde
matrix for their first finitely many powers give $a_j+b_j=0$
for every $j$. At $p=3$, use odd $n=2r+1$ and the distinct bases
$3^{2j}$ to obtain $a_j-b_j=0$. Thus $a_j=b_j=0$. The same proof
works formally from the Euler factors.
Locally the Witt product sends
$(1-\alpha T)^{-1}$ and $(1-\beta T)^{-1}$ to
$(1-\alpha\beta T)^{-1}$. Because $\chi(p)^2=1$ at odd primes,
this realizes exactly multiplication in $\mathcal A$.
The classical functional equations (DLMF 25.4.1--25.4.4 and
25.15.5, with the elementary Gauss sum for $\chi_{-4}$) give
(G006.2) simple poles for
$\Lambda_\zeta$ at $0,1$, and no others, whereas $\Lambda_\chi$
is entire and nonzero on the real axis. The latter nonvanishing
follows for positive real $s$ by pairing successive terms of the
alternating beta series, and for $s\le0$ by its functional equation;
the value at zero follows from that at one. Similarly the completed
zeta is finite and nonzero at real points other than $0,1$.
Only shifts $j=m,m-1$ therefore contribute to the signed pole order
at $m$. Starting beyond the support and proceeding downwards
recovers $a_j$. Multiplying by $s(s-1)/2$ removes the poles of
$\Lambda_\zeta$, proving the warning about $\xi$.
\hfill$\square$

\paragraph{Inputs and exact dictionary.}
This repairs and strengthens the injectivity/pole statements in
\texttt{OURS\_CONSULT\_a157422a\_COUNTING\_CONSTRUCTION.md},
Theorem 4(a),(d),(e). The zeta and primitive odd Dirichlet functional
equations have the exact classical locators DLMF 25.4.1--25.4.4
and 25.15.5 (Apostol, \emph{Introduction to Analytic Number
Theory}, equations (12)--(16), pp.\ 259--263); A is required only when
one additionally locates their nontrivial zeros on the lines
$\operatorname{Re}s=j+1/2$. The geometric objects are Tate twists
of the trivial and Gaussian quadratic compatible rank-one systems.
For each $\ell$ their local systems are lisse on
$\operatorname{Spec}\mathbb Z[1/2\ell]$; the factor at $\ell$
is read using a different auxiliary prime. Their representation
ring is $\mathcal A$; SGA 3, I, 4.7.2--4.7.3 describes the
comodule/graded representation formalism. The Witt operations
are explicitly present in cote 43, scan 37. The local factor at
$2$ is omitted for tensor multiplicativity; restoring it as an
ordinary scalar correction does not make an unramified sign
character exist at the ramified prime. Faithful encoding is not
efficient decoding and does not prove PH1 or a counting bound.

\subsection{NN-G007: all effective lifts, and the two different effective cones}
\paragraph{Statement.}
For a Laurent polynomial $c=\sum_j c_jq^j\in\mathbb Z[q^{\pm1}]$,
write $c=c_+-c_-$ coefficientwise with disjoint nonnegative parts.
Every effective class $\widehat c$ in $\mathcal A$ with
$\widehat c|_{\varepsilon=-1}=c$ is uniquely of the form
\[
 \widehat c=c_++h+\varepsilon(c_-+h),\quad
 h\in\mathbb N[q^{\pm1}],\qquad
 d:=\widehat c|_{\varepsilon=1}=|c|+2h.             \tag{G007.1}
\]
The Burnside linearization for $C_2$ is an integral ring
isomorphism
\[
 B(C_2)=\mathbb Z[x]/(x^2-2x)\ \simeq\
 R(C_2)=\mathbb Z[\varepsilon]/(\varepsilon^2-1),
 \qquad x\mapsto1+\varepsilon.                    \tag{G007.2}
\]
Its effective-set cone maps to $a+b\varepsilon$ with
$a\ge b\ge0$, strictly smaller than the effective-representation
cone $a,b\ge0$. The nonzero regular class $1+\varepsilon$ is not
an idempotent: its square is $2(1+\varepsilon)$. It maps to the
supported zero under the quotient of the \emph{effective semiring}
\[
 \mathbb N[\varepsilon]/(\varepsilon^2-1)\longrightarrow G(\mathbb Z),
 \quad a+b\varepsilon\longmapsto
 \begin{cases}\tau,&a=b=0,\\a-b\text{ supported},&a+b>0.\end{cases}
\]
This is not a quotient homomorphism from the whole representation ring.

\paragraph{Proof.}
For each coefficient solve $a-b=c_j$ with $a,b\ge0$. If $c_j\ge0$
the solution is $(c_j+h_j,h_j)$, and if $c_j<0$ it is
$(h_j,-c_j+h_j)$. This proves (G007.1), including uniqueness.
A finite $C_2$-set is a disjoint union of fixed points and free
two-point orbits. The latter square is two free orbits, proving
the Burnside presentation. Substitution $\varepsilon=x-1$
is inverse to (G007.2). A set with $r$ fixed points and $s$ free
orbits has permutation representation $(r+s)+s\varepsilon$,
which gives exactly the asserted cone. The sign representation
$\varepsilon$ is effective but lies outside that cone. The
algebraic calculation of $(1+\varepsilon)^2$ proves nonidempotence.
The displayed map preserves the operations because evaluation at
$\varepsilon=-1$ does so, while a sum or product of two effective
classes is zero only in the usual empty cases. It is surjective:
use $a=n,b=0$ for positive $n$, $a=0,b=-n$ for negative $n$,
and $a=b=1$ for the supported zero.\hfill$\square$

\paragraph{Geometric consequences.}
At an inert odd prime the free orbit is
$\operatorname{Spec}\mathbb F_{p^2}$: it has no $\mathbb F_p$
points and two $\mathbb F_{p^2}$ points. Its geometric support,
its dimension, and its rational-point count are therefore distinct
invariants. Globally its Artin representation comes from
$\operatorname{Spec}\mathbb Z[i,1/2]$ and has zeta
$\zeta^{(2)}(s)L(s,\chi)$.
Changing $h$ in (G007.1) changes the Euler product by
\[
 \prod_j\zeta_{\mathbb Q(i)}^{(2)}(s-j)^{h_j}.
\]
A puts the nontrivial zeros of all these choices on their
appropriate lines, but selects none of the choices. Negative
Tate powers are legitimate twists of an $\ell$-adic object;
$q^{-1}$ itself is not the point count of a scheme over
$\mathbb F_q$. Thus an effective Artin--Tate representation is
not automatically an effective scheme count.
The exact archive comparison is cote 133, scan 19, where passing
from equivariant sets to representations must respect this
distinction; its subset/exterior-power formula already needs the
correction recorded in ALPOGE\_GCT PL-01. The group-scheme
representation formalism is SGA 3, I, 4.7.2. This gives no way to
recover satisfiability from a virtual cancellation or an
arbitrarily chosen dimension lift.

\subsection{NN-G008: rational descent does not imply a finite Tate count}
\paragraph{Statement and proof.}
The all-rank $\mathbb Q(D)$ descent in the counting consult does
not imply Laurent-polynomial structure constants in every
$\mathbb Q(D)$-rational basis. Already take the two-dimensional
algebra with basis $1,P$ and relation $P^2=DP$, and replace
$P$ by $b=P/(D^2+1)$. Then
\[
 b^2=\frac{D}{D^2+1}\,b.                            \tag{G008.1}
\]
This rational function is not a Laurent polynomial in $q$ for
$D=q+2+q^{-1}$. Indeed a Laurent polynomial invariant under
$q\leftrightarrow q^{-1}$ belongs to $\mathbb Q[q+q^{-1}]
=\mathbb Q[D]$: induction removes the highest pair
$q^m+q^{-m}$, itself a polynomial in $q+q^{-1}$. But
$D/(D^2+1)$ is not a polynomial in $D$, by division in
$\mathbb Q[D]$.\hfill$\square$

\paragraph{Correction and dictionary.}
Theorem 1(c) of the consult correctly distinguishes rational
descent from integrality; its Theorem 4(c) and final rank-four
summary do not preserve that distinction. A rational function
can also have an even-order pole at $v=1$, so ``even vanishing
order'' requires regularity there. For example rescaling by
$D-4$ gives poles. In the geometric dictionary a generic fibre
and a rational change of basis do not choose an integral lattice
or an effective class in $K_0$; SGA 3, I, 4.7.3 is a grading
equivalence, not an integrality theorem. Exterior powers and
effective direct summands, discussed in cote 133, scan 28 and
SGA 6, VI, Theorem 3.2, require actual vector-bundle classes.
The remaining rank-four gate includes a defined canonical basis,
integral Laurent constants, and an effective compatible lift.
The formal substitution $v=\sqrt q$ constructs neither a CM
Hecke character nor a compatible arithmetic family.

\subsection{NN-G009: a quantitative size obstruction for all-word Weil realizations}
\paragraph{Statement.}
Fix a prime power $q>1$ and $k\ge1$. If a finite-dimensional
Weil object $V_k$ has traces
\[
 \operatorname{Tr}(F^n\mid V_k)=
 40(q^{n/2}+q^{-n/2})^{2k}\quad(n\ge1),
 \tag{G009.1}
\]
then its semisimplification has eigenvalue $q^a$ with multiplicity
$40\binom{2k}{k+a}$ for every $-k\le a\le k$, and
$\dim V_k=40\cdot4^k$. If instead (G009.1) is an inserted trace
$\operatorname{Tr}(A_kF^n\mid V)$ on a fixed $V$, at least
$2k+1$ distinct nonzero eigenvalues of $F$ are required. In
particular no fixed finite $V$ carries the traces for all $k$.

\paragraph{Proof.}
The binomial expansion gives
\[
 40(q^{n/2}+q^{-n/2})^{2k}
 =\sum_{a=-k}^k40\binom{2k}{k+a}(q^a)^n.
\]
Distinct nonzero exponential sequences are linearly independent:
their values at the first $r$ positive integers form an invertible
Vandermonde matrix after scaling columns. Ordinary traces are
power sums of eigenvalues with their algebraic multiplicities,
even with Jordan blocks, so uniqueness proves the first assertion.
Their sum is $40\cdot2^{2k}$. For an inserted trace, a Jordan
decomposition expresses it as a sum of polynomials in $n$ times
$\lambda^n$ over eigenvalues $\lambda$ of $F$. These
exponential-polynomial sequences are independent: applying
successive difference operators $E-\lambda$ isolates each
base, or equivalently their generating functions have poles at
distinct reciprocals. Hence every base $q^a$ must occur among
the eigenvalues of $F$, proving the dimension lower bound.
\hfill$\square$

\paragraph{Inputs, geometry and complexity.}
The word-trace identity is GCT\_NSRH F42 and counting consult
Theorem 4(c). The new assertion quantifies the necessary size;
it does not rebrand F42 as new. A realization is $40$ copies of
$H^*((\mathbb P^1)^{2k})(k)$, with all cohomological degrees even.
It has a short tensor description even though its total rank is
exponential. Tensor products and trace are the operations of
cote 2, scan 17 and SGA 5, III, 4.7. This explains why a finite
list of multiplication objects can coexist with the absence of
a single bounded-rank Weil object. Neither a short tensor
description nor A gives a polynomial-size positive polytope
formula or rank-four canonical basis.

\subsection{NN-G010: the precise category supplied by a dimension ring}
\paragraph{Statement.}
Let $B$ be a finite based unital algebra over
$\mathbb Z[q^{\pm1}]$, with its unit among the basis elements.
Suppose its constants $c_{ab}^z$ and an associative unital
nonnegative table $d_{ab}^z$ satisfy
$d\ge|c|$ coefficientwise and $d\equiv c\pmod2$.
Then the effective table
\[
 \widehat c_{ab}^z=(d_{ab}^z+c_{ab}^z)/2
                +\varepsilon(d_{ab}^z-c_{ab}^z)/2
 \tag{G010.1}
\]
is realized by a monoidal additive category whose \emph{split}
Grothendieck ring is free over $\mathcal A$ on this basis. The
construction need not be idempotent complete, abelian, rigid,
semisimple as a category, or have a simple tensor unit. It
supplies no canonical basis or uniform complexity bound.

\paragraph{Proof.}
The map $\mathcal A\hookrightarrow
\mathbb Z[q^{\pm1}]\times\mathbb Z[q^{\pm1}]$ given by
evaluation at $\varepsilon=1,-1$ is injective, with image the
pairs congruent modulo two. Thus (G010.1) is associative and
unital. Let $\mathcal T$ be finite-dimensional
$\mathbb Z\times C_2$-graded vector spaces over a characteristic
zero field, with simple objects $q^j\varepsilon^a$ and their
ordinary tensor product. Its isomorphism classes are exactly
the effective elements of $\mathcal A$.
In $\operatorname{Mat}_{|B|}(\mathcal T)$ use matrix convolution
as tensor product. For each $a\in B$ choose the matrix object
$L_a$ with $(z,b)$ entry of class $\widehat c_{ab}^z$.
Associativity of the table gives an isomorphism
\[
 L_a\otimes L_b\simeq\bigoplus_z
        \widehat C_{ab}^z\otimes L_z,
\]
because objects of $\mathcal T$ are determined by their classes.
Take the replete full additive subcategory consisting of finite
sums of shifts $q^j\varepsilon^aL_b$. It is tensor closed and
inherits the coherent associativity constraints of matrices.
The unit matrix is $L_1$.
To compute its split $K_0$, inspect the column indexed by the
unit basis element: the unit law makes that column of $L_b$
the basis vector indexed by $b$. Therefore two sums of shifted
$L_b$ are isomorphic exactly when their effective coefficients
agree. Their additive monoid is free over the effective
coefficient monoid, and group completion gives the stated
free split $K_0$.
In the ambient full matrix category the endomorphisms of the
unit are $|B|$ independent scalar diagonal entries. Their
idempotents need not split within the chosen subcategory.
Thus neither idempotent completeness nor a simple unit follows;
taking a Karoubi or abelian completion changes the question.
\hfill$\square$

\paragraph{Geometric translation.}
For $X$ a disjoint union of $|B|$ rational points, matrices of
Tate--Artin objects are kernels on $X\times X$. Convolution is
$p_{13!}(p_{12}^*(-)\otimes p_{23}^*(-))$. This is the finite
zero-dimensional instance of the external product and projection
formula on cote 2, scan 17; Artin gradings use SGA 3, I, 4.7.2.
The construction refines counting consult Theorem 5 and
Corollary 8 by stating precisely which Grothendieck group is
proved. It uses the already known structure constants to build
the objects. It cannot establish their positivity without those
constants, nor establish PH0 or PH1 for a canonical geometric
construction.

\subsection{NN-G011: a smaller composing rank-three dimension ring}
\paragraph{Statement.}
The supplied rank-three algebra admits a dimension ring with
the same 132 nonzero positions and a total of $7,276$
coefficientwise cancelling pairs, of which $6,612$ occur in
the bottom-cell coefficients. This improves the supplied
search-window value $7,952$. No claim of global minimality is
made. The parameters are
\[
 \begin{gathered}
 D=q+2+q^{-1},\quad\delta=(D-1)(D-2),\quad W=7D-8,\\
 \Lambda_1=D,\quad\Lambda_2=3D^2-10D+4,\quad
 G_1=2W\delta,\quad G_2=W\delta,\quad S=2W\delta^2,\\
 U_1=13D-16,\quad U_2=3D-4,\quad
 V_1=14D^2-17D,\quad V_2=2D^2+D-2.
 \end{gathered}                                             \tag{G011.1}
\]
For $x,y\in\{1,12,2,21\}$ write the primitive word reduction
$P_xP_y=D^\alpha\Lambda_i^\beta P_z$, where $\beta=0,1$.
Keep $\mu$ as unit and define
\[
 \begin{aligned}
 \Sigma^2&=S\Sigma,&
 \Sigma\gamma_x^i=\gamma_x^i\Sigma&=D^{|x|-1}G_i\Sigma,\\
 \gamma_x^i\gamma_y^j&=D^{|x|+|y|-2}W\Sigma&&(i\ne j),\\
 \gamma_x^i\gamma_y^i&=\delta D^\alpha\Lambda_i^\beta\gamma_z^i
       +\begin{cases}
          D^{|z|+\alpha-2}U_i\Sigma,&\beta=0,\\
          D^{|z|+\alpha-1}V_i\Sigma,&\beta=1.
        \end{cases}
 \end{aligned}                                             \tag{G011.2}
\]
The exponent in the $\beta=0$ case is nonnegative for the
eight word pairs where it occurs.

\paragraph{Proof of associativity.}
Over $\mathbb Q(D)$ take two $2\times2$ matrix blocks with
\[
 P_1=\begin{pmatrix}D&\Lambda_i\\0&0\end{pmatrix},\qquad
 P_2=\begin{pmatrix}0&0\\1&D\end{pmatrix}.
\]
They obey $P_j^2=DP_j$ and $P_1P_2P_1=\Lambda_iP_1$,
$P_2P_1P_2=\Lambda_iP_2$. The four $P_x$ are linearly
independent generically because
$\Lambda_i(D^2-\Lambda_i)\ne0$. In the direct product of
these blocks and two scalar blocks send $\gamma_x^i$ to
$\delta P_x$ in its own matrix block, to zero in the other,
to $D^{|x|-1}G_i$ in the first scalar block and to zero in
the last. Send $\Sigma$ to $S$ in the first scalar block
and zero elsewhere, and $\mu$ to the identity in every
block. These ten images are independent, successively by
the last scalar, the two matrix blocks, and the first scalar.
They are closed under products because
\[
 \frac{G_1G_2}{S}=W,\quad
 \frac{G_i(G_i-\delta D)}S=U_i,\quad
 \frac{G_i(DG_i-\delta\Lambda_i)}S=V_i.
\]
The first ratio follows immediately from (G011.1); the next
two equal $2W-D,(W-D)/2$ and
$2DW-\Lambda_1,(DW-\Lambda_2)/2$, respectively.
They are exactly the polynomials displayed there. Evaluating
the scalar block gives (G011.2). The embedded algebra is
associative over $\mathbb Q(D)$, hence all its polynomial
structure identities hold over $\mathbb Z[D]$.

\paragraph{Proof of effectivity, domination and parity.}
Write $u=q+q^{-1}=D-2$. The true cell parameters are
$\lambda_1=D$, $\lambda_2=u^2$ and
$\delta_2=-\delta_1=(D-1)(D-4)$; the latter has Laurent
coefficients $q^2-q-q^{-1}+q^{-2}$.
Our $\delta$ is its coefficientwise absolute value plus $2$.
Moreover
\[
 \Lambda_2-\lambda_2=2(q^2+q^{-2})+2(q+q^{-1})\ge0.
\]
The true bottom parameters are
$W_0=U_{1,0}=U_{2,0}=D$,
$V_{1,0}=5D-4$, $V_{2,0}=2D^2-5D+4$,
$G_{1,0}=4D(D-1)$, $G_{2,0}=D^2(D-1)$, and
$S_0=4D^2(D-1)^2$. Their differences from (G011.1) are
\[
 \begin{array}{c|l}
 W-W_0&6u+4\\
 U_1-U_{1,0}&12u+8\\
 U_2-U_{2,0}&2u\\
 V_1-V_{1,0}&14u^2+34u+16\\
 V_2-V_{2,0}&6u+6\\
 G_1-G_{1,0}&14(q^3+q^{-3})+22(q^2+q^{-2})+42u+36\\
 G_2-G_{2,0}&6(q^3+q^{-3})+8(q^2+q^{-2})+16u+12\\
 S-S_0&(u+1)^2\{14(q^3+q^{-3})+8(q^2+q^{-2})+26u\}.
 \end{array}                                               \tag{G011.3}
\]
Every difference is coefficientwise nonnegative and even.
All true bottom parameters are nonnegative Laurent
polynomials, so the new ones are also. Multiplication by
nonnegative powers of $D$ preserves these properties.
For within-cell coefficients use
$|fg|\le|f||g|$ coefficientwise, together with the inequalities
for $\delta$ and $\Lambda_i$. Congruence modulo two is
preserved by sums and products and is unaffected by the
opposite signs of the two true $\delta_i$. This proves every
dimension-ring condition for all 132 positions.

\paragraph{Exact score.}
The nineteen unit positions contribute no excess. In the
true table the sum of absolute Laurent coefficients of the
bottom constants is $4096$. This is obtained at $D=4$ from
\[
 S_0+20(G_{1,0}+G_{2,0})+200W_0
           +20(U_{1,0}+U_{2,0}+V_{1,0}+V_{2,0})=4096.
\]
For within-cell constants the types
$(\alpha,\beta)=(0,0),(1,0),(0,1),(1,1)$ have multiplicities
$2,6,6,2$. The absolute coefficient norms in cell one are
$4,8,8,24$, and in cell two $4,8,16,12$, by direct expansion
of $\delta_iD^\alpha\lambda_i^\beta$. Their weighted sum is
$328$. Thus the total true absolute norm is $4424$.
For the new table, at $D=4$ the values are
\[
 (\delta,W,G_1,G_2,S,U_1,U_2,V_1,V_2,\Lambda_1,\Lambda_2)
 =(6,20,240,120,1440,36,8,156,34,4,12).
\]
The same bottom formula gives $17320$, and the within-cell
sum is $6(26+14\cdot4)+6(26+14\cdot12)=1656$.
Hence the excess divided by two is
$(18976-4424)/2=7276$, with bottom part
$(17320-4096)/2=6612$.\hfill$\square$

\paragraph{Inputs, certificates and meaning.}
The exact input is counting consult Theorem 2 and
\texttt{OURS\_E4\_scripts/e3\_constants.json}.
\texttt{verify\_dimension\_7276.py} independently builds the
table with integer Laurent dictionaries, checks all 1000
associativity triples, and verifies all coefficient inequalities
against the supplied 132 constants. It does not import the
supplied dimension-ring builder. The full certificate is
\texttt{dimension\_certificate\_7276.json}.
NN-G010 converts the result into a finite-kernel monoidal
category, with the same SGA 3/cote 2 dictionary. This lowers
an explicit upper bound for the noncanonical cancellation
cost. It neither selects a canonical lift, proves global
optimality, nor decides PH0/PH1 at rank four. No unimodality
claim is made for this particular smaller lift.

\subsection{NN-G012: why the sign-point trace symmetry reappears}
\paragraph{Statement.}
In the specialization $q=-1$, $D=0$ of the supplied based
rank-three algebra, let $\gamma_x=\gamma_x^2$ and set
\[
 \beta_1=\gamma_1+\Sigma/4,\quad
 \beta_2=\gamma_2+\Sigma/4,\quad
 \beta_{12}=\gamma_{12},\quad\beta_{21}=\gamma_{21}.
\]
Their span is a matrix subalgebra over $\mathbb Q$, identified
with $M_2$ by
\[
 \beta_1\mapsto8E_{12},\quad\beta_2\mapsto8E_{21},\quad
 \beta_{12}\mapsto16E_{11},\quad\beta_{21}\mapsto16E_{22}.
 \tag{G012.1}
\]
Here $\Sigma^2=0$ and $\Sigma$ annihilates this subalgebra.
The matrix-trace Gram matrix, in the order $1,12,21,2$, is
\[
 \begin{pmatrix}0&0&0&64\\0&256&0&0\\0&0&256&0\\64&0&0&0\end{pmatrix}.
 \tag{G012.2}
\]
Thus the cyclic slot relation is exactly the ordinary cyclic
trace identity in a rescaled matrix-unit basis. This explains
the monomial pairing at this specialization without asserting
a symmetry of the generic table or of its six unimodality
failures.

\paragraph{Proof.}
At $D=0$ the true parameters are $\delta_2=4$,
$\lambda_2=4$; every $\Sigma\gamma_x$ and $\Sigma^2$ is zero.
Use $P_1=2E_{12}$, $P_2=2E_{21}$. The within-cell
coefficients agree with the products of $4P_x$, and the
only nonzero bottom coefficients in its sixteen products
are $4\Sigma$ in the four products whose cell output is
$16\gamma_1$ or $16\gamma_2$. Replacing those outputs by
$16\beta_1$ or $16\beta_2$ absorbs exactly $4\Sigma$;
all other products require no correction. This proves closure
and (G012.1). The four images are independent, so the map is
an isomorphism. Direct matrix multiplication gives (G012.2).
For the supplied tensor-Hecke trace, direct evaluation
with its exact arithmetic (\texttt{check\_signpoint\_Hecke\_trace.py},
reusing the supplied implementation) shows that it vanishes on $\Sigma,\gamma_1,
\gamma_2$ and takes value $16$ on $\gamma_{12},\gamma_{21}$;
it therefore restricts to this matrix trace. Finally
$\operatorname{tr}(ABC)=\operatorname{tr}(BCA)$ is a direct
index calculation.\hfill$\square$

\paragraph{Geometry and scope.}
The full sixteen-product certificate is
\texttt{verify\_signpoint\_matrix\_block.py} and
\texttt{signpoint\_matrix\_certificate.json}. The four-dimensional
kernel algebra on a two-point scheme is again the convolution
algebra of cote 2, scan 17; (G012.1) is its explicit basis change.
The trace cyclicity is the finite-dimensional case of SGA 6,
I, section 8 and SGA 5, III. At the sign point, the cyclic
permutation mixes four failed generic coefficients with two
shortest coefficients, as E4 section 6b records. Accordingly
this explains the restored trace symmetry, not generic
unimodality or the existence of a rank-four canonical basis.

\subsection{NN-G013: the missing operation in positive-formula isolation}
\paragraph{Statement.}
Suppose a nonnegative multiplicity $g(x)$ has a uniform witness
formula $g(x)=|W_x|$ with $W_x\subseteq\{0,1\}^{m(x)}$,
$m$ polynomially bounded, and polynomial-time verification.
If the parity of $|W_x\cap h^{-1}(0)|$ can be decided in
deterministic polynomial time for every affine Boolean hash
$h$, then positivity of $g$ is in RP. If this applies to one
of the unary NP-hard plethysm or Kronecker positivity families,
then $\mathrm{NP}=\mathrm{RP}$. The positive witness formula
alone only places positivity in NP.

\paragraph{Proof.}
Choose $k$ uniformly from $1,\ldots,m+2$ and an affine map
$h:\mathbb F_2^m\to\mathbb F_2^k$ uniformly. For distinct
inputs its outputs are independent and uniform. If
$N=|W_x|>0$, some such $k$ makes
$\mu=N2^{-k}\in[1/4,1/2]$. Put
$X=|W_x\cap h^{-1}(0)|$. Then
$\mathbb E X=\mu$ and
$\mathbb E X(X-1)=N(N-1)2^{-2k}\le\mu^2$.
Since $\mathbf1_{X=1}\ge X-X(X-1)$ for every nonnegative
integer $X$, the probability of a singleton is at least
$\mu-\mu^2\ge3/16$. A random $k$ therefore succeeds with
probability at least $3/(16(m+2))$. On that event parity is
one. On an empty witness set parity is always zero, so
polynomial repetition gives an RP algorithm. NP-hardness
gives $\mathrm{NP}\subseteq\mathrm{RP}$, and the usual
polynomial random string is an NP certificate for an RP
acceptance, proving equality.\hfill$\square$

\paragraph{Exact inputs, signed failure and complexity scope.}
The unary hardness locators are Fischer--Ikenmeyer,
Theorem 3.1, label \texttt{thm:plethysm-positivity-np-hard},
and Ikenmeyer--Mulmuley--Walter, \texttt{tkronnp} in the
supplied main source. The latter's positive formula for
an NP-hard subclass is \texttt{tkronpos}, so positivity of
a formula is not itself a tractability theorem.
For a signed sum isolation tests the presence of summands:
on $S_2$, $g((1,1),(1,1),(1,1))=(1-1)/2=0$, although
both character summands exist. A and M do not supply the
hashed parity oracle. E can evaluate an explicitly encoded
witness count after a parsimonious bounded-width CNF
encoding, but its exponent is measured in that encoding's
variables, including verifier and hash auxiliary variables;
the word ``polynomial'' does not preserve a linear exponent.
There is no derivation of NP=RP from the stated inputs.

\paragraph{Geometric translation.}
The witness relation and affine hashes are equations over
$\mathbb F_2$ on Boolean coordinates; verification constraints
give a finite affine scheme whose rational points are the
witnesses after parsimonious encoding. Its point-count
parity is additional information beyond support of its
terms. SGA 5, III, 4.7 interprets counts as alternating
Frobenius traces but does not evaluate them in polynomial
time. The finite-set/representation distinction is
NN-G007 and cote 133, scan 19. An arbitrary split-zero
coefficient formula is not such a witness scheme until
this identification is constructed.

\subsection{NN-G014: what a unique extremal localization actually proves}
\paragraph{Statement and proof.}
Let $K$ be a field with a valuation $\nu$ and let
$F=\sum_{i=1}^rF_i$. If exactly one summand $F_j$ has
least valuation, then
$\nu(F)=\nu(F_j)$ and its leading residue equals that of
$F_j$ in the associated graded ring. To see this divide
by $F_j$: every other ratio has positive valuation, so
$F/F_j$ has residue one. This proves $F\ne0$ as a
function, not $F(1)\ne0$ after a specialization: $1-t$
is the elementary counterexample.
For $K=\mathbb Q((t))$ and an actual coefficient formula
$g=\operatorname{CT}_tF$,
the stronger hypotheses $\nu(F_i)\ge0$ for all $i$, a unique
term of valuation zero, and positive residue for that term
prove $g>0$. More generally a unique exposed weight in a
multigraded expansion isolates its coefficient, provided
the target coefficient is exactly that exposed weight.
\hfill$\square$

\paragraph{Reduced gate and dictionary.}
For a GCT positivity algorithm this leaves three precise
tasks: a coefficient-preserving localization, polynomially
bounded certificates for the target exposure and residue,
and a procedure finding such a certificate for every yes
instance. A deterministic polynomial-time complete procedure
would yield P=NP on the known hard families; randomized
one-sided completeness yields NP=RP. A bare unique smallest
signed term does neither. The source is the question left
by the GCT support consult, section on P-F30/P-F23 and
reverse isolation. Geometrically a one-parameter subgroup
of a torus filters an equivariant module by weights; the
associated graded residue retains the exposed piece. The
exact action/graded-module dictionary is SGA 3, I, 4.7.3;
the projective/graded localization construction is EGA II,
2.2--2.5. No degeneration has been asserted to preserve
the desired multiplicity without the coefficient identity.

\subsection{NN-G015: support inclusion cannot decide multiplicity dominance}
\paragraph{Statement.}
Comparing two strictly positive succinct witness counts
is already PP-hard, even when their Boolean supports are
identically nonempty. Consequently the support inclusion
supplied by BIP cannot, solely as support information,
decide the multiplicity inequalities that remain possible
obstructions. This is a limitation of the information,
not a PP-hardness theorem for the actual GCT multiplicities.

\paragraph{Proof.}
Given a Boolean circuit on $n\ge1$ input bits take
$a(x)=\#\operatorname{CircuitSAT}(x)+1$ and
$b(x)=2^{n-1}+1$. Both have positive polynomially
verifiable witness counts: adjoin one distinguished
witness to the satisfying assignments, respectively to
all strings of length $n-1$. Then
$a(x)>b(x)$ exactly when more than half the assignments
satisfy the circuit. Majority circuit acceptance is the
defining PP computation with its polynomial-time verifier
unrolled into a polynomial-size Boolean circuit. The support bits of both counts are
always one and contain none of that comparison data.
For orbit closures $Z\subseteq\Omega$, restriction gives
the graded equivariant surjection
$\mathbb C[\Omega]\twoheadrightarrow\mathbb C[Z]$.
Complete reducibility degree by degree implies
$m_Z(\lambda)\le m_\Omega(\lambda)$. Hence a strict
reverse inequality obstructs containment, whereas
support inclusion alone asserts only that a positive
left multiplicity has a positive right multiplicity.
\hfill$\square$

\paragraph{Inputs, border scope and geometry.}
BIP, \emph{No occurrence obstructions in geometric
complexity theory}, Theorem 1.1, gives its stated
$n\ge m^{25}$ support inclusion for padded permanent
and determinant orbit closure. It does not prove
containment and does not eliminate multiplicity
obstructions or other equations. Orbit closure is the
geometric object for border complexity; confusing it
with the orbit or the support set would discard the
defining ideal. In Grothendieck language the needed
object is the equivariant closed subscheme, its
quotient coordinate algebra, and its graded locally
free representation pieces. Cote 133, scan 28 supplies
the Schur/tensor operations on their $K_0$ classes;
cote 2, scan 17 supplies restriction and tensor products.
No FGA representability theorem by itself computes the
ideal or the multiplicity comparison. The example in
this proof can be realized by finite Boolean witness
schemes with trivial group action, using succinct
equations; its explicit vector space may have
exponentially many basis vectors. Thus it makes no
unproved reduction into the BIP orbit-closure family.

\subsection{NN-G016: a cutoff projection repairs the finite-prime boundary calculation}
\paragraph{Statement and exact model.}
Fix $\ell>0$. Give $\mathbb Z\cup\{\infty\}$ the topology
in which the integers are isolated and $m\to\infty$ means
$m\to+\infty$. Put
\[
Y=(\mathbb Z\cup\{\infty\})\times\mathbb R,
\qquad T(m,y)=(m+1,y+\ell),\quad
T(\infty,y)=(\infty,y+\ell).
\]
For the crossed-product extension
\[
0\longrightarrow I=C_0(\mathbb Z\times\mathbb R)\rtimes_T\mathbb Z
\longrightarrow E=C_0(Y)\rtimes_T\mathbb Z
\longrightarrow Q=C_0(\mathbb R)\rtimes_{+\ell}\mathbb Z
\longrightarrow0,
\]
the connecting map $K_0(Q)\simeq\mathbb Z\to K_1(I)\simeq\mathbb Z$
is multiplication by $-1$ when the boundary rank-one class
is positive and increasing argument is the positive winding convention.
In particular it is an isomorphism. The proof supplies the
missing rank-one projection and a self-adjoint lift in $E$.

\paragraph{Proof: an explicit rank-one corner.}
Write $UfU^{-1}=\alpha_1(f)$, where
$\alpha_n(f)(m,y)=f(m-n,y-n\ell)$. Define the real cutoff
\[
c(y)=\begin{cases}\cos(\pi y/(2\ell)),& |y|\le\ell,\\0,&|y|>\ell.\end{cases}
\]
On every interval of length $\ell$ the two possible nonzero
terms are a sine and a cosine, so
$\sum_{n\in\mathbb Z}c(y-n\ell)^2=1$. Regarding $c$ as a
bounded continuous function on $Y$, independent of $m$, put
\[
P=\sum_{n\in\mathbb Z}c(y)c(y-n\ell)U^n\ \in M(E).
\]
This sum has only the three terms $n=-1,0,1$. It is
self-adjoint by the crossed-product involution. The coefficient
of $U^k$ in $P^2$ is
\[
\sum_n c(y)c(y-n\ell)^2c(y-k\ell)
=c(y)c(y-k\ell),
\]
so $P$ is a projection. Its boundary restriction $p$ belongs
to $Q$, because each coefficient is compactly supported in $y$.
In the boundary orbit representation indexed by
$\theta\in\mathbb R/\ell\mathbb Z$, its matrix is
\[
p_\theta(i,j)=c(\theta+i\ell)c(\theta+j\ell).
\]
It is the rank-one projection onto the unit vector with
entries $c(\theta+i\ell)$. The action on $\mathbb R$ is free
and proper. Its crossed product is locally the compact operators
over the circle; a rank-one projection nonzero on every fiber is
full. Thus $pQp\simeq C(\mathbb R/\ell\mathbb Z)$, and $[p]$
is the positive generator of $K_0(Q)$. More directly, compression
is scalar on every rank-one fiber, and multiplying $p$ by an
arbitrary continuous $\ell$-periodic function realizes every
scalar continuous section. This also exhibits the unit of the corner.

\paragraph{Proof: the lift and its winding.}
Set
\[
\varphi(t)=\tfrac12-\pi^{-1}\arctan t,
\qquad h(m,y)=\varphi(y-m\ell),\qquad h(\infty,y)=1.
\]
This is bounded, continuous and $T$-invariant. Hence $h$ is a
central multiplier of $E$. It is not asserted to belong to $E$.
Nevertheless $H=hP$ does belong to $E$: each of its three
coefficients has bounded $y$ support; as $m\to-\infty$ it tends
uniformly to zero there; as $m\to+\infty$ it tends uniformly to
its specified boundary coefficient. These are precisely the
vanishing-at-infinity and continuity conditions on this support.
Thus $H=H^*$ is a self-adjoint lift of $p$.

The exponential formula for the connecting map now gives
\[
\partial[p]=[u],\qquad
u=\exp(2\pi iH)=1+(\exp(2\pi ih)-1)P.
\]
On the open part use $t=y-m\ell$. This identifies its transformation
group with the product of $\mathbb R_t$ and the translation action
of $\mathbb Z$ on itself. Accordingly
$I\simeq C_0(\mathbb R_t,\mathcal K(\ell^2\mathbb Z))$.
In this description
\[
P_t=|v_t\rangle\langle v_t|,\quad
v_t(m)=c(t+m\ell),\quad \|v_t\|=1,
\qquad
u(t)=1+(w(t)-1)P_t,
\]
where
\[
w(t)=\exp(2\pi i\varphi(t))=\frac{t+i}{t-i}.
\]
The vectors $v_t$ are norm-continuous: locally only finitely many
coordinates can be nonzero. They trivialize a full rank-one
corner of $I$. The Morita map therefore sends $[u]$ to the scalar
class $[w]\in K_1(C_0(\mathbb R))$. In particular
$u-1\in I$, also seen directly from $|w(t)-1|\to0$ at both ends.
The continuous argument $2\pi\varphi(t)$ decreases from $2\pi$
to $0$ as $t$ increases from $-\infty$ to $+\infty$.
The resulting loop has winding number $-1$, which proves the claim.
\hfill$\square$

\paragraph{Source correction, scope and geometry.}
The explicit model is the one in corpus paper 03, lines 850--998,
and paper 04, lines 358--504, with $\ell=\log p$. Their scalar
$h$ cannot by itself lift a rank-one projection: its boundary
value is a multiplier unit in a nonunital algebra. The replacement
$hP$ removes precisely that gap. The classical analytic inputs
are Blackadar, \emph{K-Theory for Operator Algebras}, second edition,
\S9.3.2 (printed page 68), for the exponential connecting map,
and Williams, \emph{Crossed Products of C*-Algebras}, author draft
version 3.1 (6 September 2006), Corollary 4.11 and Remark 4.12
(printed page 126), for the free proper action Morita equivalence.
Both locators were checked in the primary author PDFs.

The group action on all of $Y$ is also free and proper: compact
sets have bounded $y$ coordinates, so only finitely many translates
can meet them. Thus the same free proper action theorem applies
to $E$, compatibly with its open and closed pieces. Transferring
the calculation to a larger arithmetic quotient still requires the
claimed identification of its edge extension with this exact model.
No verification of that global identification is smuggled into
the present model theorem, and no purity or complexity bound follows.

An honest algebraic-geometric comparison is available for the
scalar loop: $(t+i)/(t-i)$ is a rational function on
$\mathbb P^1_{\mathbb C}$ with divisor $[-i]-[i]$ and unit norm
on the real locus. Equivalently its real coordinates are
$((t^2-1)/(t^2+1),2t/(t^2+1))$ on the norm-one torus
$x^2+y^2=1$; the removed degree-two divisor $t^2+1=0$ has no
real points. The zero principal divisor class in the Picard
group and the nonzero winding on the real locus are different
invariants. In Grothendieck language a rank-one corner is an
actual rank-one object before taking $K_0$, not the literal unit
of a multiplier algebra. The supplied archive cote 2, scan 17
provides the analogous pullback/tensor formalism for geometric
objects; it does not identify this analytic $K_1$ class with an
\'{e}tale trace or Weil weight. There is no claimed EGA/SGA theorem
turning the crossed-product model into an algebraic scheme.

\subsection{NN-G017: fixed smooth arithmetic geometry permits exponentially distant actual grading holes}
\paragraph{Statement.}
For every integer $L\ge2$ equip the same smooth scheme
$X=\mathbb A^2_{\mathbb Z}$ with the $\mathbb G_m$ action
\[
t\cdot(x,y)=(t^Lx,t^{L+1}y).
\]
Its graded coordinate algebra has support
$S_L=\langle L,L+1\rangle$, largest missing nonnegative degree
$L^2-L-1$, conductor $L(L-1)$, and exactly $L(L-1)/2$ holes.
These are actual missing graded pieces, rather than negative
off-residue evaluations of quasipolynomial constituents.
Nevertheless for every $L$ the underlying scheme has point counts
$q^{2r}$ over $\mathbb F_{q^r}$, global zeta $\zeta(s-2)$,
and compactly supported Weil cohomology consisting of
$\mathbb Q_\ell(-2)$ in degree four. Even the induced
$\mathbb G_m$ action on that cohomology is trivial for every $L$.
Thus these arithmetic and cohomological data cannot bound the
grading conductor polynomially in the binary size of the weights.

\paragraph{Proof of the exact holes and the easy counting rule.}
The monomials $x^ay^b$ are a homogeneous basis, so
\[
H_L(z)=\frac1{(1-z^L)(1-z^{L+1})},\qquad
h_L(n)=\#\{(a,b)\in\mathbb N^2:La+(L+1)b=n\}.
\]
Write $n=QL+r$ with $0\le r<L$. Every possible $b$ has the
form $r+jL$ with $j\ge0$, and its associated $a$ is
$Q-r-j(L+1)$. Consequently
\[
h_L(n)=
\begin{cases}
0,&Q<r,\\
1+\left\lfloor\dfrac{Q-r}{L+1}\right\rfloor,&Q\ge r.
\end{cases}                                                \tag{G017.1}
\]
In residue class $r$ the missing degrees are exactly
$r,\ L+r,\ldots,(r-1)L+r$ if $r>0$, and there are none
if $r=0$. There are therefore $\sum_{r=1}^{L-1}r=L(L-1)/2$
holes. Their largest element is attained for $r=L-1$ and
is $(L-2)L+(L-1)=L^2-L-1$. Every larger degree is present,
which proves the conductor assertion. Formula (G017.1)
decides membership and computes the dimension using a constant
number of integer divisions on $O(\log(n+1)+\log L)$ bits.
It also gives a positive two-variable polytope formula with
coefficients of polynomial bitlength. A large hole depth here
does not make this positivity problem hard.

\paragraph{Proof of arithmetic and cohomological indistinguishability.}
The underlying scheme is independent of the action; its point
count is $q^{2r}$. Taking the exponential generating series
gives $(1-q^2T)^{-1}$, and multiplying over primes gives
$\zeta(s-2)$. The standard compact-support calculation for
affine space gives the displayed Tate object. To make the
classical inputs explicit, compute $H^*(\mathbb P^1)$ over
$\mathbb C$ by SGA 4, XVI, Theorem 4.1, specialize the constant
prime-to-characteristic coefficients by XVI, Corollary 2.2,
and apply the open--closed exact sequence for
$\mathbb A^1\subset\mathbb P^1$ with complement one point.
The result is $H_c^2(\mathbb A^1)=\mathbb Q_\ell(-1)$ and
zero in other degrees. SGA 4, XVII, Theorem 5.4.3 then gives
the affine-plane calculation by K\"unneth. The Tate
normalization is exactly SGA 4, XVIII, (2.8.1). The paragraph
following that formula states that connected algebraic groups
acting on affine space act trivially on this top trace sheaf;
in particular every one of the above $\mathbb G_m$ actions
induces the same trivial action on $H_c^4$. All these primary
locators were opened in this run. Since specifying $(L,L+1)$ takes $\Theta(\log L)$ bits, its quadratic conductor is
exponential in the grading description length.\hfill$\square$

\paragraph{Geometric dictionary and consequences.}
This is already a theorem about a fixed smooth normal scheme
over $\mathbb Z$ with varying diagonal group-scheme actions.
SGA 3, I, Proposition 4.7.3 and Corollary 4.7.3.1 identify
the actions with the coordinate-ring gradings. The saturated
support is $\mathbb N$: the group generated by $L,L+1$ is
$\mathbb Z$ and their rational cone is $\mathbb R_{\ge0}$.
For the associated semigroup curve
$\operatorname{Spec}\mathbb Z[t^L,t^{L+1}]$, the finite
normalization is $\operatorname{Spec}\mathbb Z[t]$ and the
conductor ideal in $\mathbb Z[t]$ is
$t^{L(L-1)}\mathbb Z[t]$. Indeed $t$ is integral because
$t^L$ belongs to the subring, the two fraction fields agree
because $t=t^{L+1}/t^L$, and $\mathbb Z[t]$ is integrally
closed. The conductor calculation follows termwise from the
already proved description of the holes. This semigroup
curve is an associated encoding of the support; it is not
the smooth scheme $X$ itself.

A makes the nontrivial zeros of the common arithmetic zeta
lie on $\operatorname{Re}s=5/2$, but it supplies no missing
weight spaces and changes none of the holes. A fixed finite
pure Weil realization of an \emph{ungraded} count, even with
its induced connected-group action, therefore cannot stand
in for the grade-indexed trace or witness objects needed by
NSRH and PH1. The two-variable formula here already has
the desired positive counting form; this example refutes a
proposed inference about conductor bounds, not PH1 itself
and not a complexity conjecture about actual plethysm.

\subsection{NN-G018: the cubic pole forces exponential rank in any stretching-as-iteration realization}
\paragraph{Statement.}
Let $f_n(r)$ be the ordinary cubic plethysm stretching sequence
in P-F31, with $n\ge4$ and $M=2^{n-1}$. Suppose that over a
characteristic-zero field there are a finite-dimensional vector
space $V_n$ and endomorphisms $A_n,F_n$ such that
\[
f_n(r)=\operatorname{Tr}(A_nF_n^r\mid V_n)
\]
for all sufficiently large integers $r$. Then
\[
\dim V_n\ \ge\ \varphi(M)=2^{n-2}.
\tag{G018.1}
\]
No commutation between $A_n$ and $F_n$ is assumed, and Jordan
blocks are allowed. The same bound applies to the sum of the
dimensions of a finite complex when alternating traces with
insertions are used. This is a rank bound for interpreting the
stretching index as iteration of one operator; it is not a
bound on descriptions of a family of separate graded objects.

\paragraph{Proof and a sharp general linear-algebra formulation.}
For any $d$-dimensional $V$ one has the formal identity
\[
\sum_{r\ge0}\operatorname{Tr}(AF^r)t^r
=\operatorname{Tr}\bigl(A(I-tF)^{-1}\bigr).
\]
The adjugate formula shows that its denominator divides
$\det(I-tF)$, which has degree at most $d$. Altering finitely
many initial coefficients adds a polynomial and changes no
finite poles or their orders. Therefore the total degree of
the reduced denominator of any rational sequence realized
eventually in this manner is at most $d$.

For completeness, this is the exact minimum dimension for
eventual realizations over an algebraically closed
characteristic-zero field. Write the partial fraction expansion
of the rational generating function after discarding its
polynomial part. A pole of order $e$ at $t=\lambda^{-1}$
contributes a sequence $\lambda^rP(r)$ with $\deg P=e-1$.
Use one Jordan block $F=\lambda I+N$ of size $e$. The identity
\[
F^r=\sum_{j=0}^{e-1}\binom rj\lambda^{r-j}N^j
\]
shows that the linear functionals $\operatorname{Tr}(AF^r)$
span all such sequences: choose entries of $A$ on the $j$th
lower diagonal to prescribe $\operatorname{Tr}(AN^j)$, and
use that the polynomials $\binom rj$ form a basis. These
diagonals are disjoint, so the choices are independent.
Taking the direct sum over poles gives dimension equal to
the sum of their orders. Together with the lower bound this
proves exactness for eventual, inserted-trace realizations.

P-F31 asserts a genuine simple pole at every primitive
$M$th root. There are $\varphi(M)=M/2$ such roots, so the
reduced denominator has at least that degree and (G018.1)
follows. One could equivalently start from a single primitive
pole, since the Hilbert series has rational coefficients and
its Galois conjugates are also poles. For an alternating
inserted trace replace the collection of spaces by their
direct sum, $F$ by its block diagonal operator and $A$ by
the block diagonal insertions with their signs. The same
denominator argument applies.\hfill$\square$

\paragraph{Ordinary traces are more restrictive.}
For $A=I$, Jordan blocks contribute only their sizes times
$\lambda^r$ to the trace. Thus a finite sum, difference or
cohomological alternating sum of ordinary traces has only
simple poles in its generating series. Any quasipolynomial
having a nonconstant polynomial constituent has a pole of
order greater than one at some root of unity: decompose its
periodic coefficient functions into their finite Fourier
series, obtaining $\sum_\rho P_\rho(r)\rho^r$; a nonconstant
constituent forces some $P_\rho$ to have positive degree.
Its pole of order $1+\deg P_\rho$ cannot cancel against
terms based at other roots. Such a sequence cannot be an
ordinary trace sequence on any fixed finite complex.
This last statement is a general criterion; the rank bound
above needs only the specifically cited P-F31 simple poles.

\paragraph{Geometric dictionary and exact limit of the conclusion.}
The input is P-F31, theorem \texttt{actualcubicpole}, part two,
in \texttt{positivity\_actual\_cubic\_pole.tex}; its family is
the Cartan multiplicity algebra of cubic $SL_n$ invariants.
Its variable $t$ records the grading in SGA 3, I, 4.7.3.1.
To reinterpret $f_n(r)$ as an extension-field trace in the
SGA 5, III, 4.7 framework requires a new identification of
that grading index with the power of Frobenius. If such an
identification is proposed using one finite cohomological
carrier, (G018.1) bounds its total rank. Neither a Tate twist,
an inserted correspondence nor virtual cancellation removes
the pole count; insertions and signs were already allowed.
It is essential to distinguish this from the Tate--Artin
realizations of individual Laurent-polynomial structure
constants in NN-G006--NN-G011, where the field size and
the basis-word or stretching parameter are separate variables.

The rank obstruction is unconditional relative to the cited
refutation and is unaffected by A or conditional M. It
rules out a polynomial-rank iteration carrier on this actual
plethysm family, but leaves short tensor descriptions,
varying graded carriers, succinct Hilbert rational functions,
and PH1 positive witness formulas open. No arithmetic-circuit
or Boolean lower bound follows from it alone.

\subsection{NN-G019: a finite optimum theorem for the integral-ratio dimension-ring class}
\paragraph{Precisely delimited class and statement.}
Consider dimension rings with the table shape (G011.2), but
allow the parameters to vary subject to the following conditions:
\[
\begin{gathered}
\delta_1=\delta_2=\delta\in\mathbb Z[D]
   \text{ is monic of degree two},\qquad \deg W\le2,\\
G_i=k_iW\delta,\qquad S=k_1k_2W\delta^2,
\qquad k_1,k_2\in\mathbb Z_{>0},\\
\Lambda_i\in\mathbb Z[D],\qquad
\Lambda_i\ge\lambda_i,\quad\Lambda_i\equiv\lambda_i\pmod2
\quad\text{coefficientwise in }q,
\end{gathered}                                              \tag{G019.1}
\]
where $D=q+2+q^{-1}$, $\lambda_1=D$ and
$\lambda_2=(D-2)^2$. As always, every table coefficient
must dominate the absolute value of the original coefficient
and have the same parity. There is no imposed degree or
coefficient cutoff on $\Lambda_i$ or on any other parameter.
Within this class the minimum number of cancelling pairs is
exactly $7276$. There are exactly two minimizing parameter
sets. Both have
\[
k_1=2,\quad k_2=1,\quad
\delta=D^2-3D+2,\quad
\Lambda_1=D,\quad\Lambda_2=3D^2-10D+4,
\]
and their two choices of $W$ are
\[
W=7D-8\qquad\text{or}\qquad W=4D^2-13D+8.
\tag{G019.2}
\]
The lower bound is a finite exact-arithmetic certificate
after the range reduction below. It is not a global optimum
over all dimension rings or over all polynomial block parameters.

\paragraph{Proof: algebraic constraints and finite bounds.}
The scalar-block identities force, with $\{i,j\}=\{1,2\}$,
\[
U_i=\frac{k_iW-D}{k_j},\qquad
V_i=\frac{k_iDW-\Lambda_i}{k_j}.                 \tag{G019.3}
\]
The coefficient constraints include
$W\ge D$, $\delta\ge |(D-1)(D-4)|$, $U_i\ge D$,
$V_i\ge K_i$, where $K_1=5D-4$ and
$K_2=2D^2-5D+4$, always with parity congruence.
Modulo two, the equation $k_2V_1=k_1DW-\Lambda_1$
therefore becomes
$k_2D=k_1D^2+D$. The Laurent polynomials $D$ and $D^2$
are linearly independent over $\mathbb F_2$, so
$k_1$ is even and $k_2$ is odd.

Put $w=W(4)$ and $d=\delta(4)$, which are the sums of
the Laurent coefficients. Let $T$ be the sum of all
nonunit dimension coefficients, evaluated at $q=1$.
The word-pair multiplicities computed in NN-G011 give
\[
\begin{split}
T={}&k_1k_2wd^2+20wd(k_1+k_2)+200w\\
 &+20\bigl(U_1(4)+U_2(4)+V_1(4)+V_2(4)\bigr)\\
 &+d\bigl(52+14\Lambda_1(4)+14\Lambda_2(4)\bigr).
\end{split}                                                \tag{G019.4}
\]
The cancelling-pair score is $(T-4424)/2$. The already
proved incumbent NN-G011 means that, to prove the optimum,
it suffices to examine $T\le18976$. Since $w,d\ge4$,
$U_i(4)\ge4$, $V_i(4)\ge16$, and $\Lambda_i(4)\ge4$,
\[
T\ge k_1k_2wd^2+20wd(k_1+k_2)+200w+800+164d.
\tag{G019.5}
\]
In particular $T\ge472w+1456$ and
$T\ge8d^2+404d+1600$. Both $w,d$ are even, so
$w\le36$ and $d\le26$. Also
$T\ge320(k_1+k_2)+2256$ gives $k_1+k_2\le52$.
These bounds are consequences of the desired score;
they are not experimental search cutoffs.

All possible $\delta$ and $W$ are now given by the finite
nonnegative integer parameters
\[
\begin{aligned}
\delta&=(q^2+q^{-2})+(1+2a)(q+q^{-1})+2b,
       &4+4a+2b&\le26,\\
W&=c_2(q^2+q^{-2})+c_1(q+q^{-1})+c_0,
       &c_0+2c_1+2c_2&\le36,
\end{aligned}
\]
where $a,b\ge0$, $c_0\ge2$ is even, $c_1\ge1$ is
odd, and $c_2\ge0$ is even. These formulas cover every
monic quadratic $\delta$ and every degree-at-most-two $W$
with their required Laurent coefficient inequalities and parity.

\paragraph{Proof: eliminate the unbounded $\Lambda_i$.}
Fix the bounded parameters and consider cell $i$.
Let $C_i=k_iDW-k_jK_i$. The requirement
$V_i\ge K_i$ with $V_i\equiv K_i\pmod2$ is equivalent to
\[
\Lambda_i\le C_i,\qquad
\Lambda_i\equiv C_i\pmod{2k_j}
\quad\text{coefficientwise}.
\]
For each exponent $e$, put $l_e=[q^e]\lambda_i$ and
$c_e=[q^e]C_i$. A possible coefficient must be at least
$l_e$, so its unique smallest value is
\[
L_e=l_e+((c_e-l_e)\bmod 2k_j),
\tag{G019.6}
\]
where the residue is in $\{0,\ldots,2k_j-1\}$.
There is no feasible value if the parity differs from
$l_e$ or if $L_e>c_e$. Otherwise the coefficients $L_e$
give a symmetric integral Laurent polynomial and hence an
element of $\mathbb Z[D]$. Coefficients outside the support
of $C_i$ and $\lambda_i$ must be zero; thus this step also
proves the required degree bound, instead of assuming one.

Increasing $\Lambda_i(4)$ while fixing the other bounded
parameters changes $T$ with strictly positive slope
$14d-20/k_j\ge36$. Consequently only the componentwise
smallest choice (G019.6) can minimize the score. This
eliminates all remaining unbounded choices.

\paragraph{Proof: the exact finite certificate.}
Enumerate the nonnegative integer tuples just displayed,
the positive even $k_1$ and positive odd $k_2$ with sum
at most $52$, and retain (G019.5)$\le18976$. There are
$1744$ such tuples. For each, divide (G019.3) coefficientwise
and test integrality, domination and parity for $U_i$;
test $G_i$ and $S$ against the true parameters; then apply
(G019.6). Exactly the following $21$ tuples survive those
tests, with $\Lambda_i$ chosen minimally:
\[
\begin{array}{rr|rr|rrr|r}
k_1&k_2&a&b&c_0&c_1&c_2&\text{pairs}\\\hline
2&1&0&1&6&3&4&7276\\
2&1&0&1&6&7&0&7276\\
2&3&0&0&10&5&0&7656\\
2&1&0&0&14&7&0&8288\\
2&1&0&1&10&3&4&9040\\
2&1&0&1&10&7&0&9040\\
2&1&0&2&6&3&4&9312\\
2&1&0&2&6&7&0&9312\\
2&1&1&0&6&3&4&9312\\
2&1&1&0&6&7&0&9312\\
2&1&0&0&10&7&4&9732\\
2&1&0&0&10&11&0&9732\\
2&1&0&0&18&3&4&9732\\
2&1&0&0&18&7&0&9732\\
2&1&0&0&6&11&4&11176\\
2&1&0&0&6&15&0&11176\\
2&1&0&0&14&3&8&11176\\
2&1&0&0&14&7&4&11176\\
2&1&0&0&14&11&0&11176\\
2&1&0&0&22&3&4&11176\\
2&1&0&0&22&7&0&11176
\end{array}
\]
This finite calculation is reproduced by
\texttt{verify\_dimension\_ratio\_optimum.py}, whose complete
source is included. It uses only integer Laurent dictionaries
and imports no supplied search or table implementation. Its
JSON certificate lists every surviving tuple, both minimizing
tables, their actual $\Lambda_i$, and the bottom and cell
masses. A separate search using the supplied polynomial
arithmetic gives the same counts and minimum.

The first minimizing table is the new $W$ in (G019.2),
and the second is NN-G011. For both, (G019.6) gives the
displayed $\Lambda_i$. The table conditions were checked
coefficientwise, and associativity follows from the same
block embedding as NN-G011: the four images in each matrix
block are independent since $\Lambda_i(D^2-\Lambda_i)\ne0$,
and (G019.3) gives exactly the required scalar-block products.
Thus both minima are attained by actual dimension rings.
Strict positivity of the slopes above proves that no
nonminimal $\Lambda_i$ gives another minimizer. This proves
the stated optimum and the two-parameter-set classification.
\hfill$\square$

\paragraph{Geometric translation and smaller remaining gate.}
By NN-G010 both minimizing rings give additive monoidal
categories of kernels on a fixed ten-point scheme with
Tate--Artin coefficient objects. The score counts the added
regular $C_2$ pairs in the structure-constant objects,
coefficient by coefficient. The relevant actual-object and
split-$K_0$ distinction is SGA 3, I, 4.7.2 and archive cote 2,
scan 17; effectivity is not equality in a virtual representation
ring. This optimization selects no canonical geometric functor.

Any improvement below $7276$ must now leave at least one
specific part of (G019.1): use unequal cell $\delta_i$,
a nonmonic or higher-degree $\delta$, $\deg W>2$,
nonintegral or nonconstant ratios $G_i/(W\delta)$,
allow $\Lambda_i$ below the true $\lambda_i$ in some
coefficients while compensating elsewhere in the table,
or change the table shape. These are genuine remaining
possibilities. The result narrows the search gate but
does not prove global minimality, canonical positivity,
PH1, or anything about rank four by itself.

\subsection{NN-G020: normal toric rays with conductor at most two and hard first coefficient}
\paragraph{Statement and input representation.}
There is a polynomial-time construction from a $3$-CNF formula $F$
on $n\geq1$ declared variables and $m$ clauses to an integral weight
matrix $A_F$, a primitive character $b$, and the diagonal action of
the split torus $T=\mathbb G_m^{n+m}$ on
\[
 S=\mathbb Z[X_1,Y_1,\ldots,X_n,Y_n,Z_1,\ldots,Z_m].
\]
The ray algebra
\[
 R_F=\bigoplus_{k\geq0}S_{kb}
 \tag{G020.1}
\]
is a normal, finitely generated, flat $\mathbb Z$-algebra and is
exactly the invariant ring $S^{\ker\chi_b}$. Its graded pieces have
canonical monomial bases and satisfy
\[
 \operatorname{rank}_{\mathbb Z}(R_F)_1=\#F,
 \qquad
 \{k:(R_F)_k\ne0\}=
 \begin{cases}
  \mathbb N,&F\text{ is satisfiable},\\
  \{0,2,3,4,\ldots\},&F\text{ is unsatisfiable}.
 \end{cases}
 \tag{G020.2}
\]
Here $\mathbb N$ includes zero. In particular the \emph{actual ray
conductor}, defined as the least $c\geq0$ with every integer $k\geq c$
in the support, is either zero or two. A rational point of the
degree-one polytope and an integral point in every degree $k\geq2$
are explicitly known independently of satisfiability.

The input for the complexity statement is this polynomial-size
\emph{diagonal-torus presentation} $(A_F,b)$, or equivalently the
inequalities below. No polynomial bound on a displayed generating
set or relations for the invariant algebra is asserted. In this
presentation, deciding whether $(R_F)_1\ne0$ is NP-complete, and
computing its rank is \#P-complete. E gives the exact subexponential
counting bounds stated below, with the original number $n$ of
Boolean variables. A places the nontrivial zeros of the common
ambient arithmetic zeta on one shifted line, while that zeta
contains no information about (G020.2).

\paragraph{Construction and homogeneous lattice-point model.}
Write every clause as three literal positions; repeating a literal
pads a shorter clause without changing any assignment. For each
clause $C$, let
\[
 L_C(x,y)=\sum_{\text{positions in }C}
 \begin{cases}x_i,&\text{the position is the literal }x_i,\\
 y_i,&\text{the position is the literal }\neg x_i.
 \end{cases}
\]
Repeated occurrences are counted with multiplicity. Define $A_F$
by the equations
\[
 x_i+y_i=k\quad(1\leq i\leq n),\qquad
 L_C(x,y)-s_C=k\quad(1\leq C\leq m),
 \tag{G020.3}
\]
and put $b=(1,\ldots,1)\in\mathbb Z^{n+m}$. The columns of $A_F$
are the weights of $X_i,Y_i,Z_C$, in that order. Its entries have
absolute value at most three, and its construction is polynomial
in the formula encoding. In the group-scheme action, a monomial
with exponent vector $z=(x,y,s)$ has character $A_Fz$; this is a
statement about comodules, not merely about rational points of
the torus over a small finite field.

At $k=1$, nonnegative integral solutions of the first equations
have $x_i\in\{0,1\}$ and $y_i=1-x_i$. The remaining equations
uniquely force $s_C=L_C(x,y)-1$. This is nonnegative exactly
when clause $C$ is true. Hence the map from satisfying assignments
to exponent vectors is a bijection. It proves the first equality
in (G020.2), including multiplicities and declared variables that
do not occur in a clause.

Set
\[
 P_F=\{z\in\mathbb R_{\geq0}^{2n+m}:A_Fz=b\}.
\]
This is a bounded rational polytope: $0\leq x_i,y_i\leq1$ and
$0\leq s_C\leq2$. It always contains the explicitly specified
point $x_i=y_i=s_C=1/2$. Its dilations give the exact homogeneous
formula
\[
 \operatorname{rank}_{\mathbb Z}(R_F)_k
   =|kP_F\cap\mathbb Z^{2n+m}|\quad(k\geq0).
 \tag{G020.4}
\]
At $k=0$, the first equations force $x=y=0$ and then $s=0$,
so the formula holds there as well. In particular $b$ is in the
weight cone for every formula; cone feasibility is already settled.

For every integer $k\geq2$, choose
\[
 x_i=\lfloor k/2\rfloor,\qquad
 y_i=\lceil k/2\rceil,\qquad
 s_C=L_C(x,y)-k.
 \tag{G020.5}
\]
Every clause has three positions, so
$L_C(x,y)\geq3\lfloor k/2\rfloor\geq k$ for $k\geq2$.
Thus these are integral nonnegative solutions for every such
degree. Combining this with the degree-zero and degree-one
calculations proves the full support formula and its conductor.

\paragraph{Exact invariant-ring identification and normality.}
The linear functional $\beta$ on the character lattice takes the
value $3m+1$ on every variable-equation row and $-1$ on every
clause-equation row. Each column for $X_i$ or $Y_i$ therefore has
$\beta$-value $3m+1$ minus its total number of literal occurrences,
which is at least one. Each slack column has $\beta$-value one.
Moreover
\[
 \beta(b)=n(3m+1)-m>0.
 \tag{G020.6}
\]
Consequently no nonzero monomial has weight zero, and no monomial
has weight $kb$ for $k<0$. This also proves finite rank of every
graded piece directly from a positive bound on its exponents.

Since $b$ is primitive, it extends to an integral basis of the
character lattice $M=\mathbb Z^{n+m}$. Thus
$K=\ker(\chi_b:T\to\mathbb G_m)$ is a split connected torus and
its character lattice is $M/\mathbb Zb$. Restricting the monomial
grading to $K$ gives, over $\mathbb Z$ itself,
\[
 S^K=\bigoplus_{k\in\mathbb Z}S_{kb}
     =\bigoplus_{k\geq0}S_{kb}=R_F.
 \tag{G020.7}
\]
The second equality uses (G020.6); omitting that step would leave
the proposed ray/invariant-ring identification unjustified.

For completeness, finite generation is the elementary rational-cone
argument. Let $H$ be the set of nonnegative integral exponent vectors
whose weights lie in $\mathbb Zb$. It is the intersection of a
rational polyhedral cone with a lattice. Choose finitely many
integral vectors generating that cone over $\mathbb R_{\geq0}$.
For a lattice vector in the cone, subtract the integer parts of
its coefficients in a nonnegative linear combination of those
vectors. The remainder is a lattice vector in the bounded region
where all those coefficients lie in $[0,1)$. There are finitely
many such remainders. They, together with the chosen vectors,
generate $H$. Thus $R_F=\mathbb Z[H]$ is finitely generated.
This is Gordan's lemma in its lattice form; see Cox--Little--Schenck,
\emph{Toric Varieties}, Proposition 1.2.17, printed p.~30.

To prove normality without an unmentioned quotient hypothesis,
let $f\in\operatorname{Frac}(S^K)$ be integral over $S^K$.
Its monic integral equation also has coefficients in $S$, so it
is integral over $S$. The polynomial ring $S$ over $\mathbb Z$
is integrally closed, hence $f\in S$. As a ratio of two invariant
elements, $f$ is invariant as a rational function, and therefore
$f\in S^K$. This proves that $R_F$ is integrally closed. Finally,
its monomial basis makes it a free $\mathbb Z$-module, so it is
flat. Normality of this algebra is entirely consistent with the
single missing degree one: saturation of its exponent semigroup
does not mean saturation of its image under the degree map.

\paragraph{Complexity consequences, with E's original parameter.}
An exponent vector of degree one has a polynomial-size encoding,
and (G020.3) verifies it in polynomial time. The established
bijection is a parsimonious reduction from \#3SAT and its
positivity version. Thus the stated NP- and \#P-completeness follow
from the classical verifier encoding, which can also be made
explicit here. Compile a polynomial-time Boolean verifier into a
polynomial-size circuit with AND and NOT gates. For each AND gate
$z=x\wedge y$ impose
$(\neg z\vee x)$, $(\neg z\vee y)$, and
$(z\vee\neg x\vee\neg y)$; for each NOT gate $z=\neg x$ impose
$(z\vee x)$ and $(\neg z\vee\neg x)$. Impose the output bit by
a unit clause, and repeat literals to pad each clause to three
positions. For every input witness, all gate bits have exactly
one consistent extension, accepted precisely when the verifier
accepts. Thus counting these assignments is parsimoniously as
hard as counting the witnesses of any \#P verifier, and the same
construction proves NP-hardness of positivity. Membership in
\#P and NP uses the degree-one witness already given. This proof
does not bound the new formula's variable count by the number
of input witness bits alone; its gate bits are counted as well.

More concretely, a deterministic polynomial-time positivity algorithm
on these presented normal toric rays would imply $P=NP$.
Polynomial-size certificates for their degree-one vanishing,
checkable in deterministic polynomial time and complete on all
these inputs, would imply $NP=coNP$: the same certificates would
put $3$-UNSAT in NP. These are implications, not assertions of
the separations $P\ne NP$ or $NP\ne coNP$.

Applying the designated E Theorem 1.1 directly to $F$, not to an
inflated auxiliary-variable formula, computes the degree-one
rank in
\[
 (n+m)^{O(1)}2^{O(n\log\log n/\sqrt{\log n})}
 \quad\text{bit operations and polynomial space}.
 \tag{G020.8}
\]
For fixed $D$ and $m\leq Dn$ the exponent improves to
$O_D(n/\sqrt{\log n})$, with polynomial prefactor and space.
The theorem is stated for $n\geq20$; the finitely many smaller
values are handled by direct enumeration. There is no hidden
replacement of $n$ by $2n+m$ in this use of E.

\paragraph{Geometric dictionary and the exact failure of zeta recovery.}
The torus action and its multigraded comodule are the objects of
SGA 3, I, 4.7.2--4.7.3 and Corollary 4.7.3.1. Equation (G020.7)
is an actual affine quotient by the indicated split subtorus;
the finite generation and normality proofs above make its scheme
meaning explicit. The ray grading also gives $\operatorname{Proj}R_F$
and its associated graded-module sheaves in the sense of EGA II,
2.4--2.5, especially 2.4.7 and 2.5.2. No identification with a
plethysm or Kronecker multiplicity algebra is asserted.

Put $N=2n+m$. The underlying \emph{ambient} scheme
$\operatorname{Spec}S=\mathbb A^N_{\mathbb Z}$ is independent of
the formula once $n,m$ are fixed. For every prime power $q$ and
integer $r\geq1$ it has $q^{Nr}$ points over $\mathbb F_{q^r}$.
Its only nonzero compact-support $\ell$-adic cohomology group
is $H_c^{2N}=\mathbb Q_\ell(-N)$, and the connected torus action
induces the trivial action on that group. The precise trace and
connected-action locator is SGA 4, XVIII, 2.8.1 and the immediately
following paragraph; the affine-space calculation also follows
from XVII, 5.4.3 and the curve trace calculation XVIII, 1.1.6--1.1.9.
Consequently
\[
 Z(\mathbb A^N_{\mathbb F_q},T)=(1-q^NT)^{-1},\qquad
 \zeta_{\mathbb A^N_{\mathbb Z}}(s)=\zeta(s-N).
 \tag{G020.9}
\]
Under A its nontrivial zeros lie on $\Re s=N+1/2$. These common
point counts and this rank-one pure Weil object do not recover
the degree-one rank, despite the normal ray algebra, known rational
half-point, positive monomial rule and conductor bound two.
For an explicit smallest comparison, the formulas
$(x\vee x\vee x)\wedge(x\vee x\vee x)$ and
$(x\vee x\vee x)\wedge(\neg x\vee\neg x\vee\neg x)$
have degree-one ranks one and zero and the same ambient
$\mathbb A^4_{\mathbb Z}$.

This last assertion concerns the ungraded ambient realization;
no formula-independent zeta or cohomology is claimed for
$\operatorname{Spec}R_F$ itself, or for grade-indexed trace objects.
Likewise the conductor here is an actual support conductor,
not GCT VI's strict off-residue saturation index. The construction
does not contradict a conjecture confined to actual plethysm or
Kronecker families. It proves that normality, a homogeneous
lattice-point model, known moment-cone membership, a uniformly
tiny actual conductor, and a finite pure ambient realization
still leave an NP-hard first-coefficient gate in this explicit
geometric family. Neither A nor the unrelated conditional M
removes that gate; E supplies precisely (G020.8).
\hfill$\square$

\subsection{NN-G021: signed auxiliary parameters lower the exact rank-three optimum}
\paragraph{Statement and the larger class.}
The auxiliary cell parameter $\Lambda_i$ is not itself a structure
constant of the rank-three table. Requiring it to dominate the
corresponding trace parameter is therefore an extra restriction.
Removing that restriction gives an associative effective dimension
ring with
\[
 7064\text{ cancelling pairs in total},\qquad
 6652\text{ in the bottom cell}.                         \tag{G021.1}
\]
An explicit choice keeps $\delta,W,G_1,G_2,S,U_1,U_2$ from NN-G011,
and replaces the four remaining parameters by
\[
 \Lambda_1=D-2,\quad \Lambda_2=3D^2-10D,\qquad
 V_1=14D^2-17D+2,\quad V_2=2D^2+D.                      \tag{G021.2}
\]
Although $\Lambda_2$ has a negative constant Laurent coefficient,
all actual multiplication coefficients are effective.

Moreover $7064$ is the exact minimum in the following class:
the same closed table shape, a common monic quadratic
$\delta\in\mathbb Z[D]$, a polynomial $W\in\mathbb Z[D]$ of
degree at most two, and arbitrary \emph{positive rational constant}
ratios $G_i/(W\delta)$; the $\Lambda_i$ may be arbitrary signed
polynomials in $\mathbb Z[D]$. Only the actual table coefficients
must be integral, dominate the absolute trace coefficients, and
have the required parity. There are eight minimizing parameter
sets in this class, described below. No optimum over unrestricted
dimension rings is claimed.

\paragraph{The explicit improvement.}
Use Laurent coefficient lists indexed by nonnegative exponents,
with negative exponents supplied by symmetry. Thus
\[
 \delta=(2,1,1),\quad D-2=(0,1),\quad
 3D^2-10D=(-2,2,3).
\]
The actual cell products are
\[
 \delta\Lambda_1=(2,3,1,1),\qquad
 \delta\Lambda_2=(6,7,6,5,3).                           \tag{G021.3}
\]
The absolute trace targets for these products are, respectively,
\[
 X_1=(2,1,1,1),\qquad X_2=(2,3,2,1,1).                 \tag{G021.4}
\]
They are dominated with the correct parity. The unchanged $G_i,S,W,U_i$
already pass their coefficient conditions by NN-G011. The new
$V_1,V_2$ equal the previous ones plus two and still dominate their
trace targets with unchanged parity. The closed table identities
remain valid because $V_i=(k_iDW-\Lambda_i)/k_j$ with
$(k_1,k_2)=(2,1)$, so decreasing $\Lambda_1$ by two and
$\Lambda_2$ by four increases both $V_i$ by two.

All the table's within-cell coefficients are $\delta D^\alpha$
or $\delta D^\alpha\Lambda_i$; its bottom coefficients are
nonnegative powers of $D$ times the parameters just checked.
Thus these base checks prove coefficientwise domination, parity
and effectivity everywhere. Associativity follows from the same
scalar-block and two-matrix-block embedding used in NN-G011;
the conditions $\Lambda_i(D^2-\Lambda_i)\ne0$ hold for (G021.2).
Alternatively the included verifier checks all $1000$ triples
directly with integer Laurent dictionaries.

The total nonunit coefficient mass is now $18552$, divided into
bottom mass $17400$ and cell mass $1152$. The true absolute masses
are $4096$ and $328$. Consequently the cancelling-pair counts are
$(18552-4424)/2=7064$ and $(17400-4096)/2=6652$.
The bottom count rises by forty while the total count falls by
$212$: forcing positivity on an auxiliary parameter had concealed
a less expensive effective multiplication table.

\paragraph{Why the larger optimization class has a finite bound.}
Write $W=gW_0$, where $g>0$ is its coefficient content and $W_0$
is primitive. Positivity and $W\equiv D\pmod2$ imply that $g$ is
odd and that
\[
 W_0=c_0+c_1(q+q^{-1})+c_2(q^2+q^{-2}),
 \quad c_0\ge2\text{ even},\ c_1\ge1\text{ odd},\
 c_2\ge0\text{ even},\ \gcd(c_0,c_1,c_2)=1.
 \tag{G021.5}
\]
The common quadratic has the exact form
\[
 \delta=2b+(1+2a)(q+q^{-1})+(q^2+q^{-2}),\qquad a,b\ge0.
 \tag{G021.6}
\]
Set $d=\delta(4)=4+4a+2b$ and $w_0=W_0(4)=c_0+2c_1+2c_2$;
evaluation at $D=4$ is evaluation at $q=1$.

The polynomial $q^2\delta$ is monic and primitive. Gauss's content
lemma shows that $W_0\delta$ and $W_0\delta^2$, after harmless
Laurent shifts, are primitive. Integrality of $G_i=k_iW\delta$
therefore forces $a_i:=gk_i$ to be a positive integer. Integrality
of $S=G_1G_2/W$ forces $g\mid a_1a_2$. The trace congruences
for $G_1,G_2$ then force $a_1$ even and $a_2$ odd: modulo two,
$W_0\delta=D^2(D-1)$ is nonzero, the first trace target is zero,
and the second equals that product. Thus
\[
 \begin{aligned}
 G_i&=a_iW_0\delta,\\
 S&=(a_1a_2/g)W_0\delta^2,\\
 U_i&=g(a_iW_0-D)/a_j,\\
 V_i&=g(a_iDW_0-\Lambda_i)/a_j\quad(j\ne i),
 \end{aligned}                                         \tag{G021.7}
\]
where $a_1a_2/g$ is a positive even integer. This accounts for
all rational constant ratios in the stated class, including
nonintegral ones.

Let $T$ be the total nonunit dimension mass. The exact table
count used in NN-G019 remains valid even for signed $\Lambda_i$:
\[
 T=S(4)+20(G_1(4)+G_2(4))+200gw_0
   +20\sum_i(U_i(4)+V_i(4))
   +52d+14d\sum_i\Lambda_i(4).
 \tag{G021.8}
\]
Since $\delta\Lambda_i\ge X_i$ and the masses of (G021.4)
are eight and sixteen, the last term is at least $336$.
Also $\sum_i(U_i(4)+V_i(4))\ge40$. Any score at most the
already proved $7276$ therefore satisfies $T\le18976$ and
\[
 (a_1a_2/g)w_0d^2+20w_0d(a_1+a_2)
       +200gw_0+1136+52d\le18976.                     \tag{G021.9}
\]
This implies the complete finite ranges
\[
 w_0\le36,\quad d\le30,\quad a_1+a_2\le52,\quad g\le23.
 \tag{G021.10}
\]
For the first bound the left side is at least $472w_0+1344$,
and $w_0$ is even. For the second it is at least
$8d^2+292d+1936$, which exceeds $18976$ at $d=32$.
For the third it is at least $320(a_1+a_2)+2272$.
The last follows already from $200gw_0\ge800g$.
All four lower bounds use only positive terms and the minimum
values $d,w_0\ge4$, $a_1+a_2\ge3$ and $a_1a_2/g\ge2$.

\paragraph{Exact minimization over every signed $\Lambda_i$.}
No degree restriction on $\Lambda_i$ has been hidden. If its
degree exceeded three, its highest Laurent coefficient would
occur with the same sign at the top of $\delta\Lambda_i$ and
with the opposite sign at the top of $V_i$, since
$\deg(DW_0)\le3$. Both cannot be effective. Hence
$\deg\Lambda_i\le3$.

For each $i$, put $a_j'=a_j/\gcd(g,a_j)$. Integrality and parity
of $V_i$ are exactly the coefficientwise congruences
\[
 \Lambda_i\equiv a_iDW_0-a_j'K_i\pmod{2a_j'},\qquad
 K_1=5D-4,\quad K_2=2D^2-5D+4,                      \tag{G021.11}
\]
and domination of $K_i$ gives the upper bounds
\[
 (\Lambda_i)_e\le
 \left\lfloor\frac{g a_i(DW_0)_e-a_j(K_i)_e}{g}\right\rfloor.
 \tag{G021.12}
\]
To check (G021.11), divide the original congruence
$g(a_iDW_0-\Lambda_i)\equiv a_jK_i\pmod{2a_j}$ by
$\gcd(g,a_j)$. The remaining coefficient of $\Lambda_i$
is odd and coprime to the modulus; its inverse is odd, so its
product with $a_j'$ equals $a_j'$ modulo $2a_j'$.

Write $h=2b$, $u=1+2a$ and
$\Lambda_i=L_0+\sum_{e=1}^3L_e(q^e+q^{-e})$.
The requirement $\delta\Lambda_i\ge X_i$, coefficientwise,
is the following complete list, with absent $X_i$ coefficients
interpreted as zero:
\[
 \begin{array}{rcl}
 L_3&\ge&0,\\
 L_2+uL_3&\ge&(X_i)_4,\\
 L_1+uL_2+hL_3&\ge&(X_i)_3,\\
 L_0+uL_1+hL_2+uL_3&\ge&(X_i)_2,\\
 uL_0+(h+1)L_1+uL_2+L_3&\ge&(X_i)_1,\\
 hL_0+2uL_1+2L_2&\ge&(X_i)_0.
 \end{array}                                           \tag{G021.13}
\]
The congruences give $\Lambda_1\equiv D$ and
$\Lambda_2\equiv D^2\pmod2$, so they enforce the required
product parity as well. Conversely that parity forces these
congruences modulo two since $\delta\not\equiv0\pmod2$.
These inequalities bound $L_3$,
then $L_2$, then $L_1$, and finally $L_0$ from below;
(G021.12) bounds each from above. When $h=0$, the last
inequality is simply tested without dividing by $h$.

For fixed $L_3,L_2,L_1$, the coefficient of $L_0$ in the score
mass is $14d-20g/a_j$. Therefore only the least allowed $L_0$
is needed when that coefficient is positive, and only the
greatest when it is negative. If it is zero every allowed
$L_0$ has the same score. This endpoint rule is used only
after all six inequalities and the congruences are imposed.
It is not an assumption that the auxiliary polynomial is
coefficientwise nonnegative.

The complete finite enumeration of (G021.5)--(G021.13),
implemented in \texttt{search\_dimension\_signed\_lambda.py},
has $2572$ tuples after (G021.9), $99$ after the structural
$U,G,S$ tests, and $33636$ checked $\Lambda_i$ endpoints.
Twenty-eight structural tuples admit feasible cell parameters.
Their minimum scores are displayed in the included JSON
certificate; the only ones at or below $7276$ are
\[
 \begin{array}{c|c|c|c|c|c}
 (a_1,a_2,g)&(a,b)&(c_0,c_1,c_2)&\text{score}
       &\#\Lambda_1\text{ choices}&\#\Lambda_2\text{ choices}\\\hline
 (2,1,1)&(0,1)&(6,3,4)&7064&2&3\\
 (2,1,1)&(0,1)&(6,7,0)&7064&1&2
 \end{array}                                           \tag{G021.14}
\]
For either row the first available $\Lambda_1$ is $D-2$.
For the quadratic-$W$ row only, there is also the symmetric
Laurent list $(4,-1,-2,2)$. Both rows allow
$\Lambda_2=(6,-2,3,0)$ or $(-2,2,3,0)$; the quadratic row
also allows $(6,-2,-1,4)$. These yield $2\cdot3+1\cdot2=8$
parameter sets. At every minimum the two slopes are strictly
positive (their numerators in (G021.7) are $64$ and $148$),
so the endpoint rule loses no additional minimizing choices.
The arithmetic source and complete survivor data are included;
the calculation imports no supplied search implementation.

All eight displayed parameter sets have nonzero
$\Lambda_i(D^2-\Lambda_i)$, so the block embedding proves
associativity for their tables. Each dominates the trace
coefficients by the exact tests above. Thus the computed
minimum is attained by actual associative dimension rings.
The explicit table (G021.2) is independently checked by
\texttt{verify\_dimension\_7064.py}, which compares all $132$
positions with the supplied trace constants and checks every
one of the $1000$ associativity triples. This proves both
the improvement and the exact optimum in the stated class.
\hfill$\square$

\paragraph{Geometric meaning and the remaining gate.}
NN-G010 turns each of these effective multiplication tables
into the same type of additive monoidal kernel category on
a ten-point scheme, with effective Artin--Tate multiplicity
objects. Only the actual structure constants specify those
objects. An auxiliary factor such as $\Lambda_2$ need not
be the class of an effective object even when its product
with $\delta$ is. This is the precise distinction between
an effective tensor structure and a chosen polynomial
factorization of its coefficients; no cancellation history
can be read from a supported-zero tag alone.
The comodule dictionary is SGA 3, I, 4.7.2; kernel products
are the construction recorded in archive cote 2, scan 17.
Neither source canonically selects these tables.

The class bound now allows signed $\Lambda_i$ and rational
constant ratios. A score below $7064$ must leave another
specified restriction: a common monic quadratic $\delta$,
degree at most two for $W$, constant ratios $G_i/(W\delta)$,
or the closed table shape. This is a smaller global
optimization gate than NN-G019. It remains a question
about noncanonical realizations, with no consequence for
PH0, PH1 or a rank-four canonical basis.

\subsection{NN-G022: unequal monic quadratic cells force at least 7958 pairs}
\paragraph{Statement.}
Keep the closed rank-three table shape, but allow two distinct monic
quadratic cell parameters $\delta_1,\delta_2\in\mathbb Z[D]$.
Let $W\in\mathbb Z[D]$ have \emph{arbitrary degree}, and suppose
\[
 G_i=k_iW\delta_i\quad(k_i\in\mathbb Q_{>0}),\qquad
 S=G_1G_2/W.
 \tag{G022.1}
\]
Allow arbitrary signed auxiliary $\Lambda_i$. If all actual
structure constants are integral Laurent polynomials, dominate
their absolute trace targets, and have the trace parity, then
the table has at least $7958$ cancelling pairs.

In particular no table of this type with unequal monic quadratic
cells can improve NN-G021's $7064$ construction, even by raising
the degree of $W$. Combined with NN-G021, the exact minimum is
$7064$ when both cells are allowed to be arbitrary monic quadratics,
$\deg W\le2$, and the ratios in (G022.1) are positive rational
constants. This is an enlargement of that result, not an assertion
about unrestricted dimension rings.

\paragraph{Integral normalization and the parity constraints.}
Every bar-invariant integer Laurent polynomial is in $\mathbb Z[D]$:
write it in the basis $1,q^j+q^{-j}$ and use the recurrence
$C_0=2$, $C_1=D-2$, $C_{j+1}=(D-2)C_j-C_{j-1}$.
Thus an actual table coefficient that is a rational function of
$D$ is an integer polynomial in $D$ as soon as it is an integral
Laurent polynomial. In particular the two $U_i$ below belong
to $\mathbb Z[D]$.

Write $W=gW_0$ with positive integer coefficient content $g$
and primitive $W_0$. The direct mixed-cell coefficient gives
$W\ge D$ and $W\equiv D\pmod2$, so $g$ is odd,
$W_0\in\mathbb Z[D]$, and $W_0\equiv D\pmod2$.
Each $q^2\delta_i$ is monic and primitive. Gauss's content lemma
in $\mathbb Z[q^{\pm1}]$ therefore gives integers
\[
 a_i=gk_i>0,\qquad g\mid a_1a_2.
\]
The actual coefficient $\delta_i$ has the same parity as
$\delta_{\rm tr}=(D-1)(D-4)$, namely $D(D-1)$ modulo two.
Consequently $W_0\delta_i\equiv D^2(D-1)\ne0$ in
$\mathbb F_2[D]$. The first $G$ trace is $4D(D-1)$ and the
second is $D^2(D-1)$. Their parities force
\[
 a_1\text{ even},\qquad a_2\text{ odd},\qquad
 a_1a_2/g\text{ a positive even integer}.             \tag{G022.2}
\]
The closed table identities now read
\[
 \begin{aligned}
 G_i&=a_iW_0\delta_i,&
 S&=(a_1a_2/g)W_0\delta_1\delta_2,\\
 U_i&=\frac{g\delta_i(a_iW_0-D)}{a_j\delta_j},&
 V_i&=\frac{g\delta_i(a_iDW_0-\Lambda_i)}{a_j\delta_j}
       \quad(j\ne i).
 \end{aligned}                                        \tag{G022.3}
\]
In particular $U_i\ge D$ and $U_i\equiv D\pmod2$.

\paragraph{The two quadratics must share an odd integer root.}
If $\delta_1$ and $\delta_2$ were coprime in $\mathbb Q[D]$,
polynomiality of $U_1$ in (G022.3) would imply
$\delta_2\mid a_1W_0-D$. Division by a monic integer polynomial
has integer quotient and remainder, so this divisibility holds
in $\mathbb Z[D]$. Reducing the quotient identity modulo two
would say that the degree-two monic polynomial $D(D-1)$ divides
$D$ in $\mathbb F_2[D]$. This is impossible. Notice that this
argument places no bound on $\deg W$.

Distinct monic quadratics sharing a root over $\mathbb Q$ have
a common linear factor; its root is rational and hence an integer.
We may therefore write
\[
 \delta_i=(D-r)(D-s_i),\qquad r,s_1,s_2\in\mathbb Z,
 \qquad s_1\ne s_2.
 \tag{G022.4}
\]
After cancelling $D-r$ in (G022.3), the two remaining linear
factors are coprime. Polynomiality of $U_1,U_2$ gives
\[
 a_1W_0(s_2)=s_2,\qquad a_2W_0(s_1)=s_1.             \tag{G022.5}
\]
The two roots of each $\delta_i$ have opposite parity, since
its reduction is $D(D-1)$. If $r$ were even, $s_2$ would be odd,
contradicting the first equality of (G022.5), whose left side
is even. Thus $r$ is odd and both $s_i$ are even.

The coefficient at $q$ in (G022.4) is $4-r-s_i$ and the
constant Laurent coefficient is $6-2(r+s_i)+rs_i$.
Their domination conditions give
\[
 r+s_i\le3,\qquad 6-2(r+s_i)+rs_i\ge0.              \tag{G022.6}
\]
Also $\delta_i(4)=(4-r)(4-s_i)>0$. Neither root equals four;
they must both be below four, because being both above four
would contradict their sum bound. Hence $r\le3$ and $s_i\le2$.
If $r=3$, the constant coefficient in (G022.6) equals $s_i$,
while the sum bound gives $s_i\le0$. Thus both $s_i$ would
be zero, contrary to their distinctness. We have proved
\[
 r\le1,\qquad r\text{ odd},\qquad s_i\le2\text{ even}.
 \tag{G022.7}
\]

\paragraph{All exceptional small roots are excluded.}
First suppose $s_2=0$. Then the first equation of (G022.5)
gives $W_0(0)=0$, so write $W_0=DH$ with $H\in\mathbb Z[D]$.
Since $s_1\ne0$, the second equation gives $a_2H(s_1)=1$.
Therefore $a_2=1$, $H(s_1)=1$, and $H-1=(D-s_1)Q$ for an
integer polynomial $Q$. The second $U$ coefficient is then
\[
 U_2=(g/a_1)D^2Q(D).
\]
As $U_2\in\mathbb Z[D]$, it is divisible by $D^2$ in that
ring, regardless of the rational scalar displayed in this
expression. Its parity therefore cannot equal $D$. This
contradicts the required parity of $U_2$, and excludes $s_2=0$.

For any even integer $s$, the congruence $W_0\equiv D\pmod2$
implies that $W_0(s)$ is even. Thus (G022.5) says that the
nonzero even integer $s_2$ is divisible by $2a_1\ge4$.
Together with $s_2\le2$, this proves $s_2\le-4$.

If $s_1=0$, then $W_0=DH$ as before. Since $s_2\ne0$,
the first equation of (G022.5) would give $a_1H(s_2)=1$,
impossible because $a_1$ is a positive even integer.
If instead $s_1=2$, the second equation and oddness of $a_2$
give $a_2=1$ and $W_0(2)=2$. Put $s_2=-a_1t$; here
$t=-W_0(s_2)>0$ is an even integer. An integer polynomial's
value difference $W_0(s_2)-W_0(2)=-t-2$ must be divisible by
$s_2-2=-a_1t-2$. But
\[
 0<t+2<a_1t+2,
\]
which makes that divisibility impossible. These cases exhaust
the even possibilities above $-2$. We have proved
\[
 s_1\le-2,\qquad s_2\le-4,\qquad
 d_1:=\delta_1(4)\ge18,\quad d_2:=\delta_2(4)\ge24.
 \tag{G022.8}
\]

\paragraph{The primitive gluing polynomial has mass at least eight.}
Let $w_0=W_0(4)$, the sum of its Laurent coefficients. The
conditions $W_0\ge0$, $gW_0\ge D$ and
$W_0\equiv D\pmod2$ show that $w_0$ is even and at least four.
If it is four, $W_0=D$; if it is six, $W_0=D+2$.
Indeed every additional noncentral symmetric pair has even
coefficients and contributes at least four to the mass.

The first possibility contradicts
$a_1W_0(s_2)=s_2$ because $a_1\ge2$ and $s_2\ne0$.
In the second, $a_2(s_1+2)=s_1$. The case $a_2=1$ is
impossible. For odd $a_2\ge3$ this requires
$s_1=-2a_2/(a_2-1)\in[-3,-2)$, which contains no even integer.
Thus
\[
 w_0\ge8.                                             \tag{G022.9}
\]
This argument again allows $W$ of arbitrary degree.

\paragraph{The score lower bound.}
The exact nonunit coefficient mass for unequal cells is
\[
 T=S(4)+20(G_1(4)+G_2(4))+200W(4)
 +20\sum_i(U_i(4)+V_i(4))
 +26(d_1+d_2)+14\sum_i(\delta_i\Lambda_i)(4).
\]
The base table targets give $U_i(4)\ge4$, $V_i(4)\ge16$
and $(\delta_1\Lambda_1)(4)\ge8$,
$(\delta_2\Lambda_2)(4)\ge16$, as in (G021.4).
Using (G022.2), (G022.8) and (G022.9), all terms are bounded
below by
\[
 \begin{aligned}
 T&\ge2w_0d_1d_2+20w_0(2d_1+d_2)+200w_0
             +26(d_1+d_2)+1136\\
  &\ge2\cdot8\cdot18\cdot24
       +20\cdot8(2\cdot18+24)+200\cdot8
       +26(18+24)+1136\\
  &=20340.
 \end{aligned}                                        \tag{G022.10}
\]
Subtracting the true absolute nonunit mass $4424$ and dividing
by two gives at least $7958$ cancelling pairs. This lower bound
does not assume $\Lambda_i$ effective, make an optimization-window
assumption, or use a numerical search. It proves the statement.
\hfill$\square$

\paragraph{Geometric interpretation and the reduced remaining gate.}
This is a rigidity consequence of the parity of actual effective
Artin--Tate multiplicity objects and integrality of the closed
kernel-composition table. The equality modulo two is the condition
that dimension and signed trace arise from one representation,
as in NN-G007 and SGA 3, I, 4.7.2. The integer roots here are
roots of auxiliary parameter polynomials; they are not asserted
to be Frobenius eigenvalues or points of an arithmetic scheme.
Archive cote 2, scan 17 supplies the honest kernel-product
dictionary through NN-G010, not an arithmetic interpretation
of those roots.

Within monic quadratic cells and positive rational constant
ratios, any realization competitive with $7064$ must have
$\delta_1=\delta_2$, even for higher-degree $W$. The remaining
optimization possibilities are higher-degree $W$ in that common
quadratic class, nonmonic or nonquadratic cells, nonconstant
ratios, or a different table shape. This improves the exact
gate without introducing a canonical construction or a
nonstandard purity theorem.

\subsection{NN-G023: the exact quadratic-cell optimum with arbitrary gluing degree}
\paragraph{Statement.}
The minimum cancelling-pair score is $7064$ in the following
entire class: the closed rank-three table shape; monic quadratic
$\delta_1,\delta_2\in\mathbb Z[D]$; arbitrary-degree
$W\in\mathbb Z[D]$; positive rational constant ratios
$G_i/(W\delta_i)$; and arbitrary rational functions
$\Lambda_i\in\mathbb Q(D)$. All actual multiplication
coefficients must be effective, dominate their absolute trace
targets, and have the correct parity.

Thus the degree restriction on $W$ in NN-G021 is removed.
The proof does not infer any bound on Laurent support locations
from a bound on coefficient mass. It leaves those locations
arbitrary and eliminates the only three potentially better
families uniformly in the location. Classification of all
minimizers in this enlarged class is not asserted.

\paragraph{Reduction to a finite mass-and-congruence relaxation.}
NN-G022 excludes unequal quadratic cells from any score below
$7064$, so put $\delta_1=\delta_2=\delta$. Normalize
$W=gW_0$ and $G_i=a_iW_0\delta$ exactly as in NN-G021.
Integrality of the auxiliary parameters is now automatic.
Each $V_i$ is an actual integral Laurent polynomial and
is a rational function of $D$, so it is in $\mathbb Z[D]$
by the bar-invariance argument in NN-G022. The identity
$\Lambda_i=k_iDW-k_jV_i$ first puts $\Lambda_i$ in
$\mathbb Q[D]$. Its actual cell coefficient
$\delta\Lambda_i$ is in $\mathbb Z[D]$; division by the
monic polynomial $\delta$ then puts $\Lambda_i$ in
$\mathbb Z[D]$. Thus rational auxiliary factors create
no additional case. NN-G022's unequal-cell bound never
used integrality or polynomiality of $\Lambda_i$.

The content and parity arguments there do not use $\deg W\le2$:
$g$ is odd, $a_1$ is positive even, $a_2$ positive odd,
and $g\mid a_1a_2$. Similarly its mass bound (G021.9)
does not use that degree hypothesis. Thus every proposed
improvement still satisfies
\[
 d:=\delta(4)\le30,\quad w:=W_0(4)\le36,
 \quad a_1+a_2\le52,\quad g\le23.
 \tag{G023.1}
\]
The quadratic is (G021.6), with $d=4+4a+2b$.

Write the primitive gluing polynomial as
\[
 W_0=c_0+c_1(q+q^{-1})+
          \sum_{j\ge2}c_j(q^j+q^{-j}),
\]
where $c_0\ge2$ is even, $c_1\ge1$ is odd and the other
$c_j$ are nonnegative even integers of finite support.
The two conditions $U_i\in D+2\mathbb N[q^{\pm1}]$ imply
\[
 M\mid c_j\quad(j\ge2),\qquad
 M=\operatorname{lcm}\left(
 \frac{2a_2}{\gcd(ga_1,2a_2)},
 \frac{2a_1}{\gcd(ga_2,2a_1)}\right).
 \tag{G023.2}
\]
Indeed the coefficient at such an exponent in $U_i$ is
$ga_ic_j/a_j$, which must be an even nonnegative integer.
The second factor in this least common multiple retains
the full two-primary part of $2a_1$, so $4\mid M$.
Put $h=\sum_{j\ge2}c_j/M$. Then
\[
 w=c_0+2c_1+2Mh.                                     \tag{G023.3}
\]
This bounds the \emph{total coefficient mass} in the outer
positions; it does not bound the indices $j$.

The coefficient at zero and the coefficient at one of each
$U_i$ must also have their exact parity and domination.
Explicitly, for $(c,t)=(c_0,2)$ and $(c_1,1)$ the two numbers
\[
 g(a_ic-t)/a_j\quad(j\ne i)
 \tag{G023.4}
\]
must be integers at least $t$ and congruent to $t$ modulo two.
Primitivity requires $\gcd(c_0,c_1)=1$ when $h=0$, and
requires $\gcd(c_0,c_1,M)=1$ when $h>0$. The latter is
used only as a necessary condition.

For an actual table set $\ell_i=\Lambda_i(4)$. Its two
effective cell products give $d\ell_1\ge8$ and
$d\ell_2\ge16$. Evaluating the $V_i$ congruences and
domination gives, with $a_j'=a_j/\gcd(g,a_j)$,
\[
 \begin{gathered}
 \ell_i\equiv4a_iw\pmod{2a_j'},\\
 \left\lceil X_i(4)/d\right\rceil\le\ell_i\le
 \left\lfloor\frac{4ga_iw-16a_j}{g}\right\rfloor,
 \qquad (X_1(4),X_2(4))=(8,16).
 \end{gathered}                                       \tag{G023.5}
\]
No degree bound on $\Lambda_i$ is used here.
The exact mass expression is (G021.8). For fixed other
parameters it is affine in $\ell_i$, with coefficient
$14d-20g/a_j$. Minimize it over the one-dimensional
congruence interval (G023.5) by taking the first allowed
integer if that coefficient is nonnegative, and the last
if it is negative. This gives a necessary lower bound on
the score, not a claim that the chosen masses arise from
actual Laurent polynomials.

All the remaining ranges are finite: positive even
$a_1\le50$, positive odd $a_2\le51$ with sum at most $52$;
odd $g\le23$ dividing $a_1a_2$; $a,b\ge0$ with
$4+4a+2b\le30$; even $2\le c_0\le36$; odd
$1\le c_1\le17$; and $0\le h\le(36-c_0-2c_1)/(2M)$.
In addition impose $G_i(4)\ge48$ and $S(4)\ge576$,
which are necessary actual coefficient-mass inequalities.
These are precisely the loops and tests in
\texttt{scalar\_dimension\_bound.py}; it uses only integer
arithmetic and records the whole relaxation.

The complete calculation gives $42$ rows whose relaxed
score is strictly below $7064$. Thirty-nine have $h=0$,
so $\deg W\le1$ and NN-G021 already excludes an actual
score below $7064$. Exactly three rows have $h>0$:
\[
 \begin{array}{c|c|c|c|c|c}
 (a_1,a_2,g)&(a,b)&(c_0,c_1,M,h)&w&d&
       \text{relaxed score}\\\hline
 (2,1,1)&(0,0)&(6,3,4,1)&20&4&5180\\
 (2,1,1)&(0,0)&(10,3,4,1)&24&4&6624\\
 (2,1,1)&(0,1)&(6,3,4,1)&20&6&6916
 \end{array}                                         \tag{G023.6}
\]
Because $h=1$, every such row has
\[
 W=c_0+3(q+q^{-1})+4(q^j+q^{-j})
 \quad\text{for one arbitrary integer }j\ge2.
 \tag{G023.7}
\]
The complete $42$-row JSON certificate is included. The
three remaining families, rather than a bounded list of
support indices, are now handled by exact algebra.

\paragraph{Excluding the two mass-four cell families.}
For $d=4$ the quadratic is
$\delta=q^2+q+q^{-1}+q^{-2}$. The nonnegative-index
Laurent coefficient list of $\delta^2$ is
$(4,2,1,2,1)$, and its coefficients at indices $j\ge2$
are at most two. Since $S=2W\delta^2$, (G023.7) gives
\[
 [q^0]S=8c_0+24+16[q^j]\delta^2\le8c_0+56.
\]
For $c_0=6$ or $10$ this is at most $104$ or $136$.
But the true bottom coefficient
$S_{\rm tr}=4D^2(D-1)^2$ has constant Laurent coefficient
$144$. These two families violate coefficientwise
domination for every integer $j\ge2$.

\paragraph{The remaining mass-six family and the fourth-root test.}
Here $\delta=(D-1)(D-2)$ and
$W=3D+4(q^j+q^{-j})$. The parity condition on
$V_2=(DW-\Lambda_2)/2$ gives
\[
 \Lambda_2\equiv DW-2K_2
  \equiv -D^2+10D-8=:C\pmod4.
 \tag{G023.8}
\]
Thus $\ell_2\equiv C(4)=16\equiv0\pmod4$.
Its effective product $A_2=\delta\Lambda_2$ has mass
$6\ell_2\ge16$, so the only possibility below eight
would be $\ell_2=4$.

We exclude that value without referring to $j$.
The coefficient list of $\delta C$ is
$(22,23,10,5,-1)$. Combining this congruence modulo four
with $A_2\ge X_2=(2,3,2,1,1)$ shows that
\[
 A_2=B+4H,\qquad
 B=(2,3,2,1,3),\qquad
 H\in\mathbb N[q^{\pm1}]\text{ symmetric}.
 \tag{G023.9}
\]
The mass of $B$ is $20$. If $\ell_2=4$, the mass of
$A_2$ is $24$, so $H$ has mass one. A symmetric effective
integer Laurent polynomial of mass one is the constant
one. Hence $A_2$ would have coefficient list
$(6,3,2,1,3)$.

At $q=i$ this last polynomial has value $6-4+6=8$.
On the other hand $\delta(i)=0$, so the identity
$A_2=\delta\Lambda_2$ requires $A_2(i)=0$.
This contradiction proves $\ell_2\ge8$.
Also $\ell_1$ is even and $6\ell_1\ge8$, so
$\ell_1\ge2$. Substituting the fixed masses of this
family into the exact score expression gives
\[
 T=17832+64\ell_1+74\ell_2
       \ge17832+64\cdot2+74\cdot8=18552.
\]
Its score is therefore at least $7064$, rather than
the unrealizable scalar-relaxation value $6916$.

Every potential improvement in (G023.6) has now been
excluded for arbitrary outer location. NN-G021 supplies
an actual table of score $7064$, and NN-G022 excludes
unequal quadratics. This proves the exact optimum in
the full class stated above.\hfill$\square$

\paragraph{Geometric dictionary and remaining global scope.}
The fourth-root evaluation is the trace of a finite
character of the grading: restrict a $\mathbb G_m$
representation to $\mu_4$ and evaluate its character at
$i$. This is the diagonalizable-group/comodule formalism
of SGA 3, I, 4.7.2--4.7.3. It is not a finite-field
cardinality, a Frobenius eigenvalue, or a zero of an
arithmetic zeta function. Effective graded objects may
have zero trace at this character. What fails in the
hypothetical small-score table is the required divisibility
of one effective character by the specified cell character.
The actual effective multiplicity objects and their
kernel products remain those of NN-G010 and archive
cote 2, scan 17.

The arbitrary-degree gluing gate is closed for monic
quadratic cells and rational constant ratios. A smaller
score must use nonmonic or nonquadratic cells, nonconstant
ratios, or another
table shape. None of these possibilities is dismissed.
The optimum remains noncanonical and has no implication
for a rank-four canonical basis, PH0, PH1 or the missing
nonstandard-purity bounds.

\subsection{NN-G024: nonprimitive common cells give the exact 5932-pair optimum}
\paragraph{Statement.}
There is an effective associative rank-three dimension ring with
$5932$ cancelling pairs, of which $5208$ are in the bottom cell.
Moreover $5932$ is the exact minimum in the entire following
class: the closed rank-three table shape, a common cell polynomial
$\delta\in\mathbb Z[D]$ of arbitrary degree and coefficient content,
arbitrary-degree $W\in\mathbb Z[D]$, positive rational constant
ratios $G_i/(W\delta)$, and arbitrary rational auxiliary functions
$\Lambda_i\in\mathbb Q(D)$. Only the actual multiplication
coefficients are required to be integral, effective, coefficientwise
dominant over the absolute trace targets, and congruent to those
targets modulo two.

This class includes all common-cell constructions in NN-G019,
NN-G021 and NN-G023 and removes their monicity and cell-degree
restrictions. The new optimum is attained by a cell polynomial
of content three and a nonintegral auxiliary $\Lambda_2$.
The theorem does not cover unequal general cells, nonconstant
ratios, or other table shapes.

\paragraph{An explicit integral table with a rational auxiliary factor.}
Put $\delta_0=(D-1)(D-2)$ and choose
\[
 \begin{aligned}
 \delta&=3\delta_0,& W&=3D,\\
 G_1&=6D\delta_0,&G_2&=9D\delta_0,&S&=18D\delta_0^2,\\
 U_1&=D,&U_2&=3D,\\
 \Lambda_1&=D-2,&\Lambda_2&=(D^2-2D-4)/3,\\
 V_1&=2D^2-D+2,&V_2&=4D^2+D+2.
 \end{aligned}                                                   \tag{G024.1}
\]
The ratios are $(k_1,k_2)=(2/3,1)$. The identities
\[
 S=G_1G_2/W,\qquad
 U_i=(k_iW-D)/k_j,\qquad
 V_i=(k_iDW-\Lambda_i)/k_j\quad(j\ne i)
\]
hold exactly. The two matrix blocks and two scalar blocks in
NN-G011 therefore give an injective associative realization
of the closed table over $\mathbb Q(D)$: neither
$\Lambda_i(D^2-\Lambda_i)$ vanishes identically.

For effectivity and integrality it is the actual product
$\delta\Lambda_i$ that matters. In symmetric Laurent
coefficient lists indexed by nonnegative exponents, the base
coefficients of this table and their absolute trace targets are
\[
\begin{array}{c|l|l}
 &\text{dimension coefficient}&\text{absolute trace target}\\\hline
 \delta &(6,3,3)&(0,1,1)\\
 \delta\Lambda_1&(6,9,3,3)&(2,1,1,1)\\
 \delta\Lambda_2&(2,5,2,3,1)&(2,3,2,1,1)\\
 W&(6,3)&(2,1)\\
 U_1&(2,1)&(2,1)\\
 U_2&(6,3)&(2,1)\\
 V_1&(12,7,2)&(6,5)\\
 V_2&(28,17,4)&(6,3,2)\\
 G_1&(36,30,18,6)&(16,12,4)\\
 G_2&(54,45,27,9)&(14,11,5,1)\\
 S&(504,450,324,180,72,18)&(144,120,68,24,4)
\end{array}                                                       \tag{G024.2}
\]
Every left list is nonnegative, coefficientwise at least the
right list, and congruent to it modulo two. The closed table
uses these coefficients multiplied by nonnegative powers of
$D$. For signed trace cell coefficients use
$|D^\alpha f|\le D^\alpha|f|$ coefficientwise. These facts
prove every dimension-ring condition. In particular the
nonintegral $\Lambda_2$ is not itself a structure constant;
its product with the nonprimitive cell is integral.

At $D=4$ the bottom mass is
\[
 2592+20(144+216)+200\cdot12+20(4+12+30+70)=14512.
\]
The two cell masses sum to
$52\cdot18+14\cdot6(6+4)=1776$, so $T=16288$.
Subtracting the true absolute masses gives
\[
 (16288-4424)/2=5932,\qquad (14512-4096)/2=5208.
 \tag{G024.3}
\]
The independent integer-Laurent verifier
\texttt{verify\_dimension\_5932.py} reconstructs all $132$
positions and checks all $1000$ associativity triples. It
handles the actual integer products directly and does not
round or formally clear an unverified denominator.

\paragraph{Content normalization for the entire common-cell class.}
Write
\[
 \delta=c\delta_0,\quad W=gW_0,
\]
where $c,g>0$ are coefficient contents and $\delta_0,W_0$
are primitive symmetric integral Laurent polynomials.
They are also in $\mathbb Z[D]$, because symmetric integral
Laurent polynomials have an integral triangular expansion
in $q^j+q^{-j}$ and hence in $D=q+2+q^{-1}$.
The actual parity conditions imply that $c,g$ are odd,
\[
 \delta_0\equiv D(D-1),\qquad W_0\equiv D\pmod2.
\]
Both polynomials are effective. In particular
$d:=\delta_0(4)\ge4$ and $w:=W_0(4)\ge4$ are even.

Gauss's content lemma applied after Laurent shifts gives
positive integers $a_i$ such that
\[
 \begin{aligned}
 G_i&=a_iW_0\delta_0,&
 S&=(a_1a_2/g)W_0\delta_0^2,\\
 U_i&=g(a_iW_0-cD)/a_j,&
 V_i&=g(a_iDW_0-\Psi_i)/a_j,
 \end{aligned}                                                     \tag{G024.4}
\]
where $\Psi_i=c\Lambda_i$ and $j\ne i$.
Here $a_i=k_i g c$, $a_1$ is even, $a_2$ is odd, and
$g\mid a_1a_2$. Indeed the primitive products force the
integrality statements, and reduction modulo two of the
$G_i$ gives their different scalar parities.

Even though $\Lambda_i$ may be rational, $\Psi_i$ is in
$\mathbb Z[D]$. First (G024.4), with integral actual $V_i$,
puts $\Psi_i$ in $\mathbb Q[D]$. The actual within-cell
coefficient is $\delta\Lambda_i=\delta_0\Psi_i\in\mathbb Z[D]$.
Primitivity of $\delta_0$ and Gauss's lemma force
$\Psi_i\in\mathbb Z[D]$. This is the exact integrality
statement; it does not force division by the content $c$.
Modulo two these actual cell products also give
$\Psi_1\equiv D$ and $\Psi_2\equiv D^2$.

Set $p_i=\Psi_i(4)$. The exact mass expression is
\[
 \begin{split}
 T={}&(a_1a_2/g)wd^2+20wd(a_1+a_2)+200gw\\
 &+20\sum_i\{U_i(4)+V_i(4)\}+52cd+14d(p_1+p_2).
 \end{split}                                                       \tag{G024.5}
\]
The two effective cell products imply $dp_1\ge8$ and
$dp_2\ge16$. Also $U_i(4)\ge4$ and $V_i(4)\ge16$.
Thus the score bound of NN-G021 applies with an extra
nonnegative $52(c-1)d$ term. For any proposed improvement
on (G024.3) it gives the following generous complete bounds:
\[
 d\le30,\quad w\le36,\quad a_1+a_2\le52,\quad g\le23.
 \tag{G024.6}
\]
Furthermore the right side of (G024.5) is at least
$3024+208c$: use $a_1a_2/g\ge2$, $a_1+a_2\ge3$,
$d,w\ge4$, and the $1136$ lower bound for the last
cell-product term and the four $U,V$ terms. Since a
proposed improvement has $T<16288<18552$, the weaker
bound $c\le73$, with $c$ odd, is valid. None of these
bounds asserts that a Laurent support index is bounded.

\paragraph{A complete finite relaxation, retaining arbitrary support locations.}
Write
\[
 W_0=c_0+c_1(q+q^{-1})+\sum_{j\ge2}c_j(q^j+q^{-j}).
\]
Here $c_0\ge2$ is even, $c_1\ge1$ is odd, and $c_j$
is nonnegative even for $j\ge2$. The same outer-coefficient
argument as in (G023.2) gives $M\mid c_j$, where
\[
 M=\operatorname{lcm}\left(
 \frac{2a_2}{\gcd(ga_1,2a_2)},
 \frac{2a_1}{\gcd(ga_2,2a_1)}\right),\quad 4\mid M.
\]
Put $h=\sum_{j\ge2}c_j/M$, so $w=c_0+2c_1+2Mh$.
For $(b,t)=(c_0,2),(c_1,1)$ the necessary zero and
first coefficient conditions on $U_i$ are
\[
 g(a_i b-ct)/a_j\in t+2\mathbb N.                     \tag{G024.7}
\]
Primitivity requires $\gcd(c_0,c_1)=1$ for $h=0$ and
$\gcd(c_0,c_1,M)=1$ for $h>0$.

Let $a_j'=a_j/\gcd(g,a_j)$. Evaluating the $V_i$
congruences and domination gives the integer intervals
\[
 \begin{gathered}
 p_i\equiv4a_iw\pmod{2a_j'},\\
 \left\lceil X_i(4)/d\right\rceil\le p_i\le
 \left\lfloor(4ga_iw-16a_j)/g\right\rfloor,
 \quad(X_1(4),X_2(4))=(8,16).
 \end{gathered}                                                       \tag{G024.8}
\]
For fixed other data the coefficient of $p_i$ in (G024.5)
is $14d-20g/a_j$. Its minimum over (G024.8) is therefore
at the first allowed integer when this coefficient is
nonnegative and at the last otherwise. This is a lower
bound using necessary conditions, not an assertion that
such a polynomial $\Psi_i$ exists.

There is one additional necessary leading-coefficient
condition when $h=0$. Then $\deg W_0=1$. If $\deg\Psi_2>2$,
its highest Laurent coefficient would have opposite signs
in the effective polynomials $\delta_0\Psi_2$ and $V_2$.
The leading coefficient of $\delta_0$ is positive, so this
is impossible. Since $\Psi_2\equiv D^2\pmod2$, its
coefficient $L$ at $q^2$ must therefore be positive odd.
The coefficient of $V_2$ at $q^2$ is required to be an
even integer at least two. Consequently
\[
 L\ge1,\quad L\equiv a_2c_1-2a_1'\pmod{2a_1'},\quad
 L\le\left\lfloor(g a_2c_1-2a_1)/g\right\rfloor.
 \tag{G024.9}
\]
This test uses no cell-degree restriction.

For reproducibility the whole finite range is as follows:
even $2\le a_1\le50$, odd $1\le a_2\le51$ with sum
at most $52$; odd $1\le g\le23$ dividing $a_1a_2$;
odd $1\le c\le73$; even $4\le d\le30$;
even $2\le c_0\le36$; odd $1\le c_1\le17$;
and $0\le h\le(36-c_0-2c_1)/(2M)$.
Impose (G024.7), the primitive-content necessary tests,
$G_i(4)\ge48$, $S(4)\ge576$, and the endpoint lower
bound from (G024.8). The included integer-only script
\texttt{verify\_common\_cell\_content\_optimum.py}
records every survivor below $T=16288$.
There are seventeen before (G024.9), and exactly five
after it:
\[
\begin{array}{c|c|c|c|c|r}
 (a_1,a_2,g)&c&d&(c_0,c_1,M,h)&w&\text{relaxed score}\\\hline
 (6,9,3)&3&4&(2,1,4,0)&4&3716\\
 (2,1,1)&1&4&(6,3,4,1)&20&5180\\
 (2,1,1)&1&4&(6,7,4,0)&20&5180\\
 (10,5,5)&3&4&(2,1,4,0)&4&5268\\
 (6,15,3)&1&4&(2,1,20,0)&4&5372
\end{array}                                                         \tag{G024.10}
\]
No outer index has been enumerated or bounded.

\paragraph{Uniform exclusion of every survivor.}
All five rows have $d=4$. The effective parity-correct
primitive polynomial of that mass is unique:
\[
 \delta_0=q^2+q+q^{-1}+q^{-2}=D(D-3).
\]
Indeed these four odd coefficients already use the entire
mass. In the first, fourth and fifth rows $W_0=D$ and
$[q^0](D\delta_0)=2$. The first and fifth rows therefore
have $[q^0]G_1=12<16$, while the fourth has
$[q^0]G_2=10<14$. Each violates an actual trace target.

For the third row the coefficient list of $\delta_0^2$
is $(4,2,1,2,1)$ and $S=2W_0\delta_0^2$, so
$[q^0]S=2(6\cdot4+2\cdot7\cdot2)=104<144$.
For the second row $h=1$ means
\[
 W_0=6+3(q+q^{-1})+4(q^j+q^{-j})\quad(j\ge2)
\]
for an arbitrary integer $j$. The same list yields
\[
 [q^0]S=48+24+16[q^j]\delta_0^2\le104<144.
\]
Thus this entire unbounded-location family also fails.
All necessary-relaxation survivors below $5932$ have
been excluded. The actual table (G024.1) attains $5932$,
proving the asserted optimum.\hfill$\square$

\paragraph{Geometric translation and the remaining gate.}
NN-G010 realizes the integral effective multiplicities of
(G024.2) by actual graded sign representations in an
additive monoidal kernel category on a ten-point scheme.
The exact dictionaries remain SGA 3, I, 4.7.2--4.7.3
for graded representations and archive cote 2, scan 17
for finite-kernel composition. There is no object of
fractional dimension attached to $\Lambda_2$; the rational
factor is merely a coordinate factorization in $\mathbb Q(D)$.
The structure constant $\delta\Lambda_2$ is the actual
integral object. The content-three factor is three copies
of its effective cell object, not division by three in
an effective Grothendieck semiring.

For the equal-cell, constant-ratio tables realized by these
chosen kernel generators, the minimum sum over the $132$
coefficient positions and grading pieces of half the excess of dimension over
absolute signed trace is exactly $5932$. This is a minimum
for the specified character-valued multiplication data and
its NN-G010 realizations, not for all monoidal categories
or all generating sets with the same Grothendieck ring.

The same supplied trace algebra therefore admits smaller
composing effective lifts after a purely auxiliary
integrality restriction is removed. This demonstrates
why formal factor parameters cannot be identified with
canonical count objects. It establishes neither a
canonical family nor a rank-four basis, PH0, PH1 or
nonstandard purity. The precise remaining lower-score
gates are unequal general cells, nonconstant ratios,
or a different multiplication-table shape; arbitrary
common-cell degree and content are now covered.

\subsection{NN-G025: the exact 4520-pair optimum for proportional cells}
\paragraph{Statement.}
The exact minimum is $4520$ cancelling pairs in the closed
rank-three table class with arbitrary proportional effective
cell polynomials $\delta_1,\delta_2\in\mathbb Z[D]$,
arbitrary-degree $W\in\mathbb Z[D]$, positive rational
constant ratios $G_i/(W\delta_i)$, and arbitrary rational
auxiliary functions $\Lambda_i$. All actual coefficients
must obey the usual integrality, domination, effectivity
and parity conditions. Proportional means
$\delta_1/\delta_2\in\mathbb Q_{>0}$; the contents need
not agree and the common primitive polynomial has no
degree restriction.

An attainer has $3532$ bottom-cell cancelling pairs.
This enlarges NN-G024's equal-cell class, whose exact
minimum $5932$ remains valid. The new global upper
bound is $4520$; global minimality without the
proportional-cell and constant-ratio restrictions is
not asserted.

\paragraph{The attainer and its exact score.}
Let $f=(D-1)(D-2)$. Set
\[
 \begin{aligned}
 \delta_1&=f,&\delta_2&=3f,&W&=3D,\\
 G_1&=6Df,&G_2&=5Df,&S&=10Df^2,\\
 U_1&=3D,&U_2&=D,\\
 \Lambda_1&=6D^2-25D+20,&
 \Lambda_2&=(D^2-2D-4)/3,\\
 V_1&=15D-12,&V_2&=2D^2+D+2.
 \end{aligned}                                                     \tag{G025.1}
\]
Here $G_i=k_iW\delta_i$ with $(k_1,k_2)=(2,5/9)$.
Direct substitution gives
\[
 \frac{G_1G_2}{S}=W,\qquad
 \frac{G_i(G_i-D\delta_i)}{S}=U_i,\qquad
 \frac{G_i(DG_i-\delta_i\Lambda_i)}{S}=V_i.
\]
The injective two-matrix-block/two-scalar-block construction
of NN-G011 works with the two different $\delta_i$ and
these nonzero $\Lambda_i(D^2-\Lambda_i)$, proving
associativity.

The changed effective base coefficients, in symmetric
nonnegative-index Laurent lists, are
\[
\begin{array}{c|l|l}
 &\text{dimension coefficient}&\text{absolute trace target}\\\hline
 \delta_1&(2,1,1)&(0,1,1)\\
 \delta_2&(6,3,3)&(0,1,1)\\
 \delta_1\Lambda_1&(22,9,17,5,6)&(2,1,1,1)\\
 \delta_2\Lambda_2&(2,5,2,3,1)&(2,3,2,1,1)\\
 G_1&(36,30,18,6)&(16,12,4)\\
 G_2&(30,25,15,5)&(14,11,5,1)\\
 S&(280,250,180,100,40,10)&(144,120,68,24,4)\\
 V_1&(18,15)&(6,5)\\
 V_2&(16,9,2)&(6,3,2)
\end{array}                                                        \tag{G025.2}
\]
Also $W=U_1=3D$ and $U_2=D$. These lists prove effectivity,
domination and parity for all base coefficients; powers
of $D$ propagate them to the full table exactly as in
NN-G024. No actual coefficient has a denominator.
The independent integer-Laurent checker verifies all
$132$ positions and all $1000$ associativity triples.

At $D=4$ the bottom mass is
\[
 1440+20(144+120)+200\cdot12+20(12+4+48+38)=11160.
\]
The cell mass is
$26(6+18)+14(6\cdot16+6\cdot4)=2304$.
Thus $T=13464$ and
\[
 (13464-4424)/2=4520,\qquad(11160-4096)/2=3532.
 \tag{G025.3}
\]

\paragraph{Normalization and complete necessary bounds.}
Write $\delta_i=c_i f$ with $f$ primitive effective and
$c_i$ positive odd integers, and $W=gW_0$ with $g$
positive odd and $W_0$ primitive effective. Put
$d=f(4)$ and $w=W_0(4)$. The same Gauss and parity
arguments as NN-G024 give positive integers $a_1$ even
and $a_2$ odd, with $g\mid a_1a_2$, and
\[
 \begin{aligned}
 G_i&=a_iW_0f,&S&=(a_1a_2/g)W_0f^2,\\
 U_i&=g(a_iW_0-c_iD)/a_j,&
 V_i&=g(a_iDW_0-\Psi_i)/a_j,
 \end{aligned}                                                      \tag{G025.4}
\]
where $\Psi_i=c_i\Lambda_i\in\mathbb Z[D]$.
Integrality of $\Psi_i$ follows from polynomiality of
$V_i$ and the primitive product $f\Psi_i$, exactly
as before. It does not require $c_1=c_2$.
The actual cell-product masses are $dp_i$ with
$p_i=\Psi_i(4)$, and $dp_1\ge8$, $dp_2\ge16$.
The exact total mass is
\[
 \begin{split}
 T={}&(a_1a_2/g)wd^2+20wd(a_1+a_2)+200gw\\
 &+20\sum_i\{U_i(4)+V_i(4)\}
       +26(c_1+c_2)d+14d(p_1+p_2).
 \end{split}                                                        \tag{G025.5}
\]
For a proposed improvement on $4520$, the complete
bounds (G024.6) still hold, since $c_1+c_2\ge2$.
Moreover
\[
 T\ge3024+104(c_1+c_2).
\]
Using even the weaker cap $T<16288$ gives
$c_1+c_2\le126$, so each positive odd $c_i\le125$.
This generous bound suffices; no degree bound on $f$
or on outer Laurent support locations is being used.

Use $b_0,b_1$ for the zero and first coefficients of
$W_0$, and $M,h$ as in (G024.7), so
$w=b_0+2b_1+2Mh$.
The complete scalar ranges are those of NN-G024,
with the one content variable replaced by
$1\le c_i\le125$, both odd. For each $i$ impose
\[
 g(a_i b-c_i t)/a_j\in t+2\mathbb N
 \quad\text{for }(b,t)=(b_0,2),(b_1,1).               \tag{G025.6}
\]
The bound from the first coefficient also permits
$c_i\le a_i b_1$ without removing any actual case.
The outer-coefficient divisibility, primitive-content
necessary tests, $S(4)\ge576$, and $G_i(4)\ge48$ are
unchanged. The intervals and endpoint rule (G024.8)
for $p_i$ are unchanged as well. When $h=0$, the
leading-coefficient test (G024.9) is still necessary.

For fixed other parameters, the coefficient of $c_i$
in (G025.5) is $26d-80g/a_j$. Thus minimize the
content contribution at the least admissible $c_i$
when that coefficient is nonnegative, and at the
greatest otherwise. The two contents can be minimized
independently for a necessary relaxation: ignoring
their joint upper bound only enlarges the possible
set. Their individual bounds are imposed. An endpoint
row is a lower bound for its whole class of content
choices, not a claim that all those choices coincide.

The exact integer-only script
\texttt{verify\_proportional\_cells\_optimum.py}
uses precisely these ranges, congruences and endpoint
rules. It leaves just four rows with relaxed score
below $4520$:
\[
\begin{array}{c|c|c|c|c|r}
(a_1,a_2,g)&(c_1,c_2)_{\rm endpoint}&d&(b_0,b_1,h)&
 (p_1,p_2)_{\rm endpoint}&\text{relaxed score}\\\hline
(6,5,3)&(1,3)&4&(2,1,0)&(6,4)&2636\\
(4,9,3)&(1,5)&4&(2,1,0)&(4,8)&3480\\
(6,9,3)&(3,3)&4&(2,1,0)&(6,4)&3716\\
(6,5,3)&(1,3)&6&(2,1,0)&(6,4)&4160
\end{array}                                                        \tag{G025.7}
\]
All have $W_0=D$. In each of the first three,
$f$ has mass four and is consequently $D(D-3)$.
Its product with $D$ has central coefficient two.
Thus their $G_1$ central coefficients are respectively
$12,8,12$, all below the trace target $16$.
These exclusions are independent of the choice of
contents within the row.

\paragraph{The last row and an exact residue obstruction.}
In the remaining row the effective parity-correct
primitive polynomial of mass six is uniquely
$f=(D-1)(D-2)$: relative to the mass-four polynomial,
the only possible addition of mass two is a central
coefficient two. Condition (G025.6) at the first
coefficient forces $c_1=1$ and $c_2=3$. Indeed
$3(6-c_1)/5$ must be a positive odd integer, giving
$c_1=1$, and $(5-c_2)/2$ must be a positive odd
integer, giving $c_2=3$.

Since $W_0=D$, effectivity of $f\Psi_i$ and $V_i$
forces $\deg\Psi_i\le2$ by their opposite leading
signs above degree two. The congruence for $V_1$ is
\[
 \Psi_1\equiv6D^2-5(5D-4)\pmod{10}.
\]
Writing $\Psi_1=L_0+L_1(q+q^{-1})+L_2(q^2+q^{-2})$
gives
\[
 L_0\equiv6,\quad L_1\equiv9,\quad L_2\equiv6
       \pmod{10}.                                      \tag{G025.8}
\]
The highest product coefficient forces $L_2\ge0$,
and $V_1=(3/5)(6D^2-\Psi_1)$ forces $L_2\le6$.
Hence $L_2=6$. The mass $p_1=L_0+2L_1+12$ is
congruent to six modulo ten; $6p_1\ge8$ makes
six its smallest possible positive value.

Suppose $p_1=6$. Then $L_0=-6-2L_1$.
The central coefficient of $f\Psi_1$ is
$2L_0+2L_1+12=-2L_1$ and must be at least two,
so $L_1\le-1$. Its coefficient at $q^3$ is
$L_1+6$ and must be at least one, so $L_1\ge-5$.
Together with (G025.8) this forces $L_1=-1$ and
$L_0=-4$. But its coefficient at $q$ would then be
$L_0+3L_1+6=-1$, below the target one.
Thus $p_1\ne6$, and consequently $p_1\ge16$.
The second mass is a multiple of four and satisfies
$6p_2\ge16$, so $p_2\ge4$.

Substitution in the exact mass identity for this row gives
\[
 T=12016+72p_1+74p_2
   \ge12016+72\cdot16+74\cdot4=13464.
\]
It cannot improve on $4520$. All four rows of the
complete necessary relaxation have now been excluded
below that score, and (G025.1) attains it. This proves
the exact minimum in the stated class.\hfill$\square$

\paragraph{Geometric meaning and residual scope.}
The two cells have the same primitive grading character
and different positive multiplicities. Their actual
coefficient objects still lie in the effective graded
sign-representation semiring, and NN-G010 supplies the
same finite-kernel monoidal realization. The precise
classical dictionaries are SGA 3, I, 4.7.2--4.7.3 and
archive cote 2, scan 17. Rational $\Lambda_2$ remains
an auxiliary factor, not a fractional effective object.
Different contents of the two cell characters do not
alter the supplied trace algebra whose effective lift
is being sought; they change the chosen dimension data.

Among the NN-G010 kernel realizations whose two cell characters
are positive integral multiples of the same primitive character
and whose specified ratios are constant, the least total
half-excess of dimension over absolute signed trace is $4520$,
with $3532$ in the bottom-cell coefficients for (G025.1).
These numerical invariants refer to the chosen table and
generators; they are not invariants of every categorification
of the abstract trace algebra.

The exact common-cell minimum and the exact proportional-cell
minimum differ, so equality of cell characters is a
substantive optimization restriction, not a consequence
of the trace constraints. A smaller score must now use
nonproportional cell polynomials, nonconstant ratios,
or another table shape. No canonicality, generic
rank-four positivity, PH0/PH1 or nonstandard purity
conclusion is inferred from this noncanonical optimum.

\subsection{NN-G026: the exact optimum for all constant-ratio closed tables}
\paragraph{Statement.}
The exact cancelling-pair minimum is $4520$ for the closed
rank-three table shape with arbitrary effective cell polynomials
$\delta_1,\delta_2\in\mathbb Z[D]$, arbitrary-degree
$W\in\mathbb Z[D]$, positive rational constant ratios
$G_i/(W\delta_i)$, and arbitrary rational auxiliary functions
$\Lambda_i\in\mathbb Q(D)$. The cells need not be proportional;
there are no degree, monicity, or content restrictions.
Every actual structure constant must still be integral,
effective, coefficientwise dominant over its absolute trace
target, and parity-correct.

The attainer is NN-G025. The new assertion is the exclusion
of every smaller table in this enlarged class. Only the
constant-ratio and closed-shape assumptions remain as
optimization restrictions; an unrestricted dimension-ring
optimum is not asserted.

\paragraph{Normalization without proportionality.}
Write $\delta_i=c_i f_i$, where $f_i\in\mathbb Z[D]$
is primitive effective, and $W=gW_0$ with $W_0$ primitive
effective. All contents $c_i,g$ are positive odd integers.
The reduction of each $f_i$ modulo two is
$\sigma=D(D-1)$, and $W_0\equiv D$.
Gauss's content lemma gives positive integers $a_1$ even,
$a_2$ odd, with $g\mid a_1a_2$, such that
\[
 \begin{aligned}
 G_i&=a_iW_0f_i,&S&=(a_1a_2/g)W_0f_1f_2,\\
 U_i&=\frac{g f_i(a_iW_0-c_iD)}{a_j f_j},&
 V_i&=\frac{g(a_iDW_0f_i-B_i)}{a_j f_j},
 \qquad B_i=\delta_i\Lambda_i.
 \end{aligned}                                                       \tag{G026.1}
\]
Here $B_i$ is the actual integral effective cell coefficient.
Unlike in the proportional case, no integrality of
$c_i\Lambda_i$ is asserted or needed.
Put $d_i=f_i(4)$, $w=W_0(4)$ and $b_i=B_i(4)$.
The masses $d_i,w$ are even and at least four, while
$b_1\ge8$ and $b_2\ge16$.

The exact nonunit mass is
\[
 \begin{split}
 T={}&(a_1a_2/g)wd_1d_2+20w(a_1d_1+a_2d_2)+200gw\\
 &+20\sum_i\{U_i(4)+V_i(4)\}
       +26(c_1d_1+c_2d_2)+14(b_1+b_2).
 \end{split}                                                        \tag{G026.2}
\]
Any improvement would have $T<13464$. The base bounds
$U_i(4)\ge4$, $V_i(4)\ge16$, and $b_1+b_2\ge24$
give the following complete finite ranges:
\[
 \begin{gathered}
 4\le w\le24,\quad4\le d_1\le50,\quad4\le d_2\le78
       \quad\text{all even},\\
 a_1\ge2\text{ even},\quad a_2\ge1\text{ odd},\quad
 a_1+a_2\le33,\quad1\le g\le15\text{ odd},\quad g\mid a_1a_2,\\
 1\le c_i\le125\quad\text{odd}.
 \end{gathered}                                                       \tag{G026.3}
\]
For clarity, these are score bounds, not degree bounds.
Indeed $T\ge472w+1344$ and
$T\ge320(a_1+a_2)+2272$ give the first and scalar-sum
bounds, while $T\ge800g$ gives the content bound for $W$.
Also
\[
 T\ge8d_1d_2+186d_1+106d_2+1936.
\]
Setting $d_2\ge4$ gives $T\ge218d_1+2360$;
setting $d_1\ge4$ gives $T\ge138d_2+2680$.
Finally $T\ge3024+104(c_1+c_2)$ gives an even stronger
bound than the generous individual cap $125$ used here.

\paragraph{An exact scalar relaxation for arbitrary cells and support indices.}
For fixed $a_i,g,w,d_i$, enumerate positive odd $c_i$
within (G026.3) and the necessary bound $c_i\le a_iw/4$.
Impose that
\[
 u_i:=\frac{g d_i(a_iw-4c_i)}{a_jd_j}
              \text{ is an even integer at least four}.             \tag{G026.4}
\]
The coefficient of $c_i$ in (G026.2) is
$d_i(26-80g/(a_jd_j))$. Choose the least admissible content
if that coefficient is nonnegative, and the greatest
otherwise. This gives a lower bound for all contents in
that row. Ignoring their joint bound only weakens the
relaxation, and does not exclude any actual table.

For $b_i$ the necessary condition that $V_i(4)$ be an
even integer at least sixteen is exactly
\[
 \begin{gathered}
 b_i\equiv4a_iwd_i\pmod{M_i},\qquad
 M_i=\frac{2a_jd_j}{\gcd(g,2a_jd_j)},\\
 X_i\le b_i\le
 \left\lfloor\frac{4g a_iwd_i-16a_jd_j}{g}\right\rfloor,
 \qquad (X_1,X_2)=(8,16).
 \end{gathered}                                                       \tag{G026.5}
\]
Its coefficient in (G026.2) is $14-20g/(a_jd_j)$,
so the same endpoint rule minimizes it on this finite
congruence interval. Also impose $G_i(4)\ge48$ and
$S(4)\ge576$. These are all necessary conditions;
no divisibility of $B_i$ by $f_i$ is assumed.

The complete integer-only enumeration
\texttt{verify\_constant\_ratio\_global\_optimum.py}
has $1776$ scalar tuples after the coarse mass bounds,
and exactly $55$ endpoint rows below $T=13464$.
The full JSON certificate lists every row. Thirty-one
have $w=4$. In all remaining twenty-four rows
\[
 (a_1,g)=(2,1),\qquad a_2\in\{1,3\}.
\]
Their mass pairs are exactly
\[
\begin{array}{c|c|l}
 a_2&w&(d_1,d_2)\\\hline
 1&8&(4,12),(6,12),(6,18),(8,16)\\
 1&10&(4,8),(6,8),(6,12),(4,16),(6,16)\\
 1&12&(4,6),(4,8),(4,10),(6,6),(4,12),(6,12)\\
 1&14&(4,8)\\
 1&16&(4,8)\\
 3&8&(4,4),(6,4),(6,6),(10,4)\\
 3&10&(6,4)\\
 3&12&(4,4),(6,4)
\end{array}                                                         \tag{G026.6}
\]
The endpoint contents in the certificate need not equal
those of an actual hypothetical table. The exclusions
below either apply to every content in the row or derive
the actual contents anew from (G026.4).

\paragraph{Three elementary polynomial observations.}
First, coprime $f_1,f_2$ are impossible in (G026.1).
If they were coprime over $\mathbb Q$, polynomiality of
$U_1$ would force $f_2\mid a_1W_0-c_1D$ over $\mathbb Q[D]$.
Primitivity and Gauss's lemma make the quotient integral.
Reducing modulo two would say that $D(D-1)$ divides $D$,
which is false.

Second, when $w=4$, effectivity and parity force $W_0=D$,
and the two cells must be proportional. To see this,
write $f_1=Hu_1$, $f_2=Hu_2$ with primitive integer
polynomials and coprime $u_1,u_2$ over $\mathbb Q$.
The nonzero factors $a_iW_0-c_iD=(a_i-c_i)D$ in
(G026.1) force both $u_1,u_2$ to divide $D$.
A zero factor would make $U_i=0$, which is forbidden.
Since $f_1\equiv f_2\equiv\sigma\ne0$ modulo two,
the nonzero reduction of $H$ gives $u_1\equiv u_2$.
The choices $1$ and $D$ cannot be different, and they
cannot both be $D$ because $u_1,u_2$ are coprime.
Thus both are constant and $f_1=f_2$. NN-G025 excludes
all thirty-one $w=4$ rows below $4520$.

Third, the only effective parity-correct primitive
polynomials of masses four and six are respectively
\[
 f^{(4)}=D(D-3),\qquad f^{(6)}=(D-1)(D-2).
\]
They are coprime. Every such polynomial of mass eight
has the form
\[
 f^{(4)}+4\quad\text{or}\quad
 f^{(4)}+2(q^j+q^{-j})\quad(j\ge1).                   \tag{G026.7}
\]
No polynomial in (G026.7) vanishes at $D=0$ or $D=3$.
For $D=0$, choose $q=-1$, so the added pair has value
$4(-1)^j\ne0$. For $D=3$, choose $q=e^{i\pi/3}$;
the pair contributes $4\cos(j\pi/3)$, which is one
of $4,2,-2,-4$. The constant addition four also
never vanishes. Hence every mass-eight polynomial
is coprime to $f^{(4)}$, uniformly in its unbounded
support location $j$.

\paragraph{Parity modulo four eliminates the nonmultiples of four.}
For the remaining rows $a_1=2$ and $g=1$. Define
\[
 Q=(a_2W_0-c_2D)/2\in\mathbb Z[D],
\]
which is integral because $a_2,c_2$ are odd and
$W_0\equiv D$. The second $U$ identity is
$f_1U_2=f_2Q$. Reducing it modulo two and cancelling
$\sigma\ne0$ gives $Q\equiv D$. Therefore
\[
 a_2W_0\equiv(c_2+2)D\pmod4.                        \tag{G026.8}
\]
At $D=4$ this implies $w\equiv0\pmod4$.
Thus all $w=10$ and $w=14$ rows in (G026.6) are
excluded, regardless of their contents.

\paragraph{The rows with $a_2=1$.}
For $w=8$, positivity of $U_2(4)$ forces $c_2<2$,
so $c_2=1$. Every mass pair listed in this group has
$d_2\ge2d_1$. Then $U_1(4)\ge4$ in (G026.4) forces
$c_1\le2$, hence $c_1=1$ as well.
Equation (G026.8) gives $W_0\equiv3D\pmod4$.
At total mass eight there is exactly one effective
such polynomial:
\[
 W_0=3D-4,
\]
with coefficient list $(2,3)$; any additional outer
pair or increase of a noncentral coefficient modulo
four would already use at least eight further units
of mass.

Write $f_i=Hu_i$ as above. Polynomiality of the two
$U_i$ now forces
\[
 u_1\mid D-2,\qquad u_2\mid5D-8.
\]
Their reductions modulo two agree. Thus either both
are constant, which is the proportional case, or
$u_1=D-2$ and $u_2=5D-8$, up to the already normalized
positive primitive units. In the latter case the
first $U$ identity gives $U_1=D-2$. Its central
Laurent coefficient is zero, below the required
coefficient two of $D$. This excludes all four
$w=8$ rows.

For $w=12$, positivity of $U_2(4)$ gives $c_2<3$,
so again $c_2=1$. Equation (G026.8) requires
$W_0\equiv3D\pmod4$. An effective polynomial of
mass twelve with this residue is uniquely $3D$:
the minimum coefficients are $(2,3)$ of mass eight,
and the remaining four units can only increase the
central coefficient by four. But $3D$ has content
three, contradicting primitivity of $W_0$. All six
rows in this group are excluded.

The only $w=16$ row has mass pair $(4,8)$, whose
primitive cells are coprime by (G026.7). It is
excluded by the first polynomial observation.

\paragraph{The rows with $a_2=3$.}
After the modulo-four test, these have $w=8$ or
$w=12$. Mass pairs $(4,4)$ and $(6,6)$ force identical
primitive cells and are covered by NN-G025; $(6,4)$
is coprime and impossible. The only remaining pair
is $(d_1,d_2)=(10,4)$ at $w=8$.
In this row
\[
 U_1(4)=(10/12)(16-4c_1),\qquad
 U_2(4)=(24-4c_2)/5.
\]
Requiring both to be at least four and the contents
to be positive odd gives $c_1=c_2=1$.
Equation (G026.8) gives $W_0\equiv D\pmod4$.
At mass eight this uniquely gives $W_0=D+4$.
Now $f_2=D(D-3)$ and polynomiality of
\[
 U_1=\frac{f_1(D+8)}{3f_2}
\]
forces $f_2\mid f_1$, because $f_2$ is coprime
to $D+8$. The quotient is integral by Gauss's lemma.
Its value at $D=4$ would be $d_1/d_2=10/4$, which
is not an integer. This last row is impossible.

All fifty-five necessary-relaxation rows have been
excluded below $4520$. The NN-G025 attainer belongs
to the stated class, so the minimum is exactly
$4520$.\hfill$\square$

\paragraph{Geometric dictionary and the smaller remaining gate.}
The parity and modulo-four arguments are constraints
on actual effective graded sign representations:
dimension and signed trace agree modulo two, while
the composing table imposes further integral divisions.
SGA 3, I, 4.7.2--4.7.3 provides the diagonalizable-group
character dictionary, and the NN-G010 realization is
the finite-kernel composition recorded in archive
cote 2, scan 17. The evaluations in (G026.7) restrict
the grading character to finite roots of unity;
they are not Frobenius counts or arithmetic zeta zeros.
No rational auxiliary factor is promoted to an
effective object unless its actual coefficient has
been verified integral and nonnegative.

In geometric terms, all effective graded multiplicity objects
for the two cells may now vary independently, subject to this
closed kernel-table presentation and its constant rational
ratio identities. Among the resulting NN-G010 realizations,
the least total half-excess of dimension over the absolute
signed trace coefficients is exactly $4520$. The trace
specialization is fixed, but no canonical choice of objects,
coherence data, or generators is claimed. This optimum is
presentation-dependent, not a universal categorical invariant.

Arbitrary nonproportional cell polynomials and
unbounded support locations have now been treated
within the constant-ratio closed-table class. A
smaller score must use nonconstant ratios
$G_i/(W\delta_i)$ or a different table shape.
This is an exact realization theorem with explicit
scope, not a canonical-basis theorem or a consequence
for PH0, PH1, rank four, or nonstandard purity.

\subsection{NN-G027: nonconstant ratios produce smaller explicit realizations}
\paragraph{Statement.}
There is a closed-table effective associative dimension ring
with $3452$ cancelling pairs, of which $2852$ are in the
bottom cell. The addendum below improves these numbers to
$3352$ and $2752$, respectively. Both constructions have
nonconstant rational functions
$G_i/(W\delta_i)$, so both lie outside the exact $4520$
class of NN-G026. Thus the constant-ratio restriction is
substantive. No exact optimum over the enlarged rational-function
class is asserted.

\paragraph{Construction and associativity.}
Put
\[
 \sigma=D(D-3),\qquad h=\sigma+1,\qquad B=5\sigma+4.
\]
Choose
\[
 \begin{aligned}
 \delta_1&=\sigma+2,&\delta_2&=3\sigma+2,&W&=3D,\\
 G_1&=6Dh,&G_2&=DB,&S&=2DhB,\\
 U_1&=3D,&U_2&=D,&V_1&=15D-6,&V_2&=2D^2+D,\\
 A_1&=6D^2h-B(5D-2),&
 A_2&=D^2B-2h(2D^2+D),&
 \Lambda_i&=A_i/\delta_i.
 \end{aligned}                                                       \tag{G027.1}
\]
The actual within-cell coefficients involving $\Lambda_i$
are $A_i$ multiplied by nonnegative powers of $D$.
They are integral polynomials even though the auxiliary
$\Lambda_i$ are rational functions.

All three scalar-block identities follow directly:
\[
 \frac{G_1G_2}{S}=3D=W,\qquad
 G_1-D\delta_1=G_2,\qquad
 G_2-D\delta_2=G_1/3.
\]
Consequently $G_i(G_i-D\delta_i)/S=U_i$.
The two definitions of $A_i$ give
\[
 \frac{G_1(DG_1-A_1)}S=3(5D-2)=V_1,\qquad
 \frac{G_2(DG_2-A_2)}S=V_2.
\]
Also $\Lambda_i(D^2-\Lambda_i)\ne0$: the leading
terms of $\Lambda_1,\Lambda_2$ are respectively
$6D^2$ and $D^2/3$. The injective generic block
construction of NN-G011 therefore proves associativity
over $\mathbb Q(D)$. As all actual structure constants
are integral, its identities hold over $\mathbb Z[D]$.
This argument makes no semisimplicity claim at the
special values where an auxiliary factor has a pole.

\paragraph{Coefficient verification and exact score.}
The following symmetric Laurent lists give every base
coefficient needed for the effectivity, domination and
parity checks:
\[
\begin{array}{c|l|l}
 &\text{dimension coefficient}&\text{absolute trace target}\\\hline
 \delta_1&(2,1,1)&(0,1,1)\\
 \delta_2&(2,3,3)&(0,1,1)\\
 A_1&(14,5,1,5,6)&(2,1,1,1)\\
 A_2&(2,3,4,3,1)&(2,3,2,1,1)\\
 G_1&(24,24,18,6)&(16,12,4)\\
 G_2&(18,19,15,5)&(14,11,5,1)\\
 S&(172,152,114,78,40,10)&(144,120,68,24,4)\\
 V_1&(24,15)&(6,5)\\
 V_2&(14,9,2)&(6,3,2)
\end{array}                                                         \tag{G027.2}
\]
The remaining base coefficients are $W=U_1=3D$ and
$U_2=D$. Every entry in the middle column dominates
the right column coefficientwise with the same parity.
The closed multiplication table only multiplies these
base coefficients by nonnegative powers of $D$;
for signed traces use the coefficientwise triangle
inequality. Hence all $132$ positions satisfy the
required conditions. The included independent
integer-Laurent verifier also checks all $1000$
associativity triples and reports zero failures.

At $D=4$ the bottom mass is
\[
 960+20(120+96)+200\cdot12+20(12+4+54+36)=9800.
\]
The two cell characters have masses $6,14$, and the
actual $A_1,A_2$ have masses $48,24$. Thus their
contribution is $26(6+14)+14(48+24)=1528$.
The total is $T=11328$, yielding
\[
 (11328-4424)/2=3452,\qquad (9800-4096)/2=2852.
 \tag{G027.3}
\]
Finally
\[
 \frac{G_1}{W\delta_1}=\frac{2(\sigma+1)}{\sigma+2},
 \qquad
 \frac{G_2}{W\delta_2}=\frac{5\sigma+4}{3(3\sigma+2)},
\]
which are both nonconstant. The lower score is therefore
consistent with NN-G026, and proves that extending
its optimum beyond constant ratios would be false.
\hfill$\square$

\paragraph{A verified refinement to 3352 pairs.}
The same identities permit a smaller realization. Keep
$H=D^2-3D+1$, and now put
\[
 J=5D^2-17D+10.
\]
Define
\[
 \begin{aligned}
 \delta_1&=6H-J=D^2-D-4,&
 \delta_2&=J-2H=3D^2-11D+8,\\
 G_1&=6DH,&G_2&=DJ,&S&=2DHJ,\\
 W&=3D,&U_1&=3D,&U_2&=D,\\
 V_1&=15D,&V_2&=2D^2+D-4,\\
 A_1&=6D^2H-5DJ,&
 A_2&=D^2J-2H(2D^2+D-4),&\Lambda_i&=A_i/\delta_i.
 \end{aligned}                                                       \tag{G027.4}
\]
Again $G_1-D\delta_1=G_2$ and
$G_2-D\delta_2=G_1/3$, while the definitions of $A_i$
give the two $V_i$ identities. Neither
$\Lambda_i(D^2-\Lambda_i)$ is identically zero, so the
same injective block construction proves associativity.
The required base coefficients are
\[
\begin{array}{c|l|l}
 &\text{dimension coefficient}&\text{absolute trace target}\\\hline
 \delta_1&(0,3,1)&(0,1,1)\\
 \delta_2&(4,1,3)&(0,1,1)\\
 A_1&(6,5,1,5,6)&(2,1,1,1)\\
 A_2&(6,5,6,1,1)&(2,3,2,1,1)\\
 G_1&(24,24,18,6)&(16,12,4)\\
 G_2&(18,17,13,5)&(14,11,5,1)\\
 S&(156,140,106,70,36,10)&(144,120,68,24,4)\\
 V_1&(30,15)&(6,5)\\
 V_2&(10,9,2)&(6,3,2)
\end{array}                                                         \tag{G027.5}
\]
They prove every base condition; powers of $D$ propagate
them to the full table. The independent verifier
\texttt{verify\_dimension\_3352.py} checks all $132$
positions and all $1000$ triples with zero failures.

The bottom mass is
\[
 880+20(120+88)+200\cdot12+20(12+4+60+32)=9600.
\]
The cell mass is $26(8+12)+14(40+32)=1528$.
Consequently
\[
 T=11128,\qquad
 (11128-4424)/2=3352,\qquad(9600-4096)/2=2752.
 \tag{G027.6}
\]
The two ratios are $2H/(6H-J)$ and
$J/(3(J-2H))$, both nonconstant. This strengthens the
explicit upper bound to $3352$, retaining the original
$3452$ example as the first checked construction.
NN-G028 proves an exact optimum for the indicated
fixed-bottom subclass; no unrestricted optimum is
claimed here.\hfill$\square$

\paragraph{Geometric meaning and residual scope.}
The actual polynomials in (G027.2) are effective
Artin--Tate multiplicity objects and define another
NN-G010 finite-kernel monoidal realization of the
same trace table. With its chosen generators, the
sum over multiplication positions and grading pieces
of half the excess of dimension over absolute signed
trace is $3452$ for the first construction, and $3352$
for (G027.4), with bottom-cell contributions $2852$
and $2752$, respectively.
This is an explicit upper bound for a presentation-dependent
quantity, not a universal invariant of all categories
realizing the abstract algebra.

The rational $\Lambda_i$ and the two nonconstant ratios
are generic coordinate factorizations; no object of
fractional dimension or meromorphic sheaf is inferred.
Only the verified actual coefficients define objects.
The precise dictionary is SGA 3, I, 4.7.2--4.7.3 and
archive cote 2, scan 17 through NN-G010. The example
leaves the exact optimum for arbitrary rational ratios
and for other table shapes open. It does not give a
canonical family, a rank-four basis, PH0/PH1 or
nonstandard purity.

\subsection{NN-G028: fixed bottom coefficients force an exact 3352-pair optimum}
\paragraph{Statement.}
Among all effective associative closed rank-three tables with
\[
 W=3D,\qquad U_1=3D,\qquad U_2=D,                    \tag{G028.1}
\]
the exact cancelling-pair minimum is $3352$.
Cell polynomials, contents and degrees are arbitrary;
$G_i/(W\delta_i)$ and $\Lambda_i$ may be arbitrary
rational functions of $D$. Only the actual table
coefficients are required to be integral, effective,
dominant and parity-correct. There are exactly two
minimizing parameter sets within this closed presentation.
One is (G027.4); the other replaces its $V_2$ by
$2D^2-D+4$ and adjusts $A_2$ accordingly.

The fixed coefficients (G028.1) are a substantive
restriction. This theorem closes a rational-function
optimization problem inside that restriction, and
makes no lower-bound assertion when $W$ or the two
$U_i$ vary.

\paragraph{A forced two-polynomial parameterization.}
The scalar identities used in NN-G011 are necessary
for associativity in this closed table: associativity
with $\Sigma$ gives
\[
 WS=G_1G_2,\qquad U_iS=G_i(G_i-D\delta_i),\qquad
 V_iS=G_i(DG_i-A_i),\quad A_i=\delta_i\Lambda_i.
\]
For example, $\Sigma(\gamma_1^1\gamma_1^2)$ gives
the first identity; multiplication by $\Sigma$ of
a same-cell two-letter product gives the second,
and of a three-letter reduction gives the third.
Thus no generic-embedding assumption is added in
the lower-bound argument.

Substituting (G028.1) yields
\[
 G_1-D\delta_1=G_2,\qquad
 3(G_2-D\delta_2)=G_1.
\]
Consequently
\[
 \begin{gathered}
 H=(\delta_1+\delta_2)/4,\qquad
 J=(\delta_1+3\delta_2)/2,\\
 G_1=6DH,\quad G_2=DJ,\quad S=2DHJ,\qquad
 \delta_1=6H-J,\quad\delta_2=J-2H.                 \end{gathered}\tag{G028.2}
\]
Both $H,J$ are effective integral symmetric Laurent
polynomials, hence belong to $\mathbb Z[D]$.
For $J$, integrality follows directly from
$\delta_1\equiv\delta_2$ modulo two. To prove it for
$H$, first put $K=(\delta_1+\delta_2)/2\in\mathbb Z[D]$.
The actual coefficient $G_1=3DK$ is zero modulo two,
so cancellation in $\mathbb F_2[D]$ gives $K\equiv0$;
therefore $H=K/2$ is integral. Effectivity follows
from their displayed positive combinations of the cells.
The $G_2$ parity gives $J\equiv D(D-1)$ modulo two.

Let $\sigma=q^2+q+q^{-1}+q^{-2}=D(D-3)$ denote
the absolute cell trace coefficient. The actual cell
bounds in (G028.2) are exactly
\[
 2H+\sigma\le J\le6H-\sigma                       \tag{G028.3}
\]
coefficientwise. In particular $H$ has coefficients
at least one at each of $q^{\pm1},q^{\pm2}$.
Write
\[
 m=H(4)\ge4,\qquad n=J(4)\ge2m+4.
\]
Here $m$ is an integer and $n$ is a positive even integer.
No degree or Laurent-support bound has been imposed.

\paragraph{A mass inequality removes every $m\ge6$.}
Let $v_i=V_i(4)$ and $a_i=A_i(4)$ in this paragraph;
these symbols are masses, not the ratio numerators of
NN-G026. The actual coefficients obey
\[
 a_1=96m-(n/3)v_1\ge8,\qquad
 a_2=16n-2mv_2\ge16.
\]
The exact total mass becomes
\[
 \begin{split}
 T={}&8mn+584m+80n+2720+
       \frac{5760m}{n}+\frac{160n}{m}\\
    &+\left(14-\frac{60}{n}\right)a_1
       +\left(14-\frac{10}{m}\right)a_2.
 \end{split}                                                        \tag{G028.4}
\]
Both coefficients on the last line are positive,
since $m\ge4$ and $n\ge12$. Therefore
\[
 T\ge L(m,n):=8mn+584m+80n+3056
       +\frac{5760m-480}{n}+\frac{160n-160}{m}.       \tag{G028.5}
\]
For $m\ge6$ and $n\ge2m+4$, this is minimized at
$(m,n)=(6,16)$. Here is a direct verification of that
global assertion, rather than a bounded numerical test.
The derivative with respect to $n$ is
\[
 8m+80+160/m-(5760m-480)/n^2.
\]
Its value at $n=2m+4$, multiplied by $m(m+2)^2$,
is
\[
 8\{m^2(m^2+14m-116)+135m+80\}>0
 \quad(m\ge6).
\]
It is increasing in $n$, so $L(m,n)\ge L(m,2m+4)$.
The latter expression is
\[
 16m^2+776m+6576-\frac{6000}{m+2}+\frac{480}{m},
\]
whose derivative
$32m+776+6000/(m+2)^2-480/m^2$ is positive for
$m\ge6$. Thus
\[
 T\ge L(6,16)=11138>11128.                          \tag{G028.6}
\]
The new attainer has total mass $11128$, so every
competing minimizer has $m\le5$.

If $m=4$, (G028.3) forces $H=\sigma$.
Then $[q^0]G_1=[q^0](6D\sigma)=12<16$, violating
its trace target. If $m=5$, the only possible addition
to those four mandatory coefficients is one central
unit, so
\[
 H=(1,1,1)=D^2-3D+1.                                \tag{G028.7}
\]
This proves a degree reduction from the score and
actual coefficient inequalities; it does not infer
a degree bound from mass alone.

\paragraph{Only four $J$ polynomials remain.}
By (G028.3) the support of $J$ is contained in that
of $H$, so write its symmetric coefficient list as
$(j_0,j_1,j_2)$. Its bounds and parity give
\[
 j_0\in\{2,4,6\},\qquad j_1,j_2\in\{3,5\}.
\]
Effectivity of $A_i=DG_i-(G_j/W)V_i$ forces
$\deg V_i\le2$: above that degree its highest term
would have negative sign, since $H,J,V_i$ have
positive leading coefficients. The coefficient of
$V_2$ at $q^2$ is even and at least two. The top
coefficient of $A_2=D^2J-2HV_2$ is therefore at most
$j_2-4$, forcing $j_2=5$; its $V_2$ coefficient is
then exactly two.

The central coefficient of $S=2DHJ$ equals
$8j_0+16j_1+60$ and must be at least $144$.
This leaves precisely
\[
 (j_0,j_1,j_2)=(2,5,5),(4,5,5),(6,3,5),(6,5,5).
 \tag{G028.8}
\]
Each $J$ in this list is primitive. Since
$A_1=6D^2H-(J/3)V_1$ is integral, Gauss's lemma
forces $V_1=3C$ with $C$ an integral effective
polynomial. The parity and target of $V_1$ give
$C_0\ge2$ even, $C_1\ge3$ odd, and $C_2$ even.
The leading inequality $6-5C_2\ge0$ forces
$C_2=0$. Thus $C=(x,y)$ with $x\ge2$ even and
$y\ge3$ odd.

\paragraph{Complete coefficient optimization of the two remaining factors.}
The symmetric list of $6D^2H$ is
$(96,90,66,30,6)$. Subtracting $JC$ and comparing
with $X_1=(2,1,1,1)$ gives exactly
\[
 \begin{gathered}
 y\in\{3,5\},\qquad
 5x+j_1y\le65,\\
 j_1x+(j_0+5)y\le89,\qquad
 j_0x+2j_1y\le94.                                  \end{gathered}\tag{G028.9}
\]
For fixed $y$, the largest even $x$ satisfying these
three bounds maximizes $C(4)=x+2y$.

Write $V_2=(z,t,2)$, where $z\ge6$ is even and
$t\ge3$ is odd. Expanding $D^2J-2HV_2$ and comparing
with $X_2=(2,3,2,1,1)$ gives exactly
\[
 \begin{gathered}
 t\le(j_1+15)/2,\qquad
 z+2t\le3j_0+4j_1,\\
 z+2t\le2j_0+(7j_1+13)/2,\qquad
 z+t\le j_0/2+2j_1+12.                             \end{gathered}\tag{G028.10}
\]
Here $t$ ranges only over $3,5,7,9$. For fixed $t$
take the largest allowed even $z$. These lists
include every coefficient condition: the remaining
top coefficients are $6$ and $1$, and parity is
already enforced. The total mass is strictly decreasing
in $v_1$ and $v_2$, with coefficients
$20-14n/3$ and $20-28m$, so maximizing these masses
independently is the exact minimizing rule.

The finite arithmetic is displayed in full by
\texttt{verify\_fixed\_bottom\_optimum.py}; its four
rows are
\[
\begin{array}{c|c|c|c|c|r}
 J&\max C(4)&\max V_2(4)&a_1&a_2&\text{score}\\\hline
 (2,5,5)&18&30&84&52&3720\\
 (4,5,5)&18&36&48&24&3452\\
 (6,3,5)&20&32&40&32&3352\\
 (6,5,5)&16&38&64&36&3728
\end{array}                                                         \tag{G028.11}
\]
There is only one minimizing $J$. In that row (G028.9)
forces the maximizing $C=(10,5)$, so $V_1=15D$.
Equation (G028.10) has exactly two maximizing choices,
\[
 V_2=(10,9,2)=2D^2+D-4,
 \quad\text{or}\quad
 V_2=(14,7,2)=2D^2-D+4.
\]
The first is the verified construction (G027.4).
For the second the actual $A_2$ list is
$(6,5,2,5,1)$; it has the same mass and satisfies
all required conditions. The other data are unchanged,
and the generic block construction again proves
associativity. This proves both attainment and the
complete two-parameter-set minimum in the class
(G028.1).\hfill$\square$

\paragraph{Geometric translation and the next exact gate.}
Fixing (G028.1) fixes three actual multiplicity objects
in the NN-G010 finite-kernel presentation. Among all
its effective graded sign-representation tables with
those fixed objects and the specified trace specialization,
the sum over multiplication positions and grading pieces
of half the dimension-minus-absolute-trace excess is
exactly $3352$. There are two minimizing character-valued
parameter sets in this presentation; no uniqueness of
categories or coherence data is asserted. This is an
optimization of chosen multiplication data, not an
intrinsic minimum for every categorification of the
abstract trace algebra.

The exact source dictionary remains SGA 3, I,
4.7.2--4.7.3 for effective grading characters and
archive cote 2, scan 17 for kernel composition.
The mass reduction respects actual coefficient
objects and never treats a rational auxiliary factor
as an effective object. A smaller closed-table score
must change at least one of $W,U_1,U_2$ in (G028.1);
other table shapes are a further separate possibility.
No canonicality, rank-four, PH0/PH1 or purity theorem
is inferred from this exact restricted optimum.

\subsection{NN-G029: an existence-preserving degree cutoff makes the unrestricted optimum decidable}
\paragraph{Statement.}
Fix a finite based trace algebra with bar-invariant Laurent
structure constants $c_{ij}^k(q)\in\mathbb Z[q^{\pm1}]$,
including a specified unit. An effective symmetric dimension
lift consists of bar-invariant $d_{ij}^k(q)\in\mathbb N[q^{\pm1}]$
that are associative and unital, dominate $|c_{ij}^k|$
coefficientwise, and have the same coefficients modulo two.
Extra positions with zero trace are permitted unless explicitly
forbidden. Its cancelling-pair cost is
\[
 C(d)=\frac12\sum_{i,j,k,e}
             \bigl(d_{ij,e}^k-|c_{ij,e}^k|\bigr).
\]
Suppose one such lift of cost at most $C_0$ is known.
Put
\[
 N_0=\sum_{i,j,k,e}|c_{ij,e}^k|,\qquad N=N_0+2C_0,
 \qquad K=\max\bigl(1,\max\{|e|:c_{ij,e}^k\ne0\}\bigr),
\]
Let $E$ consist of the positive trace exponents together with one,
put $h=|E|$, and set $n_* =\lfloor N/2\rfloor+h$.
Every lift of cost at most $C_0$ has another lift of exactly
its cost, with the same coefficient lists and the same nonzero
structure-constant positions, whose Laurent degrees are at most
\[
 B=n_*(K+1)4^{\,n_*-1}.                              \tag{G029.1}
\]
Consequently a finite terminating search computes the true
minimum over all symmetric effective associative lifts in
this presentation, without a closed-table or constant-ratio
restriction.

This is an existence-preserving compression statement.
It does \emph{not} bound the degrees of every existing lift,
assert a practical search algorithm, or identify the numerical
unrestricted optimum.

For the supplied rank-three trace table, $N_0=4443$ and $K=5$.
The NN-G027 incumbent has $C_0=3352$, so $N=11147$ and
$n_*=5578$. Every competing lift therefore has an equal-score
representative of Laurent degree
\[
 \le33468\,4^{5577}<2^{11170}.                       \tag{G029.2}
\]
One may search all $10^3$ structure-constant positions, including
new even coefficients at positions of zero trace, or retain any
specified forbidden positions. Both variants are finite and
complete for their stated presentation.

\paragraph{Elementary small-integer-point lemma.}
Let $P\subset\mathbb R^n$ be a nonempty polyhedron defined by
finitely many integer inequalities $a_r\cdot x\ge b_r$.
Equalities may be represented by their two inequalities.
Assume $P\subseteq\{x_i\ge1\}$, every row has Euclidean
norm at most four, and $|b_r|\le K$ with $K\ge1$.
If $P$ contains an integer point, it contains one satisfying
\[
 1\le x_i\le n(K+1)4^{n-1}\quad\text{for every }i.
 \tag{G029.3}
\]

Here is a complete polyhedral proof. Include the inequalities
$x_i\ge1$ in the presentation and homogenize:
\[
 \mathcal C=\{(x,t):t\ge0,\ a_r\cdot x\ge b_rt\}.
\]
Because $x_i\ge t\ge0$, this cone lies in the nonnegative
orthant and has no nonzero line. Its section
$\sum_i x_i+t=1$ is a compact polytope. Every point of
this section is a convex combination of its vertices:
if its smallest face has positive dimension, choose a
nonzero direction in that face's affine span, follow it
in both directions to the two boundary endpoints of
the compact face, and write the point as their convex
combination. The endpoints lie in faces of smaller
dimension, so induction terminates. There are only
finitely many vertices, since each is determined by
a full-rank subset of the finite active constraints.
Thus $\mathcal C$ is generated by finitely many extreme
rays.

An extreme ray with $t>0$, rescaled to $t=1$, has the
form $(v,1)$ with $v$ a vertex of $P$: a nontrivial
segment through $v$ in $P$ would contradict extremality
of that ray. An extreme ray with $t=0$ has the form
$(r,0)$ with $r$ an extreme ray of the recession cone
\[
 \operatorname{rec}(P)=\{r:a_r\cdot r\ge0\}.
\]
In particular $r_i\ge0$. Applying the cone decomposition
to an existing integer point $(x,1)$ gives
\[
 x=\sum_j\lambda_jv_j+\sum_s\mu_s r_s,
 \qquad \lambda_j,\mu_s\ge0,\quad\sum_j\lambda_j=1.
 \tag{G029.4}
\]
This proves the needed vertex-plus-recession decomposition
also when $P$ has lower affine dimension. It explains
explicitly the rays on the $t=0$ face; none is lost.

A vertex $v$ has $n$ linearly independent active integer
rows. Otherwise a small positive and negative displacement
in their common kernel would remain feasible, contradicting
that it is a vertex. Orthogonalizing the rows by Gram--Schmidt shows that the
absolute determinant is the product of the orthogonalized
row lengths, each at most its original length; this proves
the determinant bound used below. Cramer's rule then gives
\[
 |v_i|\le nK4^{n-1}:                                  \tag{G029.5}
\]
the denominator is a nonzero integer of absolute value at
least one, and each numerator is a sum of $n$ numbers
$b_r$ times cofactors bounded by $4^{n-1}$.

Likewise an extreme recession ray has $n-1$ independent
active homogeneous rows. Otherwise its active common
kernel would have a direction transverse to the ray,
and small displacements would split it inside the cone.
Its integer cofactor vector spans the ray. Choose the
nonnegative orientation and divide by the greatest common
divisor of its entries to obtain a primitive integer
ray $r$, with
\[
 0\le r_i\le4^{n-1}.                                 \tag{G029.6}
\]
The assertion also covers $n=1$, using the empty determinant
one. If the recession cone is zero there are simply no
such rays.

The ray part of (G029.4) needs at most $n$ rays. If more
are used, they are linearly dependent. A nonzero linear
relation among nonnegative nonzero rays has coefficients
of both signs. Subtract a suitable multiple of that
relation from the nonnegative coefficients until one
coefficient becomes zero, preserving all the others'
nonnegativity, and repeat. The vertex convex combination
is left unchanged.

Now replace each of these at most $n$ coefficients $\mu_s$
by its fractional part. Equivalently set
\[
 x'=x-\sum_s\lfloor\mu_s\rfloor r_s.
\]
This is an integer vector, because $x$ and the chosen
primitive rays are integral. It remains in $P$, because
it is the same convex combination of vertices plus a
nonnegative recession combination. The coefficients in
that combination are less than one, so (G029.5)--(G029.6)
give
\[
 x'_i\le nK4^{n-1}+n4^{n-1}=n(K+1)4^{n-1}.
\]
Since $P\subseteq\{x_i\ge1\}$, this proves (G029.3).
Notice that no rounding of the vertex convex combination
was performed. Integer flooring was applied only to
integer recession directions.\hfill$\square$

\paragraph{Encoding every relevant exponent collision and fixed anchor.}
Take any lift $d$ of cost at most $C_0$. Its total
coefficient mass is $N_0+2C(d)\le N$. Keep every
coefficient and every structure-constant position fixed.
For each occurrence of a positive exponent in a symmetric
pair $a(q^e+q^{-e})$, introduce one positive integer
variable $x_\nu$, and replace the pair's exponents by
$x_\nu,-x_\nu$. There are at most $\lfloor N/2\rfloor$
such variables: each pair uses coefficient mass at least
two. Central coefficients need no variable. Distinct
occurrences may use separate variables even if their
original numerical exponents agree.

Adjoin one anchor variable $y_r$ for each $r\in E$ and require
\[
 y_r=r\quad(r\in E).
\]
Thus the total number of variables is at most $n_*$.
There are $h$ such variables. These are affine integer equalities with constants
bounded by $K$; they will not be rescaled. Every recession
direction has zero anchor coordinates, so the flooring
step preserves them literally.

Form the finite collection $\mathcal L$ of all linear
forms $0$, all positive and negative individual variables,
and all sums of two such signed variables. Repetitions
such as $x_\nu+x_\nu$ are allowed. Each form has coefficient
$\ell^1$ norm at most two. It includes every exponent
appearing in a product of two structure constants,
as well as every single exponent and every fixed trace
weight $\pm r$ through its anchor.

Use the original integer exponent assignment to record
the complete order pattern of these forms. If two forms
have the same value, impose their equality. If one is
larger, impose that their difference be at least one.
Also impose positivity of every variable and the fixed
anchor equations. These are finitely many integer linear
constraints, and the original assignment is an integer
solution. Each non-anchor row has $\ell^1$ norm at most
four and right side zero or one; an anchor row has norm
one and right side at most $K$ in absolute value.
Consequently the small-integer-point lemma supplies
another solution with all coordinates at most (G029.1).

This construction preserves the exact order pattern,
not only a set of selected equalities. In particular it
prevents previously different sums from colliding and
prevents a non-anchor monomial from newly occupying a
fixed trace weight. The constraints also retain equalities
between distinct occurrences of the same original weight.

\paragraph{Why the new exponent assignment is another lift of the same cost.}
Use the small solution to redefine the monomial exponents,
leaving their coefficients unchanged. Positivity of the
variables keeps each symmetric pair away from exponent
zero; their negatives still give bar invariance. The
single-exponent order constraints keep distinct monomials
within each coefficient distinct. Thus each coefficient
list, each nonzero table position, and the total mass are
unchanged.

Domination and parity at the prescribed trace weights
are unchanged, because equality or inequality with every
anchor $\pm y_r$ and with zero was preserved and the
anchors still equal the original integers $r$.
Every monomial outside the fixed trace support stays
outside it and retains its even coefficient. The specified
unit has only its unchanged constant coefficient, so
unitality is preserved.

For associativity, fix three basis elements and an output
basis element. Each side of
\[
 \sum_h d_{ij}^h d_{hk}^{\ell}
       =\sum_h d_{jk}^h d_{ih}^{\ell}
\]
is a sum of monomial products. Their numerical coefficients
are products of the original unchanged coefficients.
Their exponents are forms in $\mathcal L$. The original
identity says that the coefficient totals agree in every
exponent-collision class. Exactly the same classes occur
after compression, so every coefficient total still
agrees. This verifies all associativity identities.
The trace coefficient norm is fixed and the dimension
mass is unchanged, hence the cancelling-pair cost is
exactly the same. This proves (G029.1).

\paragraph{The finite algorithm and its exact scope.}
There are only finitely many arrays of nonnegative integer
Laurent coefficients with total mass at most $N$ and
exponents in $[-B,B]$. Enumerate them, impose bar symmetry,
the specified unit, coefficient domination and parity,
and test associativity by exact finite polynomial
arithmetic. Optional prescribed zero positions are
retained throughout. The known incumbent ensures the
accepted set is nonempty. Minimize the cost over it.

Any better unbounded-degree lift would compress to an
accepted lift of the same cost in this finite set.
Therefore the computed minimum is the unrestricted
symmetric-lift minimum for the chosen basis and unit,
not merely a degree-window minimum. The search bound is
immense: neither polynomial time nor any useful running
time estimate follows from this argument.

For the numerical specialization, NN-G011 proves the
absolute nonunit trace mass $4424$, and the nineteen
unit positions add $19$, giving $N_0=4443$.
The supplied trace coefficient table has maximum
Laurent exponent five. The checked NN-G027 lift has
mass $11128+19=11147$. Hence $n_*\le5578$ and
$n_*(K+1)4^{n_*-1}\le33468\,4^{5577}<2^{11170}$.
The small certificate
\texttt{verify\_degree\_cutoff\_constants.py} checks
these input norms and constants; it does not perform
the enormous finite search or substitute for the proof.

\paragraph{Sparse anchors give polynomial-size certificates without an incumbent.}
The anchor set need only contain the positive exponents
actually present in the trace, together with one; all
other trace coefficients are zero. With
\[
 E=\{1\}\cup\{e>0:c_{ij,e}^k\ne0\text{ for some }i,j,k\},
 \qquad h=|E|,
\]
the proof uses $h$ anchor variables, not $K$ variables.
The same small-point lemma applies with right-hand-side
bound $K=\max E$, and gives
\[
 n\le\lfloor N/2\rfloor+h,\qquad
 \log_2 B=O(N+\log K).                              \tag{G029.7}
\]
Indeed $h\le\lfloor N_0/2\rfloor+1\le N+1$ whenever
the trace norm is at most the allowed dimension mass.
Preserving comparison with every exponent in $E$ is
sufficient: a monomial at a nonzero trace weight stays
there, and a monomial outside all those weights remains
outside the nonzero trace support. Adding the weight-one
anchor fixes the marked character used in the geometric
interpretation. The rank-three numerical bound is unchanged,
since its positive trace exponents are exactly $1,\ldots,5$.

This proves a uniform NP-membership statement. The input
is a finite based trace table encoded sparsely: rank and
basis indices, integer coefficients and signed exponents
are written in binary, and a \emph{total coefficient-mass
budget $N$} is written in unary. Equivalently one may
consider an input regime in which $N$ is polynomially
bounded by the input length. The decision problem asks
whether a symmetric effective associative unital lift
of total mass at most $N$ exists. Trace-algebra identities
and normalization of its sparse coefficient lists can
be checked as part of the input verification.

No known incumbent is needed for this decision result.
If the answer is yes, apply the compression proof to
that yes-witness with the prescribed budget $N$.
A unit on rank $r$ already uses at least $2r-1$ coefficient
units, so a yes-instance has $r\le(N+1)/2$. The trace
norm must also be at most $N$; otherwise the answer is
immediately no. These observations make both the rank
and the sparse anchor count polynomially bounded.

The compressed witness has at most $N$ nonzero monomials.
Their coefficients have $O(\log N)$ bits, their table
positions have $O(\log r)$ bits, and (G029.7) gives
$O(N+\log K)$ bits per exponent. Sorting the monomials
inside each table coefficient gives a canonical sparse
encoding of polynomial length. A verifier checks total
mass, bar symmetry, the unit, domination and parity by
sparse comparisons with the input table. It checks all
associativity identities by multiplying sparse polynomials,
adding their binary exponents, and collecting equal sums.
There are polynomially many such products and terms in
$r,N$; all coefficients and exponent sums have polynomial
bitlength. Thus the full verification is polynomial in
the input and witness lengths, proving membership in NP.

This is a certificate theorem for the existence of an
effective associative lift under a mass budget. It is
not NP-hardness, a polynomial-time algorithm, or PH1 for
actual plethysm or Kronecker multiplicities. With a binary
budget that may be exponentially large, the stated
polynomial certificate bound does not follow.

\paragraph{Geometric translation and the reduced optimization gate.}
The exponent variables are characters in a common finite
presentation of a diagonalizable group. More explicitly,
start with a free character group on these symbols and
one marked symbol $e_0$. Impose the homogeneous character
relations corresponding to the exponent collisions and
the anchor relations $y_r-r e_0=0$. An integral cocharacter
to this diagonalizable group, with value one on $e_0$,
assigns the integer grading weights. The strict order
conditions select an open rational polyhedral region
in that affine cocharacter slice. The original and the
compressed solutions are two integral cocharacters in
the same finite collision region, both fixing the marked
trace weights.

SGA 3, I, 4.7.2--4.7.3 identifies these character
presentations with diagonalizable-group comodules.
Pulling back their finite multiplicity objects along
the two cocharacters preserves the displayed finite
multiplication relations; NN-G010 and archive cote 2,
scan 17 then give the corresponding finite-kernel
realizations. Thus the actual conclusion is existence
of an equal-cost realization with bounded grading
weights and the same specified trace table. No equivalence
of the two categories or canonical choice of grading
is asserted.

The reassignment is not a homomorphism of the original
whole weight lattice $\mathbb Z$ to itself, and is not
a single Veronese scaling: the weight-one anchor remains
one. It preserves the finite character relations needed
for this multiplication presentation. The unrestricted
symmetric-lift optimum is now a finite effective search
problem, and unary-mass-budget existence has polynomial-size
verifiable certificates. Its value and any practical
optimization method remain open. Nothing here computes actual GCT multiplicities,
proves PH0/PH1, or transfers a nonstandard purity statement.

\subsection{NN-G030: associativity forces modulo-four constraints beyond trace parity}
\paragraph{Statement.}
For every integral associative closed rank-three table with
the specified trace parities, put $\sigma=D(D-1)$ and
$A_2=\delta_2\Lambda_2$, the actual second cell-product
coefficient. No constant-ratio, degree, or auxiliary-integrality
hypothesis is needed. Then the following are polynomial
congruences in $\mathbb Z[D]$:
\[
 S\equiv\sigma G_1,\qquad
 G_2-D\delta_2\equiv G_1,\qquad
 DG_2-A_2\equiv G_1\pmod4.                           \tag{G030.1}
\]
In fact
\[
 \begin{aligned}
 G_1&\in(4,2D),\\
 G_2,A_2&\in(4,2D,D^2),\\
 S&\in(4,2D^2(D-1)).
 \end{aligned}                                                       \tag{G030.2}
\]
Thus $G_1,G_2,S,A_2$ take values divisible by four
at every even integer $D$, and $S$ does so at every
integer $D$. In particular their ordinary dimension
masses at $D=4$ are multiples of four.
These necessary conditions use associativity and
integral trace parity alone, before effectivity or
domination is tested.

\paragraph{Proof.}
The supplied trace table gives
\[
 \begin{gathered}
 W\equiv D,\quad\delta_2\equiv\sigma,\quad
 G_1\equiv S\equiv0,\quad G_2\equiv D\sigma,\\
 U_2\equiv V_2\equiv D,\qquad A_2\equiv D^2\sigma
       \pmod2.
 \end{gathered}
\]
Write $G_1=2H$ and $S=2K$, and define the integral
polynomials
\[
 P=(G_2-D\delta_2)/2,\qquad
 Q=(DG_2-A_2)/2.
\]
The three necessary associativity identities reduce to
\[
 WK=HG_2,\qquad U_2K=G_2P,\qquad V_2K=G_2Q.
\]
Reducing the first modulo two and cancelling the
nonzero polynomial $D$ in $\mathbb F_2[D]$ gives
$\overline K=\overline H\sigma$.
The second and third then give
\[
 D\overline H\sigma=D\sigma\overline P
                       =D\sigma\overline Q.
\]
Cancelling $D\sigma\ne0$ gives
$\overline P=\overline Q=\overline H$.
Multiplying these statements by two proves (G030.1).
This does not require $\overline H$ itself to be nonzero.

At $D=0$, (G030.1) says
\[
 S(0)\equiv0,\quad G_2(0)\equiv G_1(0),\quad
 A_2(0)\equiv-G_1(0)\pmod4.
\]
The integers $G_1(0),G_2(0)$ are even, and $W(0)$
is even. Since $WS=G_1G_2$, their product at zero
is divisible by eight. If their common residue
modulo four were two, their product would instead
be four modulo eight. Hence $G_1(0),G_2(0),A_2(0)$
are all divisible by four.

Every coefficient of $G_1$ is even, so its divisible-by-four
constant term gives $G_1\in(4,2D)$. The parities
$G_2\equiv D^2(D-1)$ and
$A_2\equiv D^3(D-1)$ show that their linear
coefficients are even; their constant terms are
multiples of four. This gives their ideal membership
in (G030.2). Finally substitute $G_1=4a+2DB$
in the first congruence of (G030.1):
\[
 S\equiv2D^2(D-1)B\pmod4,
\]
which proves the last ideal membership. Evaluation
at even integers gives the claimed divisibilities.
For arbitrary integer $D$, the product $D^2(D-1)$
is even, proving the additional assertion for $S$.
\hfill$\square$

\paragraph{A low-mass scalar candidate that these constraints exclude.}
Ordinary mass parity and the scalar associativity
identities alone admit the following hypothetical data:
\[
 \begin{array}{c|rrrrrrrrrrrr}
 &w&g_1&g_2&s&d_1&d_2&u_1&u_2&v_1&v_2&a_1&a_2\\\hline
 \text{values}&10&70&90&630&4&12&6&6&30&48&10&24
 \end{array}
\]
Here lower-case symbols denote evaluations at $D=4$,
with $a_i=A_i(4)$. They obey all target mass lower
bounds and are even. Directly,
\[
 ws=g_1g_2,
 \qquad u_i s=g_i(g_i-4d_i),\qquad
 v_i s=g_i(4g_i-a_i),
\]
For example both $u_i s$ equal $3780$;
$v_1s=18900$ and $v_2s=30240$.
Their relaxed total mass would be $8522$, or
$(8522-4424)/2=2049$ pairs. But $g_1=70$, $g_2=90$
and $s=630$ are not multiples of four. Thus no integral
polynomial table with the required trace parity can
realize these scalar data. This is a strict exclusion
from the scalar relaxation, not an assertion that the
hypothetical data were an actual lower-cost lift.

The constraints do not say that every coefficient
of $G_1$ or $S$ is divisible by four. In the actual
$3352$ attainer the coefficient of $D$ in $G_1$
is six and the coefficient of $D^2$ in $S$ is
$-94$. Both are two modulo four. The polynomial
ideals in (G030.2), rather than coefficientwise
divisibility by four, are the proved statements.
The small exact checker verifies the displayed
congruences on that attainer and verifies all scalar
identities in the excluded mass example.

\paragraph{Geometric translation and algorithmic scope.}
At the level of integer congruences, individual effective
graded sign representations require dimension and signed
trace to agree modulo two.
Associativity of their finite-kernel composition imposes
the additional congruences (G030.1). In the symmetric
character ring $\mathbb Z[D]$, they say that the listed
coefficient characters vanish in the explicit quotient
rings $\mathbb Z[D]/(4,2D)$,
$\mathbb Z[D]/(4,2D,D^2)$, and
$\mathbb Z[D]/(4,2D^2(D-1))$, as appropriate.
These are exact scheme-theoretic congruence statements
over the nonreduced characteristic-four base, not
purity or point-counting assertions.

The character interpretation is SGA 3, I, 4.7.2--4.7.3;
$D=4$ is evaluation at the identity grading character,
while $D=0$ is evaluation at $q=-1$.
The composing objects remain the NN-G010 kernel
realization of archive cote 2, scan 17. These conditions
provide inexpensive necessary tests before a bounded
search or NP verification, but are not sufficient for
an effective lift and do not determine its optimum,
canonicality, rank-four positivity, or PH0/PH1.

\subsection{NN-G031: a competitive closed table must have gluing mass at least eight}
\paragraph{Statement.}
Every effective associative symmetric closed table of the specified trace
with at most $3352$ cancelling pairs has $W(4)\ge8$.
This permits arbitrary cell degrees and contents and arbitrary
rational ratios and auxiliary parameters, subject only to
integrality and effectivity of the actual coefficients.
In particular the entire $W=D+2$ branch in this presentation is excluded, without
enumerating support locations or assuming a degree bound.

\paragraph{Mass notation and immediate bounds.}
Put $\psi=D(D-3)=q^2+q+q^{-1}+q^{-2}$, the absolute
cell trace, and use lower-case letters for evaluation at $D=4$.
For $A_i=\delta_i\Lambda_i$, write $a_i=A_i(4)$.
All these masses are positive even integers. The trace imposes
\[
 d_i\ge4,\quad u_i\ge4,\quad v_i\ge16,\qquad
 a_1\ge8,\quad a_2\ge16.
\]
The nonunit mass, with the unit contributions cancelled from
the pair count, is
\[
 T=s+20(g_1+g_2)+200w+20(u_1+u_2+v_1+v_2)
       +26(d_1+d_2)+14(a_1+a_2).                    \tag{G031.1}
\]
The proposed score cap is exactly $T\le11128$.
Associativity gives
\[
 ws=g_1g_2,\qquad u_i g_j=w(g_i-4d_i)\quad(i\ne j).
                                                               \tag{G031.2}
\]
Consequently $u_1u_2<w^2$. Since $u_i\ge4$, $w=4$
is impossible. Parity and domination give $w\ge4$ even,
so it remains to exclude $w=6$.

An effective symmetric polynomial of mass six dominating
$D=q+2+q^{-1}$ and congruent to it modulo two is uniquely
$W=D+2$: its surplus mass two must be central.
Similarly $u_i=4$ forces $U_i=D$, and $u_i=6$ forces
$U_i=D+2=W$.

\paragraph{A useful divisibility consequence of $U_i=D$.}
For this $W$, the identities are
\[
 U_1G_2=W(G_1-D\delta_1),\qquad
 U_2G_1=W(G_2-D\delta_2).                           \tag{G031.3}
\]
If $U_1=D$, reduction of the first at $D=0$ implies
$D\mid G_1$, and coprimality of the monic polynomials
$D,W$ implies $W\mid G_2$. Write $G_1=DR$, $G_2=WT$;
all quotients are integral. The second identity at $D=0$
then implies $T(0)=0$, because $W(0)^2=4$.
Thus
\[
 G_2=DW C_2,\qquad G_1=D\delta_1+D^2C_2.
\]
The parity $G_1\equiv0$ and
$\delta_1\equiv D(D-1)$ in $\mathbb F_2[D]$ gives
$C_2\equiv D-1$ after cancellation of $D^2$.
Therefore $g_2/24$ is an odd integer.
The same argument with indices interchanged shows that
$U_2=D$ implies
\[
 G_1=DW C_1,\qquad G_2=D\delta_2+D^2C_1,
 \qquad C_1\equiv0\pmod2,                          \tag{G031.4}
\]
where the last parity follows from $G_1\equiv0$.
Hence $48\mid g_1$. None of these quotient polynomials
is asserted to be effective.

\paragraph{The five possible $u$ pairs.}
Since $u_1,u_2$ are even and at least four, their product
less than $36$ leaves only
\[
 (4,4),\ (4,6),\ (6,4),\ (4,8),\ (8,4).
\]
Solving (G031.2) at $w=6$ gives
\[
 g_1=\frac{24(6d_1+u_1d_2)}{36-u_1u_2},\qquad
 g_2=\frac{24(u_2d_1+6d_2)}{36-u_1u_2}.             \tag{G031.5}
\]
In either pair containing eight, $g_1+g_2\ge576$,
so $20(g_1+g_2)>11128$ already.

For $(u_1,u_2)=(6,4)$, the actual identities with
$U_1=W,U_2=D$ yield
\[
 G_1=DW K,\qquad K=(\delta_1+\delta_2)/2.
\]
The polynomial $K$ is integral and effective, with
$K\ge\psi$ coefficientwise. Since $G_1\equiv0$ and
$DW\equiv D^2$ modulo two, $K$ has even coefficients.
Its coefficients at $q^{\pm1},q^{\pm2}$ are therefore
at least two. Hence $t=d_1+d_2=2K(4)\ge16$.
Formula (G031.5) gives $g_1=12t$ and
$g_2=8t+4d_2\ge8t+16$. Consequently
\[
 T\ge s+20(g_1+g_2)+1200
    \ge16t^2+432t+1520
    \ge12528>11128.
\]

For $(u_1,u_2)=(4,6)$, the identities instead give
$G_2=DW K$ with the same effective
$K=(\delta_1+\delta_2)/2\ge\psi$.
Now $G_2\equiv D^2(D-1)$ implies $K\equiv D-1$.
In Laurent coefficients its central coefficient is
odd and at least one, its coefficients at $q^{\pm1}$
are odd and at least one, and those at $q^{\pm2}$
are even and at least two. Therefore $K(4)\ge7$,
so $t\ge14$. The preceding symmetric scalar estimate,
including the remaining lower bounds in (G031.1), gives
\[
 T\ge16t^2+458t+2696
    \ge12244>11128.
\]

\paragraph{The remaining pair $(4,4)$.}
Here (G031.5) becomes
\[
 g_1=\frac{12}{5}(3d_1+2d_2),\qquad
 g_2=\frac{12}{5}(2d_1+3d_2),\qquad
 g_1+g_2=12t.
\]
Also
\[
 s=\frac{24}{25}(6t^2+d_1d_2)\ge\frac{144}{25}t^2.
\]
Using the remaining mass bounds in (G031.1),
\[
 T\ge\frac{144}{25}t^2+266t+2336.                  \tag{G031.6}
\]
At $t=24$ this exceeds $11128$, and it is increasing
for $t\ge0$. Thus $t<24$.

Write $d_1=2a,d_2=2b$. The integrality of $g_1,g_2$
gives $b=a+5k$ for some integer $k$. The stronger
conditions $48\mid g_1$ and $g_2/24$ odd show
that $a$ is even and $k$ is odd:
\[
 g_1=24a+48k,\qquad g_2=24a+72k.
\]
With $a,b\ge2$ and $d_1+d_2<24$, the only possibilities
are $(d_1,d_2)=(4,14)$ and $(16,6)$.
For completeness, $k\ge1$ forces $t=4a+10k\ge18$;
if $k\ge3$ it is at least $38$, while $k=1$ forces
$a=2$. If $k\le-1$, the inequality $b\ge2$ gives
$a\ge2-5k$; evenness improves this to $a\ge8$ when
$k=-1$, leaving $a=8$, and $k\le-3$ forces $t\ge42$.

For $(4,14)$ one has $(g_1,g_2)=(96,120)$ and
$\delta_1=\psi$. In (G031.4), write $C_1=2H$;
then $H$ is integral and $H(4)=g_1/48=2$.
Since both $U_i=D$, the first identity gives
\[
 \delta_1=W C_1-D C_2=2W H-D C_2.
\]
As $D\mid\delta_1$ and $W(0)=2$, this forces
$H(0)=0$. Thus $D\mid H$ over $\mathbb Z[D]$, and
$H(4)$ is a multiple of four, contradicting $H(4)=2$.

For $(16,6)$ one has $(g_1,g_2,s)=(144,120,2880)$.
The remaining identity gives
\[
 v_1=\frac{576-a_1}{20}.
\]
Since $v_1$ is even, $a_1\equiv16\pmod{40}$;
with $a_1\ge8$ it follows that $a_1\ge16$.
Hence $14a_1+20v_1=576+13a_1\ge784$.
The other terms of (G031.1) then imply
\[
 T\ge2880+5280+1200+160+572+784+544
       =11420>11128.
\]
This excludes the last branch and proves the statement.
\hfill$\square$

\paragraph{Geometric translation and scope.}
In the NN-G010 realization by graded sign-representation
kernels, $w$ is the ordinary dimension of the gluing
coefficient object. The theorem says that any such
closed-table realization costing at most $3352$
regular-sign pairs must use a gluing object of dimension
at least eight. A dimension-six gluing object has the
unique character $D+2$; its associative kernel products
would force the incompatible character congruences and
mass bounds proved above. This is an actual dimension
obstruction for the effective realization, not a statement
about Frobenius weights or point counts. The character
dictionary is SGA 3, I, 4.7.2--4.7.3; the composition
dictionary is archive cote 2, scan 17, as realized in
NN-G010. The result shrinks the remaining optimization
gate but gives no global optimum or canonical GCT basis.


\subsection{NN-G032: bounded-mass lift existence is fixed-parameter tractable}
\paragraph{Statement.}
For the symmetric effective associative lift problem of NN-G029,
let $N$ be a bound on the total coefficient mass, including
the unit, and let $L$ be the bit length of the sparse input
trace table, its rank and its specified unit and forbidden
positions, and the budget. There are a computable function
$F$ and an absolute constant $c$ such that lift existence
with mass at most $N$ is decidable in time
\[
                    F(N)(L+1)^c.                  \tag{G032.1}
\]
Thus the problem is fixed-parameter tractable in the numerical
mass budget. This is valid even when the budget is encoded
in binary; it is not a polynomial-time claim for unbounded
$N$. A lift can also be produced in the same type of bound.
The claim concerns the unrestricted symmetric table problem;
it does not identify that optimization with a canonical GCT
basis or imply any positive multiplicity formula.

The proof eliminates integer exponent variables symbolically.
The new consequence is the parameterized algorithm for this
lift problem; the integer-linear quantifier elimination used
below is a standard method, proved here self-containedly.
It never searches through all integers up to the largest input
trace exponent. Its constants can be extremely large, so it
does not make the $N=11147$ example a practical computation.

\paragraph{Finite discrete patterns and integer formulas.}
First canonically combine repeated sparse monomials, and
check the prescribed symmetry, unit and input format. These
operations take polynomial time in $L$. If the absolute trace
mass exceeds $N$, reject. A unital rank-$r$ lift contains at
least $2r-1$ unit coefficients, so reject also if $2r-1>N$.
Consequently every nonrejected input has $r\le(N+1)/2$,
at most $N$ nonzero trace monomials, and trace coefficient
absolute values at most $N$.

The following discrete information consequently has only
a number of possibilities bounded by a computable function
of $N$: the rank; the occupied trace positions and their
integer coefficients; the forbidden positions; and the
coefficient-and-position patterns of all possible effective
lift tables of total mass at most $N$.
For a lift pattern, every noncentral symmetric monomial pair
costs at least two units of mass. Assign a distinct formal
positive integer variable to each occurrence of such a pair,
and retain the central coefficients as fixed integers.
There are at most $\lfloor N/2\rfloor$ exponent variables.
Within each position, require that separate pairs have
distinct exponents. Coefficients are selected as part of
the finite pattern, so there are no unknown coefficient
products in what follows.

Let $k_1,\ldots,k_h$ be the distinct positive exponents of
the input trace. They are free integer parameters; their
binary magnitudes are not included in the discrete pattern.
One has $h\le\lfloor N/2\rfloor$.
As in NN-G029, form zero, the signed individual exponent
variables and trace parameters, and all sums of two signed
individual forms. Enumerate their possible total order
and equality patterns. Each chosen pattern is expressed
by a conjunction of linear equalities and strict
inequalities with integer coefficients bounded in terms
of $N$. Since the variables are integers, a strict
inequality is an ordinary weak inequality with an added
constant one.

For a fixed collision pattern, domination and parity are
comparisons of fixed coefficients, and every associativity
coefficient is a fixed integer sum of fixed integer products.
Discard a pattern if these tests fail, or if it violates a
unit or forbidden-position condition. Every remaining
pattern yields a formula
\[
 \exists x_1,\ldots,x_m\in\mathbb Z\quad
       \Phi(x_1,\ldots,x_m;k_1,\ldots,k_h),          \tag{G032.2}
\]
where $m\le N/2$ and $\Phi$ is a finite conjunction of
integer linear inequalities, including $x_i\ge1$.
The numbers of patterns, atoms and variables, and all
coefficients and constants in these formulas, are bounded
by functions of $N$ alone. The input trace exponents
appear only as free parameters $k_j$.

Conversely, a satisfying integer assignment to any retained
formula gives an actual lift: its monomials have the selected
positive exponents; its constants dominate the trace with
the correct parity; and all coefficient identities for
associativity hold by the selected collision partition.
Thus the finite disjunction of (G032.2) is equivalent to
the desired bounded-mass existence question. This step
does not require an incumbent or the numerical cutoff
from NN-G029.

\paragraph{Elementary elimination of one integer variable.}
We prove the needed elimination rule completely. Consider
any finite Boolean combination of integer linear inequalities
and congruences in one quantified integer $y$, other integer
variables, and free parameters. All coefficients and moduli
are fixed integers. Equalities are pairs of inequalities.
Negating an inequality changes it to an inequality with
constant shift one. A negated congruence modulo a fixed
positive $m$ is a disjunction over its other $m-1$ residues.
After distributing conjunctions and disjunctions, it is
enough to eliminate $y$ from one conjunction of inequalities
and congruences. This may greatly increase size, but the
increase depends only on the formula, not on parameter values.

Let $B\ge1$ be the least common multiple of the absolute
values of all nonzero coefficients of $y$ in that conjunction.
If none occurs, set $B=1$. Substitute $z=By$ and add
$z\equiv0\pmod B$.
An atom $ay+t\ge0$, where $t$ does not involve $y$,
becomes, for $a\ne0$,
\[
 \operatorname{sgn}(a)z+(B/|a|)t\ge0.
\]
A congruence $ay+t\equiv0\pmod m$ becomes
\[
 \operatorname{sgn}(a)z+(B/|a|)t
                   \equiv0\pmod{mB/|a|}.
\]
These transformations are equivalent under $z=By$:
the inequality is multiplied by a positive integer, and
$m\mid u$ is equivalent to $cm\mid cu$ for every
positive integer $c$. Atoms with $a=0$ are unchanged.
Conversely $z\equiv0\pmod B$ recovers an integer $y$.

All remaining atoms involving $z$ are therefore lower bounds
$z\ge\ell_i$, upper bounds $z\le u_j$, or congruences
$z\equiv r_a\pmod{m_a}$, where the bounds and residues
are integer linear forms in the other variables and free
parameters. Let $M$ be the least common multiple of the
fixed positive moduli. Enumerate $a=0,\ldots,M-1$ and
impose the compatibility atoms
\[
                      a\equiv r_b\pmod{m_b}
                 \quad\hbox{for every }b.
\]
These atoms exactly say that the chosen residue class
$z\equiv a\pmod M$ satisfies every congruence.

If either the lower-bound list or the upper-bound list is
empty, that residue class has an integer satisfying every
remaining one-sided bound. Thus no further condition is
needed in that case. If both lists are nonempty, enumerate
an index $i$ for a largest lower bound and a number
$b\in\{0,\ldots,M-1\}$, and impose
\[
 \ell_i\ge\ell_{i'}\ (\hbox{all }i'),\qquad
 \ell_i+b\equiv a\pmod M,\qquad
 \ell_i+b\le u_j\ (\hbox{all }j).
\]
These conditions are necessary and sufficient: the smallest
integer in the chosen residue class at least $\ell_i$ is
exactly $\ell_i+b$ for its unique such $b$.
The resulting finite disjunction contains no $z$ or $y$
and consists only of integer linear inequalities and
congruences. This proves elimination of one integer
existential quantifier, with an explicit finite procedure.

\paragraph{Runtime and witness recovery.}
Apply this elimination successively to all the exponent
variables in every retained formula (G032.2). Its initial
size, coefficients, moduli and number of variables depend
only on $N$. At every step the resulting size, coefficients
and moduli still depend only on $N$, even though they may
grow rapidly. The result is a finite Boolean combination
of integer linear inequalities and fixed-modulus
congruences in the input exponents $k_j$ alone.
Formula construction and elimination therefore cost some
computable function of $N$, independently of the binary
magnitudes of the $k_j$.

Substitute the given binary exponents into this formula.
There are only $F_0(N)$ atoms and coefficients, each of
bit length at most $F_0(N)$ for a computable $F_0$.
Integer addition, multiplication and division by these
fixed coefficients and moduli can all be carried out by
elementary bit algorithms in time bounded by
$F_1(N)(L+1)^2$. The finite initial parsing and sorting
contribute another fixed polynomial in $L$.
This proves (G032.1), with an absolute polynomial exponent.

One may also retain the selected disjunct at every elimination
step. Recover variables in reverse order, selecting a true
disjunct after substituting the free parameters and the variables
already recovered at that stage. With two-sided
bounds use the recorded integer $z=\ell_i+b$; with only
lower bounds choose the least compatible residue representative
above their maximum; with only upper bounds choose the
largest representative below their minimum; and with no
bounds choose the selected residue itself. Then set $y=z/B$.
The compatibility atoms ensure divisibility by $B$.
Repeating this construction yields the integer exponents
of a lift. The number of steps and every arithmetic
coefficient are bounded by functions of $N$, so recovery
obeys the same type of parameterized time and bit-length
bound. In particular every recovered bit length is
$O_N(L+1)$, since each recovery step uses a bounded number
of linear operations and fixed-modulus divisions with
coefficients depending only on $N$.\hfill$\square$

\paragraph{Geometric translation and limitations.}
For a fixed coefficient pattern, the grading weights lie
in the cocharacter lattice of a finitely presented
diagonalizable group, with the associativity collisions
imposing integral linear relations. Domination and the
chosen order of weights specify finitely many lattice
chambers. Eliminating exponent variables projects these
integer lattice conditions; the resulting congruences
record finite-index lattice obstructions. Their character
duals are finite diagonalizable groups such as $\mu_m$.
The actual conclusion is that, for bounded total dimension
mass, admissibility of the prescribed trace weights is
decidable by finitely many such linear inequalities and
congruences, computed uniformly from the mass budget.

This is an arithmetic statement about cocharacter lattices,
not a claim that inequalities define Zariski-open subsets
or that an effective lift is a canonical geometric object.
The grading dictionary is SGA 3, I, 4.7.2--4.7.3, and
the actual composing kernel realization is NN-G010,
using archive cote 2, scan 17. The algorithm gives no
polynomial bound on the mass budget needed by varying
GCT inputs, no practical small value of $F$, and no
rank-four positivity, purity, PH0 or PH1 theorem.


\subsection{NN-G033: an effective realization with 2744 cancelling pairs}
\paragraph{Statement.}
The same rank-three trace table has an effective associative
symmetric closed-table realization with exactly $2744$
cancelling pairs, including $1296$ in the bottom cell.
This improves the unrestricted symmetric realization upper
bound from $3352$. All actual structure coefficients are
integral; the second auxiliary $\Lambda$ remains rational.
No global optimum is asserted in this construction theorem.

\paragraph{Construction and proof of associativity.}
As throughout, $D=q+2+q^{-1}$. Put
\[
 C=D-1,\qquad H=2D(D-2),\qquad
 W=2D^2-7D+4=D+2(q^2+q^{-2}).
\]
Define the actual base coefficients by
\[
 \begin{aligned}
 \delta_1&=D(D-3),&\delta_2&=WC-H,\\
 G_1&=DH,&G_2&=DWC,&S&=D^2CH,\\
 U_1&=D,&U_2&=W,&V_1&=2D^2-3D,&V_2&=3W,\\
 A_1&=D^2(D-3),&A_2&=D^2WC-3DH,&
 \Lambda_i&=A_i/\delta_i.
 \end{aligned}                                                       \tag{G033.1}
\]
The symbol $\Lambda_i$ is only a rational coordinate
factor; the actual cell products use $\delta_i$ or
$A_i$, multiplied by nonnegative powers of $D$.
In particular $\Lambda_1=D$.

Here $H-\delta_1=DC$. Directly,
\[
 \frac{G_1G_2}{S}=W,\qquad
 G_1-D\delta_1=D^2C,\qquad
 G_2-D\delta_2=DH=G_1.
\]
It follows that
\[
 \frac{G_1(G_1-D\delta_1)}S=D=U_1,
 \qquad
 \frac{G_2(G_2-D\delta_2)}S=W=U_2.
\]
The remaining identities are
\[
 DG_1-A_1=DCV_1,\qquad DG_2-A_2=3DH,
\]
and therefore $G_i(DG_i-A_i)/S=V_i$ for both $i$.
Neither $\Lambda_i$ nor $D^2-\Lambda_i$ is the zero
rational function. For the second index this follows
explicitly from
\[
 D^2\delta_2-A_2=DH(3-D)\ne0.
\]
The injective generic block realization proved in
NN-G011 thus proves associativity over $\mathbb Q(D)$.
All actual coefficients are integral, so these identities
hold in $\mathbb Z[D]$. No regularity of the auxiliary
rational coordinate at every specialization is used.

\paragraph{Every coefficient condition and the exact score.}
A list $(a_0,\ldots,a_m)$ below denotes
$a_0+\sum_{j=1}^m a_j(q^j+q^{-j})$.
The base coefficient checks are completely recorded by
\[
\begin{array}{c|l|l}
 &\text{dimension coefficient}&\text{absolute trace target}\\\hline
 W&(2,1,2)&(2,1)\\
 U_1&(2,1)&(2,1)\\
 U_2&(2,1,2)&(2,1)\\
 \delta_1&(0,1,1)&(0,1,1)\\
 \delta_2&(0,1,1,2)&(0,1,1)\\
 G_1&(16,14,8,2)&(16,12,4)\\
 G_2&(18,17,13,7,2)&(14,11,5,1)\\
 S&(164,144,96,46,14,2)&(144,120,68,24,4)\\
 V_1&(6,5,2)&(6,5)\\
 V_2&(6,3,6)&(6,3,2)\\
 A_1&(2,3,3,1)&(2,1,1,1)\\
 A_2&(22,23,26,23,11,2)&(2,3,2,1,1)
\end{array}                                                         \tag{G033.2}
\]
Each left entry dominates the right entry coefficientwise
and has the same parity. Multiplication by the required
nonnegative powers of $D$ preserves parity and domination
of the absolute trace, by the coefficientwise triangle
inequality. This proves the conditions for the full
closed multiplication table. The included independent
integer-Laurent verifier checks all $132$ positions and
all $1000$ associativity triples with zero failures.

At $D=4$, the masses are
\[
 (w,u_1,u_2,v_1,v_2)=(8,4,8,20,24),\qquad
 (g_1,g_2,s)=(64,96,768),
\]
and $(d_1,d_2,a_1,a_2)=(4,8,16,192)$.
Hence the bottom mass is
\[
 768+20(64+96)+200\cdot8+20(4+8+20+24)=6688.
\]
The cell mass is $26(4+8)+14(16+192)=3224$.
The nonunit total is $9912$, giving
\[
 \frac{9912-4424}{2}=2744,\qquad
 \frac{6688-4096}{2}=1296.                         \tag{G033.3}
\]
Thus there are $1448$ pairs outside the bottom cell.
For comparison, the $3352$ realization had $2752$
bottom pairs and $600$ elsewhere. The new construction
reduces the total while moving more of the cancellation
into the two nonbottom cells; the objective is the total
over the specified presentation.

The ratios
\[
 \frac{G_1}{W\delta_1}
   =\frac{2D(D-2)}{W(D-3)},\qquad
 \frac{G_2}{W\delta_2}=\frac{D(D-1)}{WC-H}
\]
are nonconstant rational functions. Thus this realization
continues to lie outside the constant-ratio class of
NN-G026 and does not contradict its exact $4520$ optimum.
\hfill$\square$

\paragraph{A smaller complete search cutoff.}
Including the $19$ unit coefficients, the new total mass
is $N=9931=4443+2\cdot2744$. There are still five positive
trace anchors $1,\ldots,5$, so NN-G029 gives
\[
 n\le\lfloor9931/2\rfloor+5=4970,\qquad
 \deg d\le29820\cdot4^{4969}<2^{9953}.              \tag{G033.4}
\]
Accordingly, any equal-or-better unrestricted symmetric
realization has an equal-score representative below
this new cutoff, allowing extra zero-trace positions.
This improves the earlier $2^{11170}$ existence cutoff;
it does not bound the degree of every existing realization
or make the finite search practical.

\paragraph{Geometric translation and limitations.}
The effective Artin--Tate coefficient objects in
(G033.2), with the NN-G010 kernel-composition realization,
give an actual additive monoidal model of the same
graded sign trace. The excess of ordinary dimension
over absolute signed trace is precisely twice $2744$
regular-sign pairs in its chosen multiplication
presentation, with $1296$ pairs in the bottom cell.
The gluing character has dimension eight and includes
weight-two characters absent from $D$.

This construction concerns actual coefficient objects,
not formal $K_0$ cancellation or a count obtained by
evaluating a virtual class. Its class is governed by
SGA 3, I, 4.7.2--4.7.3 and the archive cote 2, scan 17
composition dictionary implemented in NN-G010.
The rational auxiliary $\Lambda_2$ does not represent
a fractional-dimensional object. The global minimum,
canonicality, rank-four extensions and PH0/PH1 remain
separate questions. The older exact subclass bounds
and the NN-G031 implication for all costs at most
$3352$ remain valid with their stated hypotheses.


\subsection{NN-G034: the exact optimum with one fixed first-cell realization is 2744}
\paragraph{Statement and class.}
In the symmetric closed-table presentation, impose only
\[
 \delta_1=D(D-3),\qquad G_1=2D^2(D-2),\qquad U_1=D.   \tag{G034.1}
\]
Allow all remaining actual coefficient polynomials to vary,
with arbitrary degrees and contents, and allow rational
auxiliary $\Lambda_i$ and nonconstant ratios.
Among effective associative lifts in this class, the exact
minimum is $2744$ cancelling pairs. There is exactly one
minimizing set of actual base polynomials, namely NN-G033.
This is an exact class theorem, not the unrestricted
$1000$-position optimum and not the optimum over all closed
tables with their first-cell data free.

\paragraph{Forced parameters and the first-cell optimum.}
Put $C=D-1$ and $H=2D(D-2)$, so $G_1=DH$.
The identities $U_1G_2=W(G_1-D\delta_1)$ and
$WS=G_1G_2$ force
\[
                  G_2=DWC,\qquad S=D^2CH.            \tag{G034.2}
\]
Thus $(g_1,g_2,s)=(64,12w,768)$ at $D=4$.
The first $V$ identity forces
\[
 A_1=D^2H-DCV_1,\qquad a_1=256-12v_1.
\]
Since $v_1$ is even and $a_1\ge8$, one has $v_1\le20$.
The first-cell contribution satisfies
\[
 20v_1+14a_1=3584-148v_1\ge624,
\]
with equality exactly when $v_1=20,a_1=16$.

There is exactly one allowable $V_1$ at this equality.
It dominates the trace character $5D-4=(6,5)$,
has the same parity, and has surplus mass four.
Since $A_1$ is effective and its first term $D^2H$
has degree four, $V_1$ has degree at most two.
The only possibilities are
\[
 5D,\qquad7D-8,\qquad5D-4+2(q^2+q^{-2})=2D^2-3D.
\]
The respective central coefficients of $A_1$ are
$-10,-6,2$. The last choice gives
$A_1=D^2(D-3)=(2,3,3,1)$ and satisfies all conditions;
the first two fail effectivity.

\paragraph{A complete mass reduction.}
For a competitor with score at most $2744$, the nonunit
mass is $T\le9912$. Formula (G031.1) and the preceding
first-cell bound give
\[
 \begin{aligned}
 T&\ge2856+440w+20u_2+20v_2+26d_2+14a_2,\\
 u_2&=\frac{w(3w-d_2)}{16},\qquad
 v_2=\frac{w(48w-a_2)}{64}.                         \end{aligned} \tag{G034.3}
\]
The omitted amount is exactly $148(20-v_1)\ge0$.
Using $u_2\ge4,v_2\ge16,d_2\ge4,a_2\ge16$ gives
$T\ge3584+440w$, hence $w\le14$.
NN-G031 rules out $w=4,6$ at this score.
Eliminating $u_2,v_2$ in (G034.3) gives
\[
 T\ge2856+440w+\frac{75}{4}w^2
       +\left(26-\frac{5w}{4}\right)d_2
       +\left(14-\frac{5w}{16}\right)a_2.           \tag{G034.4}
\]
For $w=12,14$ both final coefficients are positive,
and even the first three terms are at least
$10836>9912$. Thus only $w=8$ or $w=10$ remain.

Write $T_j=q^j+q^{-j}$, an integral polynomial in $D$.
Every effective symmetric $W\equiv D\pmod2$
dominating $D$ has the following exhaustive small-mass forms:
\[
 \begin{array}{c|l}
 w=8&D+4\quad\hbox{or}\quad D+2T_j\quad(j\ge1),\\
 w=10&D+6\quad\hbox{or}\quad D+2+2T_j\quad(j\ge1).
 \end{array}                                                       \tag{G034.5}
\]
Indeed the even surplus has mass four or six; a noncentral
surplus pair costs four, leaving only the displayed central
remainder. Each polynomial in (G034.5) is primitive in
$\mathbb Z[D]$, since its Laurent coefficients have gcd one.

\paragraph{An equal-mass divisibility lemma.}
Suppose $W,U$ both belong to the same row of (G034.5),
and $U=WR$ for $R\in\mathbb Z[D]$ with $R\equiv1\pmod2$.
Then $U=W$.
If $\deg W\ge2$, its leading coefficient is two.
A nonconstant quotient would require $\deg U\ge2$,
also with leading coefficient two; thus $R$ would have
leading coefficient one, contradicting $R\equiv1$
in positive degree. A constant quotient has $R(4)=1$
and is one. For the degree-one choice with leading
coefficient three, a nonconstant integral quotient would
have to have leading coefficient $2/3$, impossible.

It remains to consider $W=D+\ell$ with $\ell=4$ or six.
A nonconstant quotient requires
$U=D+\ell-4+2T_j$ with $j\ge2$.
At $D=-\ell$, divisibility would give $T_j(-\ell)=2$.
This is impossible. To see it directly, set $x=\ell+2\ge6$;
the positive sequence
$a_j=(-1)^jT_j(-\ell)$ has $a_0=2,a_1=x$ and
$a_{j+1}=xa_j-a_{j-1}$. Induction gives
$a_{j+1}>a_j>2$ for $j\ge1$.
This completes the lemma without bounding $j$.

\paragraph{The complete $w=8$ branch.}
Each $W$ in its row of (G034.5) is coprime over
$\mathbb Q[D]$ to $D(D-2)$: at $D=0$ its value is
$4$ or $\pm4$, and at $D=2$ its value is among
$6,2,-2$. The two second-cell identities therefore
force
\[
 U_2=WR,\qquad V_2=WZ\qquad(R,Z\in\mathbb Z[D]).
\]
Here polynomial divisibility follows first over
$\mathbb Q[D]$ and then integrally by primitivity.
Cancelling $D$ modulo two gives $R\equiv Z\equiv1$.
Formula (G034.3) shows $4\le u_2\le10$, so the
positive odd integer $R(4)=u_2/8$ must be one.
The equal-mass lemma gives $U_2=W$, and then
\[
                 \delta_2=WC-H,\qquad d_2=8.         \tag{G034.6}
\]
Likewise $16\le v_2\le46$, so $Z(4)$ is three or five.
The corresponding $a_2$ are $192$ and $64$.

For $W=D+4$, (G034.6) has negative leading coefficient,
so it is not effective. For $W=D+2T_j$, one has
\[
             \delta_2=-D(D-3)+2(D-1)T_j.
\]
Domination of $D(D-3)$ forces $j\le2$: when $j\ge3$
the product $(D-1)T_j$ has no coefficient at $q^1$,
whereas $D(D-3)$ does. Thus only $W=3D-4$ and
$W=D+2T_2$ remain.

The second $V$ identity is now
\[
              A_2=D^2\bigl(WC-2(D-2)Z\bigr).         \tag{G034.7}
\]
If $W=3D-4$, effectivity of $V_2=WZ$ makes the leading
coefficient of a nonconstant $Z$ positive. Its parity
makes that coefficient even. For $\deg Z=1$,
the leading coefficient inside (G034.7) is at most
$3-4<0$, and higher degrees also give a negative
leading coefficient. Thus $Z$ must be constant.
But then $V_2$ has degree one and cannot dominate
the trace coefficient two at $q^2$. This excludes
the degree-one choice.

For $W=D+2T_2=(2,1,2)$, the same leading-coefficient
argument gives $\deg Z\le1$: degree two would have
leading coefficient at least two, against the leading
coefficient two of $WC$.
Write $Z=(z-2a)+a(q+q^{-1})$, where $z=Z(4)\in\{3,5\}$.
Here $a\ge0$ is even, while the central coefficient
is allowed to be negative. The central coefficient
of $V_2=WZ$ is $2z-2a$ and must be at least six.
If $z=3$, this forces $a=0$ and $Z=3$.
If $z=5$, it permits only $a=0,2$.
The central coefficient of $A_2$ is
$70-16z+4a$, which is negative in both $z=5$ cases.
Consequently only $Z=3$ survives. Its $a_2=192$
gives, from (G034.4),
\[
                  T\ge7576+16\cdot8+\frac{23}{2}\cdot192
                    =9912.                         \tag{G034.8}
\]
Equality forces the unique first-cell optimum already
proved, and gives exactly the parameters of NN-G033.

\paragraph{The complete $w=10$ branch: scalar constraints.}
Here (G034.3) reads
\[
 u_2=\frac58(30-d_2),\qquad
 v_2=\frac5{32}(480-a_2).
\]
Evenness and positivity force $d_2\equiv14\pmod{16}$
and $d_2<30$, hence $d_2=14,u_2=10$.
They also force $a_2\equiv32\pmod{64}$.
The choice $a_2=32$ gives the lower bound $T\ge9668$,
whereas the next choice $a_2=96$ already gives
$T\ge10364>9912$. Thus a competitor must have
\[
                  d_2=14,\quad u_2=10,\quad
                  a_2=32,\quad v_2=70.               \tag{G034.9}
\]

\paragraph{Coprime polynomials in the $w=10$ branch.}
No $W$ in this row is divisible by $D$.
It is divisible by $D-2$ exactly when
$W=D+2+2T_j$ with $j\equiv2\pmod4$, since
$T_j(2)$ is zero for odd $j$ and
$2(-1)^{j/2}$ for even $j$.
First suppose $W$ is coprime to $D(D-2)$.
Then $U_2=WR,V_2=WZ$ integrally, with
$R\equiv Z\equiv1$, and $R(4)=1$.
The equal-mass lemma gives $U_2=W$ and
$\delta_2=WC-H$.

The choice $W=D+6$ again gives a negative leading
coefficient in $\delta_2$. For
$W=D+2+2T_j$, domination of $\delta_2$ requires
\[
                       (D-1)(1+T_j)\ge D(D-3).
\]
The coefficient at $q^2$ forces $j\le3$.
The coprime possibilities are $j=1$ and $j=3$.
For $j=1$, $W=3D-2$ and the same leading-coefficient
argument in (G034.7) forces $Z$ constant, so $V_2$
cannot dominate its required $q^2$ coefficient.

For $j=3$, effectivity in (G034.7) forces
$\deg Z\le2$: a degree-three leading coefficient
would be positive even and produce a negative leading
term, and larger degrees are also impossible.
Write $Z=(z_0,z_1,z_2)$ in symmetric Laurent notation.
By (G034.9), $Z(4)=7$.
The coefficient of $q^2$ in $D^2WC$ is $54$.
Since $G_1=(16,14,8,2)$, the coefficient of $q^2$
in $G_1Z$ is
\[
        8z_0+16z_1+16z_2=8Z(4)=56.
\]
Thus $A_2=D^2WC-G_1Z$ has coefficient $-2$ at $q^2$,
contradicting effectivity. This excludes every coprime
$w=10$ choice, including all unbounded locations $j$.

\paragraph{The $D-2$ factor in the $w=10$ branch.}
It remains to consider $W=D+2+2T_j$ with
$j\equiv2\pmod4$. Write $W=(D-2)B$.
The polynomial $B$ is primitive, $B(0)=-3$, and
$B(2)=1$: for even $j$, $T_j$ is an even polynomial
in $D-2$, so $W'(2)=1$.
Thus $B$ is coprime to $D(D-2)$; modulo two,
$B\equiv1$. The identities now force
\[
 U_2=BR,\quad V_2=BZ,\qquad
 R\equiv Z\equiv D\pmod2,
 \qquad R(4)=2,\quad Z(4)=14.                       \tag{G034.10}
\]
The leading coefficient of $B$ is two and
$\deg B=j-1$. If $R$ were constant, its parity
would contradict $R\equiv D$.
Otherwise $\deg U_2\ge j\ge2$, so the mass-ten
classification gives leading coefficient two for $U_2$.
Hence $R$ has leading coefficient one, and its parity
forces $\deg R=1$. The only mass-ten polynomial of
degree $j$ in (G034.5) is $W$ itself. Therefore
$U_2=W$, $R=D-2$ and $\delta_2=WC-H$.
The preceding $q^2$ domination argument gives $j\le3$;
thus necessarily $j=2$.

We have reached
\[
 W=2D^2-7D+6=(D-2)(2D-3),\qquad B=2D-3.
\]
The second $V$ identity reduces to
\[
                    A_2=D^2(WC-2Z).                 \tag{G034.11}
\]
Effectivity forces $\deg Z\le2$. Indeed its leading
coefficient is positive when nonconstant; degree three
would make it positive even by (G034.10), exceeding
the leading coefficient two of $WC$ after the factor
two in (G034.11), and higher degrees also fail.
Write $Z=(z_0,z_1,z_2)$. Its parity says that $z_2$
is even, $z_1$ is odd, and $z_0$ is even; its mass says
$z_0=14-2z_1-2z_2$.
The relevant coefficients of (G034.11) are
\[
 \begin{array}{c|l}
 q^1&-25+2z_1+8z_2,\\
 q^2&32-4z_1-8z_2,\\
 q^4&11-2z_2.
 \end{array}
\]
Domination of their trace targets $3,2,1$ implies
\[
 z_1+4z_2\ge14,\qquad z_1+2z_2\le7,\qquad z_2\le5.
\]
The first two inequalities force $z_2\ge7/2$.
Evenness then gives $z_2=4$, and oddness gives
$z_1=-1$, hence $z_0=8$.
But the central coefficient of $V_2=BZ$ is then
$z_0+4z_1=4$, below its required value six.
This excludes the final $w=10$ branch.

The $w=8$ equality realization of NN-G033 is therefore
the unique minimizer in (G034.1), proving the theorem.
\hfill$\square$

\paragraph{Geometric translation and remaining gate.}
Fixing (G034.1) fixes the first-cell coefficient object,
its bottom action character, and one gluing product in
the NN-G010 category of graded sign-representation kernels.
Among all effective associative completions in this closed
presentation, the least total excess of dimension over
absolute trace is exactly $2\cdot2744$, achieved by the
single set of coefficient characters of NN-G033.
The proof allows arbitrary grading locations in the other
objects; the mass-ten $D-2$ factor is treated by an actual
integral divisibility argument, not discarded as generic.

The finite-character evaluations $D=0$ and $D=2$
are restrictions at $q=-1$ and $q=i$, under SGA 3,
I, 4.7.2--4.7.3. They are not assertions about Frobenius
eigenvalues. The monoidal composition dictionary is
archive cote 2, scan 17, via NN-G010.
Any smaller closed-table score must change at least one
of the three first-cell data in (G034.1); other table
shapes remain outside this class. This is a strict
remaining gate, with no canonical-basis, rank-four,
PH0/PH1 or purity conclusion.


\subsection{NN-G035: a complete first-cell mass gate for competitive closed tables}
\paragraph{Statement.}
Every symmetric closed-table effective associative lift with
at most $2744$ cancelling pairs lies in the complete necessary
mass relaxation below. It has $w\in\{8,10,12\}$ and $s\ge588$.
Its first-cell mass triple $(g_1,d_1,u_1)$ belongs to the
following set of $21$ triples, grouped by $g_1$:
\[
\begin{array}{c|l}
 g_1&\text{possible }(d_1,u_1)\\\hline
 64&(4,4)\\
 72&(4,6),(6,4)\\
 76&(4,6)\\
 80&(4,8),(8,4),(8,6),(12,4)\\
 84&(6,6),(10,6),(14,4)\\
 88&(8,8)\\
 96&(4,10),(4,12),(6,6),(12,4),(14,6)\\
 100&(4,10),(6,10),(10,6)\\
 128&(12,8)
\end{array}                                                         \tag{G035.1}
\]
The full relaxation has $43$ base rows and $467$ possible
auxiliary increment pairs, all explicitly listed in the
certificate. Its minimum relaxed score is $1700$.
Consequently the minimum over all symmetric closed tables
is between $1700$ and $2744$.
Surviving mass tuples are not asserted realizable; no
Laurent degree or support-location bound is inferred.
This theorem does not apply to the broader unrestricted
$1000$-position problem of NN-G029.

\paragraph{Complete finite bounds.}
Use the actual-coefficient mass notation of (G031.1).
The known lower bounds are
\[
 s\ge576,\quad g_i\ge48,\quad d_i\ge4,\quad
 u_i\ge4,\quad v_i\ge16,\quad a_1\ge8,\quad a_2\ge16.
\]
All are even; NN-G030 also makes $g_1,g_2,s,a_2$
multiples of four. The score cap is $T\le9912$.
NN-G031 gives $w\ge8$, and the mass formula implies
\[
              20(g_1+g_2)\le9912-3520,
              \qquad g_1+g_2\le316.                 \tag{G035.2}
\]
Thus each $g_i\le268$.
If $w\ge20$, the relation $g_1g_2=ws\ge576w$
and the arithmetic-geometric mean inequality give
$g_1+g_2\ge48\sqrt{20}>214$.
Since the sum is divisible by four, it is at least
$216$. The mass formula would then imply
\[
 T\ge576+20\cdot216+200\cdot20+800+208+336
       =10240>9912.
\]
Hence $w$ is one of $8,10,12,14,16,18$.

Fix one such $w$ and multiples of four $g_i\ge48$
with sum at most $316$. Then $s=g_1g_2/w$ must be
an integer multiple of four at least $576$.
For each index $i$, with $j\ne i$, enumerate the
even integers $4\le d_i\le\lfloor g_i/4\rfloor$,
retaining exactly those for which
\[
             u_i=\frac{w(g_i-4d_i)}{g_j}
\]
is an even integer at least four. These finite bounds
contain every possible actual cell mass.

The auxiliary masses satisfy
$v_i=w(4g_i-a_i)/g_j$. Set $t_1=2,t_2=4$, and
\[
 M_i=\operatorname{lcm}\left(t_i,
                   \frac{2g_j}{\gcd(w,2g_j)}\right).
\]
Their necessary congruence is exactly
$a_i\equiv4g_i\pmod{M_i}$, together with
$a_1\ge8,a_2\ge16$.
Let $a_i^0$ be the least integer satisfying this lower
bound and congruence, and put
$v_i^0=w(4g_i-a_i^0)/g_j$. If $v_i^0<16$, discard
the tuple, since increasing $a_i$ only decreases $v_i$.
Otherwise every actual auxiliary mass is
\[
 a_i=a_i^0+M_i k_i,\qquad
 v_i=v_i^0-\frac{wM_i}{g_j}k_i,\qquad k_i\ge0.       \tag{G035.3}
\]
The mass increment per unit $k_i$ is
\[
 \Delta_i=14M_i-20\frac{wM_i}{g_j}>0,
\]
because $w\le18$ and $g_j\ge48$ make
$14-20w/g_j\ge13/2>0$.
Thus the least auxiliary masses give a rigorous lower
bound $T_0$ for the tuple, and all remaining choices
satisfy the explicit finite inequalities
\[
 \Delta_1k_1+\Delta_2k_2\le9912-T_0,\qquad
 \frac{wM_i}{g_j}k_i\le v_i^0-16.                   \tag{G035.4}
\]

\paragraph{The extra $U_i=D$ tests for masses eight and ten.}
For $w=8,10$, (G034.5) shows that $W$ is primitive
and coprime to $D$.
If $u_1=4$, then $U_1=D$. The proof of (G031.4)
gives $G_2=DWC$ for an integral $C\equiv D-1\pmod2$;
therefore $g_2/(4w)$ is an odd integer.
If $u_2=4$, the analogous quotient is even, so
$8w\mid g_1$. These tests use only the stated mass
of $W$ and integral associativity, not a constant-ratio
or a degree assumption.

\paragraph{Why the minimal bottom polynomial is impossible.}
Suppose $s=576$. Coefficientwise domination and equality
of masses force
\[
                  S=4D^2(D-1)^2.                    \tag{G035.5}
\]
Write $C=D-1$ temporarily. From $WS=G_1G_2$ and
$G_1\equiv0,G_2\equiv D^2C\pmod2$, reduction after
dividing by two shows that $G_1$ is divisible by four.
Writing $G_1=4J$ gives
\[
             JG_2=WD^2C^2,\qquad J\equiv DC\pmod2.
\]
Since $W\equiv D\pmod2$, $W(1)$ is odd, so $C$
does not divide $W$. Neither $J$ nor $G_2$ can be
divisible by $C^2$, by their displayed reductions
modulo two. Their product contains exactly $C^2$;
therefore each contains exactly one factor $C$.
Write
\[
          J=CA,\qquad G_2=CB,\qquad
          AB=WD^2,\qquad A\equiv D, B\equiv D^2\pmod2.
\]
The actual $D$-adic orders of $A,B,W$ are at most
$1,2,1$, respectively, again by these nonzero reductions.
If $D\nmid W$, their first two orders sum to two;
if $D\mid W$, they sum to three. In either case
at least one is maximal: either $A=DA_1$ with
$A_1\equiv1\pmod2$, or $B=D^2B_1$ with
$B_1\equiv1\pmod2$.
Thus at least one of $g_1,g_2$ is $48$ times a
positive odd integer.

It cannot equal $48$: using $g_1g_2=576w$ in
the corresponding $u_i$ identity gives
$u_i=4-d_i/3<4$.
Hence some $g_i\ge144$.
For that index, eliminate $v_i$ using $s=576$:
\[
 20v_i+14a_i=\frac{5g_i^2}{36}
                   +\left(14-\frac{5g_i}{144}\right)a_i.
\]
The second coefficient is positive because $g_i\le268$.
Since $a_i\ge8$, the right side is at least
$5g_i^2/36+112-5g_i/18$, an increasing function for
$g_i\ge144$, and hence at least $2952$.
For the other index $20v_j+14a_j\ge432$.
Also $g_1+g_2\ge192$ and $w=g_1g_2/576\ge12$.
Together these estimates imply
\[
 T\ge576+3840+2400+160+208+2952+432
       =10568>9912.
\]
This contradiction excludes every minimal-bottom case;
no finite computation was used for this polynomial step.

\paragraph{Exact enumeration and completeness of the certificate.}
The verifier enumerates precisely the finite ranges in
(G035.2)--(G035.4), using integer gcd and lcm arithmetic.
It first tests $s$, the lower mass bound, and the two
$U_i=D$ divisibility tests, then enumerates every $d_i$
and uses the least admissible $a_i^0$.
There are $1445$ coarse $(w,g_1,g_2)$ tuples before
the cell tests and $55$ surviving base rows with
$T_0\le9912$. Exactly $12$ have $s=576$ and are
excluded by the preceding proof. The remaining $43$
rows have exactly $467$ increment pairs (G035.4).
The certificate lists every row and every pair.
It explicitly includes the actual NN-G033 mass tuple.

Projecting these rows gives (G035.1), $w\in\{8,10,12\}$,
and least $s=588$. Their least $T_0$ is $7824$, giving
the relaxed lower bound $(7824-4424)/2=1700$.
Every actual realization at or below the incumbent
maps to one of these rows and increments by the
proved equations and bounds. Larger scores are already
above $1700$. This proves the stated closed-table
lower bound and all finite necessary restrictions.
\hfill$\square$

\paragraph{Geometric translation and scope.}
For the actual graded sign-representation kernels of
NN-G010, this theorem bounds the possible ordinary
dimensions of the first-cell coefficient object, its
bottom action object, and its correction product by
the $21$ listed triples. Their grading locations are
still unrestricted. The bottom coefficient object
cannot have exactly the absolute trace character
$4D^2(D-1)^2$ in any competitive realization; integral
character factorization and parity force a dimension
contradiction. The finite list records necessary
dimension data, not actual objects for every row.

The dictionary is SGA 3, I, 4.7.2--4.7.3 and archive
cote 2, scan 17 via NN-G010. No equality of formal
$K_0$ classes is treated as an effective realization.
The result shrinks the closed-table search to explicit
mass data, while NN-G029--NN-G032 separately address
the wider unrestricted presentation. Neither statement
supplies a canonical GCT basis, rank-four positivity,
or PH0/PH1.


\subsection{NN-G036: the exact unique 2744-pair optimum for the symmetric closed table}
\paragraph{Statement.}
The minimum cancelling-pair score in the specified symmetric
closed-table presentation is exactly
\[
                                2744.                \tag{G036.1}
\]
The actual polynomial parameter realization attaining it is
unique within that presentation and is NN-G033.
All coefficient degrees and contents, all grading support
locations, and all rational ratios and auxiliary parameters
are allowed, subject to integrality and effectivity of the
actual structure constants. This is not a numerical lower
bound for unrestricted extra-position multiplication tables.

The proof retains its checked stages in their original order.
The first stage forces every competitive $g_1=64$ table into
NN-G034, which has the unique $2744$ attainer. The remaining
ordinary-dimension rows are then eliminated by uniform
integral polynomial arguments. Intermediate intervals and
row counts below describe successive stages, not the final
gate. The last proof and complete projection leave no other
competitive row.

\paragraph{Uniform irreducibility of the small gluing characters.}
Retain $T_j=q^j+q^{-j}$ as an integral polynomial in $D$.
Every mass-eight character in (G034.5) is irreducible
over $\mathbb Q[D]$.
Every mass-ten character there is irreducible except
\[
 W=D+2+2T_j\quad(j\equiv2\pmod4),
\]
in which case
\[
               W=(D-2)B,\qquad B\text{ primitive and irreducible},
               \qquad B\equiv1\pmod2,\quad B(4)=5.   \tag{G036.2}
\]
These statements include all support indices $j$.

Here is an elementary proof. The degree-one choices
are immediate. Every remaining $W$ is primitive,
has leading coefficient two, and reduces to $D$
modulo two. In any nontrivial primitive integral
factorization, one factor has leading coefficient
$\pm1$ and the other $\pm2$. The first factor's
degree is unchanged upon reduction modulo two;
since the product of the reductions is $D$, it
must have degree one. Up to a sign it is $D-r$,
with $r$ even. The other factor reduces to one,
so its constant term is odd.

For mass eight, $W(0)=\pm4$, hence $r=\pm4$.
But $W(4)=8$, and $W(-4)=-4+2T_j(-4)\ne0$.
The last assertion follows from the recurrence
used in NN-G034: at $D=-4$, the absolute values
of $T_j$ are at least six and strictly increase.
There is therefore no such factor.

For mass ten, $W(0)$ is $6$ or $-2$, hence
$r\in\{2,-2,6,-6\}$.
At $D=-2$ the value of $W$ is $2T_j(-2)\ne0$;
at $D=6$ it is positive; and at $D=-6$ the
same recurrence, with initial absolute value eight,
makes $-4+2T_j(-6)$ nonzero.
At $D=2$ the value is $4+2T_j(2)$, which is
zero exactly for $j\equiv2\pmod4$.
This proves that $D-2$ is the only possible factor.
After division by it, $B$ is primitive with leading
coefficient two and reduction one modulo two.
If $B$ factored nontrivially, a factor with odd
leading coefficient would retain positive degree
modulo two, impossible when the product reduces
to one. Thus $B$ is irreducible.
Finally $B(4)=W(4)/2=5$.

\paragraph{The consequence of $g_1=64$.}
The complete mass list NN-G035 gives
$d_1=4,u_1=4$ and $w=8$ or ten in this branch.
Therefore
\[
                 \delta_1=\psi:=D(D-3),\qquad U_1=D.
\]
The primitive $W$ is coprime to $D$ in both small-mass
families. The integral divisibility argument of
NN-G031 gives polynomials $C,Q\in\mathbb Z[D]$ with
\[
 \begin{aligned}
 G_2&=DWC,&G_1&=D\delta_1+D^2C=D^2Q,\\
 C&\equiv D-1\pmod2,&Q&=C+D-3\equiv0\pmod2,\\
 C(4)&=3,&Q(4)&=4.
 \end{aligned}                                                       \tag{G036.3}
\]
No effectivity of $C$ or $Q$ is asserted here.

If $W$ is irreducible, it cannot divide $Q$:
primitivity would make the quotient integral,
so $W(4)=8$ or ten would divide $Q(4)=4$.
Thus $W$ is coprime to $G_1$ and the second-cell
identity implies $U_2=WR$ integrally, with
$R\equiv1\pmod2$.
For $w=8$, the scalar relation
$u_2=12-d_2/2$ gives $4\le u_2\le10$.
Its positive odd ratio $u_2/8$ is therefore one.
The equal-mass lemma of NN-G034 implies $U_2=W$
and then $d_2=8$.
For $w=10$, evenness and positivity of
$u_2=5(30-d_2)/8$ give $d_2=14,u_2=10$.
The same equal-mass lemma gives $U_2=W$.

In the exceptional mass-ten case (G036.2), the
irreducible $B$ cannot divide $Q$, since $B(4)=5$
does not divide four. It is also coprime to $D$.
Thus $U_2=BR$ integrally; modulo two, $R\equiv D$.
The scalar relation still gives $u_2=10$.
Since $\deg B=j-1$ and its leading coefficient is
two, the same leading-coefficient argument as in
NN-G034 forces $\deg R=1$: $R$ cannot be constant
by its parity, and the mass-ten $U_2$ then has
leading coefficient two, so $R$ has leading
coefficient one, which its parity allows only
in degree one.
Therefore $U_2$ has degree $j$ and is the unique
mass-ten character of that degree, namely $W$.
We have proved $U_2=W$ in every case, with no
generic coprimality assumption left unexamined.

\paragraph{Small mass forces the polynomial $C$ itself.}
The two first identities now imply
\[
 \delta_1+\delta_2=(W-D)C.
\]
Put $F=(\delta_1+\delta_2)/2$.
It is an effective symmetric integral polynomial,
dominates $\psi$ coefficientwise, and has mass
six for $w=8$ or nine for $w=10$.
The character $C$ has mass three and satisfies
$C\equiv D-1\pmod2$.

If $W=D+4$ or $D+6$, then $F=2C$ or $3C$.
In particular $C$ is effective. Since $F\ge\psi$,
its coefficients at $q^{\pm1},q^{\pm2}$ force
those of $C$ to be at least one. This would give
$C(4)\ge4$, a contradiction.

For the remaining mass-eight characters,
$F=T_jC$. Its leading coefficient is positive.
If $\deg C>1$, that coefficient in $C$ is even
and at least two by parity. The degree of $F$
then exceeds two, so its outer pair contributes
at least four mass units disjointly from the four
units already required by $\psi$. This contradicts
$F(4)=6$. Thus $\deg C=1$.

For mass ten, $F=(1+T_j)C$. If $\deg C>1$,
the same argument produces an outer pair of
coefficient at least two outside $\psi$.
Since $F(4)=9$, the only possibility is exactly
\[
                         F=\psi+1+2T_k
\]
for some $k>2$: four mass units come from $\psi$,
four from that outer pair, and the remaining
unit must be central. This is impossible modulo
two at $D=1$. Its displayed value is odd there,
whereas $F=(1+T_j)C$ is even because $C(1)$ is
even. Therefore $\deg C=1$ also in this case.

Write $C=aD+b$, with $a$ positive odd and
$4a+b=3$. Its symmetric central coefficient is
$3-2a$ and its $q^1$ coefficient is $a$.
If $w=8,j=1$, the $q^1$ coefficient of $T_1C$
is $3-2a$ and must be at least one, forcing $a=1$.
If $w=10,j=1$, the corresponding coefficient of
$(1+T_1)C$ is $3-a$, again forcing the positive
odd $a$ to be one.
For $j\ge2$, the leading pair is outside $\psi$;
an $a\ge3$ would give total mass at least
$4+2a\ge10$, exceeding six or nine.
Thus in every case
\[
                  C=D-1,\qquad Q=2D-4,\qquad
                  G_1=2D^2(D-2).                    \tag{G036.4}
\]
The three fixed first-cell data of NN-G034 are now
forced, rather than assumed. That theorem proves
that the sole competitive realization in this
branch is NN-G033, of cost exactly $2744$.

\paragraph{The improved finite gate and lower bound.}
Remove from the NN-G035 certificate the five base
rows with $g_1=64$. Every strict improvement must
lie in the remaining $38$ rows, with $348$
auxiliary increment pairs and $20$ first-cell
mass triples. The least remaining relaxed score
is $1932$; the exact projection checker verifies
these four numbers. The removed branch has the
unique actual score $2744$ just proved. This establishes
the intermediate lower bound $1932$; the combined
refinement below strengthens it to $2572$.
\hfill$\square$

\paragraph{Geometric translation and scope.}
Within the NN-G010 realization by graded sign-representation
kernels, a competitive first-cell bottom-action object
of ordinary dimension $64$ determines its entire character
and, in turn, the unique $2744$ completion in this closed
presentation. An improvement requires dimension at least
$72$ for that object and initially obeys $20$
dimension triples, reduced to $15$ by the addendum below. This statement has no bound on the
grading support of the other surviving objects.
The irreducibility argument is a statement in the
integral grading-character ring, not a claim of
irreducibility of the representations themselves.

The dictionary is SGA 3, I, 4.7.2--4.7.3; the character
evaluation at $q=-1$ and the parity reductions are
integral restrictions of this ring. Actual composition
uses archive cote 2, scan 17 through NN-G010.
The proof neither identifies a canonical GCT basis nor
applies this lower bound to the general extra-position
problem. At this intermediate stage only a lower bound was established.
The later exclusions below close the entire closed-table gate;
the extra-position optimization remains separate.


\paragraph{Combined refinement: the 2572 lower bound for closed tables.}
This independently reviewed addendum strengthens the
preceding intermediate bound. The full argument below
excludes twelve more necessary base rows and proves the
closed-table interval $[2572,2744]$.
Its first step excludes the two
NN-G035 base rows
\[
 (w,g_1,g_2,d_1,d_2,u_1,u_2)
   =(8,72,96,6,6,4,8),\quad(8,80,96,8,4,4,8),
\]
with every allowed auxiliary increment, not only their
least auxiliary masses. The mass-ten rows with first-cell
triple $(80,8,4)$ are not excluded by this argument.

In either row, $U_1=D$ and the earlier integral quotient
argument gives $G_1=DH$, $G_2=DWC$,
$H=\delta_1+DC$, $C\equiv D-1\pmod2$, and $C(4)=3$.
The mass-eight $W$ is irreducible and coprime to $D$.
It cannot divide $H$, whose mass is respectively $18$
or $20$, because eight divides neither number.
Thus $W$ is coprime to $G_1$ and both second-cell
identities give $U_2=WR,V_2=WZ$ integrally.
Their parities give $R\equiv Z\equiv1$.
The equal-mass lemma forces $U_2=W$, and hence
\[
             \delta_1+\delta_2=(W-D)C.               \tag{G036.R1}
\]

For the first row, both cells have their unique
mass-six character $\delta_i=\psi+2$, where
$\psi=D(D-3)$. If $W=D+4$, (G036.R1) would give
$C=(\psi+2)/2$, which is not an integral polynomial.
Otherwise $W=D+2T_j$ and
$T_jC=\psi+2=(D-1)(D-2)$.
Since $C\equiv D-1$ has positive degree, this forces
$j=1$, $C=D-1$, $W=3D-4$ and $H=2(D-1)^2$.
Now
\[
                 A_2=D^2WC-DHZ.
\]
If $Z$ has degree one, its leading coefficient is
positive even, so the leading coefficient of $A_2$
is at most $3-4<0$. Higher degrees also fail.
Thus $Z$ is constant, but then $V_2=WZ$ has degree
one and cannot dominate its required $q^2$ coefficient
two. This excludes the first row.

For the second row, $\delta_2=\psi$ and
$F=(\delta_1+\delta_2)/2$ is effective, has mass six,
and dominates $\psi$. The mass-eight part of the
small-mass argument in NN-G036 forces $C=D-1$;
the only possible $W$ are $3D-4$ and $D+2T_2$.
For $W=3D-4$ one has $\delta_1=\psi+4$ and
$H=2D(D-2)+4$. Its degree and leading coefficient
give the same contradiction for $Z$ as in the first row.

For $W=D+2T_2$, the forced polynomials are
\[
 \delta_1=\psi+2T_3,\qquad
 H=2D(D-2)+2T_3,\qquad C=D-1.
\]
Here $D^2WC$ has degree five and leading coefficient
two, while $DH$ has degree four and leading coefficient
two. A nonconstant $Z\equiv1\pmod2$ would again give
a negative leading coefficient in $A_2$.
Consequently $Z$ is a positive odd integer.
The scalar relations give $v_2=8Z$ and
$a_2=384-10v_2$. Since $v_2\ge16$ and $a_2\ge16$,
the only possible odd value is $Z=3$, so
$v_2=24,a_2=144$.

The first auxiliary relation is $a_1=320-12v_1$.
It gives $v_1\le26$; and substitution of every other
now-fixed mass into the total gives
\[
                    T=9760+148(26-v_1).
\]
Since $v_1$ is even, the cap $T\le9912$ forces
$v_1=26$ and $a_1=8$.
We show that these masses cannot satisfy coefficientwise
domination.

Effectivity of $A_1=D^2H-D(D-1)V_1$ gives
$\deg V_1\le3$. Write $V_1=(d,a,b,c)$ in symmetric
Laurent notation. Domination and parity require
$d\ge6$ even, $a\ge5$ odd, $b\ge0$ even, and
$c\ge0$ even. The leading coefficient of $A_1$
also gives $c\le2$.
The mass condition is $d+2a+2b+2c=26$.
The relevant coefficients of $A_1$, after eliminating
$d$, are
\[
 \begin{array}{c|l}
 q^0&-44+2a+6b+8c,\\
 q^1&-24+a+3b+5c,\\
 q^2&14-a-2b-c,\\
 q^3&24-a-3b-4c.
 \end{array}
\]
The target lower bounds at degrees zero and three,
respectively two and one, force
$a+3b+4c=23$. The degree-one target then gives
$c-1\ge1$, so $c=2$ and $a+3b=15$.
The degree-two target gives $b-3\ge1$, hence $b\ge4$.
This implies $a\le3$, contradicting $a\ge5$.
The second row is therefore impossible as well.
\hfill$\square$

Deleting these two rows alone leaves $36$ base rows,
$311$ auxiliary increment pairs, and $19$ first-cell
mass triples, with least relaxed score $1956$.
The following uniform arguments exclude three more
mass blocks.

\paragraph{A common-factor observation.}
The identities used throughout are
\[
 WS=G_1G_2,\qquad U_1G_2=W(G_1-D\delta_1),\qquad
 U_2G_1=W(G_2-D\delta_2).                            \tag{G036.R2}
\]
If a primitive irreducible factor $P$ of $W$ divides
neither $U_1$ nor $U_2$, then it divides both $G_1$
and $G_2$. Indeed it divides at least one from the
first identity, and the appropriate one of the next
two forces it into the other. All divisibilities
first obtained over $\mathbb Q[D]$ are integral by
primitivity. In particular divisibility by $P$ forces
$P(4)$ to divide the corresponding ordinary mass.

\paragraph{The entire mass block $(w,g_1,g_2)=(10,80,80)$.}
Here $u_i=10-d_i/2$, so $u_i\in\{4,6,8\}$.
If $W$ is irreducible, its mass ten divides neither
$u_i$, and (G036.R2) gives $W\mid G_1,G_2$.
If $W=(D-2)B$ is the exceptional case of (G036.2),
the primitive irreducible $B$ has mass five, again
dividing neither $u_i$. Thus write $G_i=BP_i$.
The second and third identities give
\[
 U_1P_2=(D-2)(BP_1-D\delta_1),\qquad
 U_2P_1=(D-2)(BP_2-D\delta_2).
\]
An effective character of mass four, six or eight,
dominating and congruent to $D$, has nonzero value
at $D=2$: these values are $2$, $4$, or one of
$6,2,-2$, respectively. Therefore $D-2$ divides
both $P_i$. We have proved $W\mid G_1,G_2$ also
in this exceptional case.

Now write $G_i=WP_i$ in every case. Then
\[
 P_1(4)=P_2(4)=8,\qquad P_1\equiv0,\quad
 P_2\equiv D(D-1)\pmod2.
\]
The polynomial $P_2$ cannot be divisible by $D$:
its integral quotient would have value two at $D=4$
but would reduce to $D-1$ modulo two, whose value
there is odd. Thus $P_2(0)\ne0$.
The divided identities at $D=0$ show also
$P_1(0)\ne0$, because $W(0)\ne0$.
Multiplying them gives
\[
                        W(0)^2=U_1(0)U_2(0).        \tag{G036.R3}
\]
For the mass-ten characters, $W(0)$ is six or
minus two. For the possible $U_i$, the values are
zero, two, four or minus four. Equality (G036.R3)
is possible only when $W(0)=-2$ and both $U_i(0)=2$.
Consequently
\[
 U_1=U_2=D+2,\qquad d_1=d_2=8,\qquad
 W=D+2+2T_j\quad\hbox{with }j\hbox{ odd}.             \tag{G036.R4}
\]
Each cell character is then $\delta_i=\psi+2R_i$,
where $R_i$ has mass two and is either the constant
two or some $T_k$.

Subtracting the two divided identities in (G036.R2),
and putting $E=R_1-R_2$, gives
\[
 A_0(P_1-P_2)=E,\qquad
 A_0=\frac{D+2+T_j}{D}.                             \tag{G036.R5}
\]
The quotient is integral since $j$ is odd and
$T_j(0)=-2$. The difference $P_1-P_2$ is nonzero,
as its reduction is $D(D-1)$, so $E\ne0$.
Every nonzero difference between two members of
$\{2,T_1,T_2,\ldots\}$ has leading coefficient
$1$ or $-1$, and its value at $D=0$ has absolute
value at most four.

If $j=1$, $A_0=2$, contradicting this odd leading
coefficient of $E$.
For odd $j\ge3$, $A_0$ is monic of degree $j-1$.
Thus $P_1-P_2$ has leading coefficient $\pm1$;
its reduction $D(D-1)$ forces its actual degree
to be exactly two. Write it $aD^2+bD+c$, with
$a=\pm1$, $b$ odd and $c$ even.
The recurrence $T_{j+1}=(D-2)T_j-T_{j-1}$ gives
\[
 T_j'(0)=(-1)^{j-1}j^2,\qquad A_0(0)=1+j^2.
\]
For completeness, the derivative recurrence is
$r_{j+1}=2(-1)^j-2r_j-r_{j-1}$ with
$r_0=0,r_1=1$; the displayed formula satisfies it
by direct substitution.
Evaluating (G036.R5) at zero now gives
$|(1+j^2)c|\le4$, so $c=0$.
But $P_1(4)=P_2(4)$ then implies
$16a+4b=0$, making $b=-4a$ even, a contradiction.
This excludes the entire mass block, including
all six base rows and all their auxiliary choices.

\paragraph{The entire mass block $(w,g_1,g_2)=(8,96,96)$.}
Here $u_i=8-d_i/3$ is four or six.
The mass-eight $W$ is irreducible, and its mass
divides neither $u_i$, so (G036.R2) again gives
$G_i=WP_i$. Now $P_i(4)=12$, with
$P_1\equiv0$ and $P_2\equiv D(D-1)$ modulo two.
The polynomial $P_1$ cannot be divisible by $D$:
the quotient would have even coefficients but
ordinary mass $12/4=3$.
The divided identities at zero consequently
force both $P_i(0)\ne0$ and give (G036.R3).
But $W(0)^2=16$, whereas $U_i(0)$ are zero or two,
whose products are zero or four. This contradiction
excludes both base rows in this block.

\paragraph{The two asymmetric mass-ten blocks.}
The finite certificate has one base row for
$(w,g_1,g_2)=(10,76,100)$, with
$(d_1,d_2,u_1,u_2)=(4,6,6,10)$, and its index-swapped
row with $(g_1,g_2)=(100,76)$.
Let $G_-$ denote the smaller action character and
$G_+$ the larger. Their cell characters are
respectively $\psi$ and $\psi+2$, and the smaller
index has $U_-=D+2$.

If $W$ is irreducible, it cannot divide $G_-$,
since ten does not divide $76$.
It therefore divides $G_+$ and then $U_+$ by
(G036.R2). Since $u_+=10$, the equal-mass lemma
forces $U_+=W$.
If $W=(D-2)B$, the irreducible mass-five $B$ cannot
divide $G_-$, so the same argument gives $B\mid U_+$.
Writing $U_+=BR$, its reduction is $R\equiv D$.
The leading-coefficient argument from NN-G036 forces
$\deg R=1$ and then the mass-ten classification gives
$U_+=W$ in this case as well.

The larger-index identity now gives
$G_+=G_-+D(\psi+2)$.
Substitution into the smaller-index identity yields
\[
 (W-D-2)G_-=D\bigl(W\psi+(D+2)(\psi+2)\bigr).       \tag{G036.R6}
\]
If $W=D+6$, its denominator is four and the resulting
$G_-$ has leading coefficient $1/2$, contradicting
integrality.
For $W=D+2+2T_j$, equation (G036.R6) gives
\[
 G_-=D\psi+\frac{N(D)}{T_j(D)},\qquad
 N(D)=D(D+2)(D^2-3D+1).
\]
Integrality forces the monic $T_j$ to divide $N$.
Thus $j\le4$. None does:
$N(2)=-8$ excludes $T_1=D-2$ and also $T_3$,
which contains $T_1$;
the remainder modulo $T_2=D^2-4D+2$ is $16D-10$;
and $T_4=D^4-8D^3+20D^2-16D+2$ is not the
monic degree-four polynomial $N$.
This excludes both asymmetric blocks, uniformly
over all original support indices.
\hfill$\square$

\paragraph{The resulting finite dimension gate.}
In addition to the five rows already removed in
NN-G036, this addendum removes two isolated mass-eight
rows, six rows in the $(10,80,80)$ block, two in the
$(8,96,96)$ block, and the two asymmetric rows.
The remaining complete necessary projection has
$26$ base rows, $42$ auxiliary increment pairs,
and $15$ first-cell mass triples. Its least relaxed
score is $2572$. The proved exact attainer still
has score $2744$, so the closed-table minimum lies
in $[2572,2744]$.

In the NN-G010 geometry, these are exclusions of
ordinary-dimension patterns for actual composing
graded sign-representation objects. The tests at
$D=0$ are integral character evaluations at $q=-1$;
the derivative there is a character-polynomial
calculation, not a Weil-weight assertion.
They preserve the SGA 3, I, 4.7.2--4.7.3 and
archive cote 2, scan 17 dictionary of NN-G036.
The surviving rows are not declared realizable,
and this lower bound is not transferred to arbitrary
extra-position lifts or canonical GCT multiplicities.


% Further NN-G036 refinement: root whole-proof PASS 2026-11-22T00:03:47Z.
\paragraph{The last mass-eight competitor.}
The remaining mass-eight base row has
\[
 (g_1,g_2,d_1,d_2,u_1,u_2)=(128,80,12,4,8,4).
\]
Thus $\delta_2=\psi$ and $U_2=D$. The integral divisibility
argument already used in NN-G035 gives
\[
 G_1=DWC,\qquad G_2=D(\psi+DC),\qquad
 C\equiv0\pmod2,\quad C(4)=4.
\]
For clarity, $DG_1=W(G_2-D\psi)$ and $\gcd(W,D)=1$
first give $G_1=WQ$ and $G_2=D(Q+\psi)$ integrally.
The other cell identity at zero forces $G_1(0)=0$,
so $Q(0)=0$ and $Q=DC$. Reducing $G_1=DWC$ modulo
two gives the displayed parity of $C$.
The irreducible mass-eight $W$ cannot divide $G_2$:
it is coprime to $D$, whereas the integral polynomial
$\psi+DC$ has value twenty at four, not divisible by eight.
Consequently $U_1G_2=W(G_1-D\delta_1)$ forces $W\mid U_1$.
The equal-mass lemma gives $U_1=W$. Cancelling now yields
\[
 \delta_1+\delta_2=(W-D)C.
\]
Write $C=2H$ with $H\in\mathbb Z[D]$ and $H(4)=2$.
The effective character
$F=(\delta_1+\delta_2)/2=(W-D)H$ has mass eight and
dominates $\psi$.
If $W=D+4$, then $F=4H$, so $H$ is effective;
domination forces its coefficients at both positive
exponents one and two to be at least one. Its mass
would therefore be at least four, contradicting two.
Otherwise $W=D+2T_j$ and $F=2T_jH$ has all Laurent
coefficients even. Its coefficients at exponents
$1,-1,2,-2$ must all be at least two. These exhaust
the mass eight, hence $F=2\psi$ and $T_jH=\psi$.
Since $T_j$ is monic of degree $j$, this requires
$j\le2$. But $T_1=D-2$ does not divide $\psi$, and
$T_2=D^2-4D+2$ is not the monic quadratic $\psi$.
This excludes the last mass-eight competitor.
Together with the preceding exclusions and NN-G034,
every competitive mass-eight closed table is the unique
$2744$ attainer. No bound on its initial support indices
was imposed.

\paragraph{The mass-ten block with equal action masses one hundred.}
The row is
\[
 (g_1,g_2,d_1,d_2,u_1,u_2)=(100,100,10,10,6,6).
\]
Thus both $U_i=D+2$. If $W$ is irreducible, its mass
ten divides neither $u_i$, so (G036.R2) forces
$W\mid G_1,G_2$. In the exceptional factorization
$W=(D-2)B$, the irreducible $B$ has mass five and
again divides neither $U_i$; it divides both $G_i$.
After cancelling $B$ in the cell identities, evaluation
at $D=2$, where $U_i(2)=4\ne0$, forces the remaining
factor $D-2$ into both $G_i$. Thus in every case
\[
 G_i=WP_i,\qquad P_i(4)=10,\qquad
 P_1\equiv0,\quad P_2\equiv D(D-1)\pmod2.
\]
The polynomial $P_2$ cannot be divisible by $D$,
since its quotient would have nonintegral value $10/4$
at four. The divided cell identities at zero then
force both $P_i(0)\ne0$ and give $W(0)^2=4$.
Consequently $W=D+2+2T_j$ with $j$ odd.
Write $\delta_i=\psi+2R_i$. Now $R_i$ is an effective
character of mass three, hence is either three or
$1+T_k$. A nonzero difference $E=R_1-R_2$ still
has leading coefficient $\pm1$ and $|E(0)|\le4$.
It is nonzero because the parity difference of
$P_1,P_2$ is nonzero.
Subtracting the divided identities gives exactly
\[
 \frac{D+2+T_j}{D}(P_1-P_2)=E.
\]
For $j=1$, the factor is two, contradicting the odd
leading coefficient of $E$. For odd $j\ge3$, it is
monic, with value $1+j^2$ at zero. Hence
$P_1-P_2=aD^2+bD+c$, where $a=\pm1$, $b$ is odd,
and $c$ is even: the degree follows from its leading
coefficient and its reduction $D(D-1)$.
The bound on $E(0)$ forces $c=0$.
Equal values of the two $P_i$ at four then give
$16a+4b=0$, contradicting that $b$ is odd.
The row is impossible.

\paragraph{The two mass-ten blocks with action masses eighty-four and one hundred.}
In either orientation the smaller action has
$\delta_-=\psi+2$, $U_-=D+2$, while the larger
has $\delta_+=\psi$ and $u_+=10$.
An irreducible $W$ of mass ten cannot divide $G_-$
of mass eighty-four; (G036.R2) consequently forces
$W\mid G_+$ and $W\mid U_+$, whence $U_+=W$
by the equal-mass lemma.
If $W=(D-2)B$, the mass-five irreducible factor
$B$ likewise cannot divide $G_-$. Thus $B\mid U_+$.
Writing $U_+=BR$, reduction gives $R\equiv D\pmod2$.
The leading coefficient argument in NN-G036 forces
$\deg R=1$; the mass-ten classification then gives
$U_+=W$ also in this case.
The larger-index identity now says $G_+=G_-+D\psi$.
Substituting in the smaller-index identity gives
\[
 (W-D-2)G_-=D\bigl(W(\psi+2)+(D+2)\psi\bigr).
\]
For $W=D+6$ this would make the leading coefficient
of $G_-$ equal to one half. Otherwise
$W=D+2+2T_j$ and
\[
 G_-=D(\psi+2)+\frac{N(D)}{T_j(D)},\qquad
 N=D(D+2)(D^2-3D+1).
\]
The same integral divisibility obstruction as in
(G036.R6) applies: $T_j\mid N$ requires $j\le4$,
but $N(2)=-8$, $T_1\mid T_3$,
$N\bmod T_2=16D-10$, and $N\ne T_4$.
Both orientations are excluded, with every auxiliary
increment allowed by their scalar relaxation.
\hfill$\square$

\paragraph{Scope and geometric content.}
These four base-row exclusions concern the ordinary
dimensions of actual graded $C_2$-representation
structure objects in the matrix-kernel category of
NN-G010, with the same symmetric closed multiplication
table and character constraints. In particular, within
the mass-eight gluing stratum the actual smallest
effective lift is the $2744$-pair object of NN-G033,
uniquely within that presentation. The integral torus
character interpretation is SGA~3, I, 4.7.2--4.7.3;
kernel composition is the supplied archive transcription,
cote~2, scan~17. No new claim about a canonical GCT
representation family or a different multiplication
table follows.

\paragraph{A mass-twelve gluing lemma.}
Throughout, $D=q+2+q^{-1}$ and $T_j=q^j+q^{-j}$.
Every effective $W\ge D$ with $W\equiv D\pmod2$ and $W(4)=12$
is one of
\[
 D+8,\quad D+4+2T_j,\quad D+4T_j,\quad
 D+2T_i+2T_j\ (i<j).
\]
Indeed $(W-D)/2$ is an effective symmetric character of mass four,
and these are its complete partitions into a central coefficient
and paired nonzero exponents. Hence $W(0)\in\{-8,0,8\}$.
If $W(0)=0$, write $W=DH$. Then
\[
 H\in\mathbb Z[D],\quad H(4)=3,\quad H\equiv1\pmod2,
 \quad H(0)\ne1,-1.
\]
Either $H=3$ (the case $W=3D$), or $H$ is primitive and
irreducible over $\mathbb Q[D]$, with leading coefficient two.
For the last assertion, every nonconstant case has leading
coefficient two and odd constant coefficient. In a nontrivial
primitive integral factorization, one factor has odd leading
coefficient and would keep its positive degree modulo two,
contradicting that the product reduces to one.
For the assertion at zero, the zero-valued cases in the list are
$D+4+2T_o$ with $o$ odd and $D+2T_o+2T_e$ with $o$ odd and
$e$ even. The derivative identity from (G036.R5) gives
\[
 H(0)=1+2o^2\quad\hbox{or}\quad 1+2(o^2-e^2).
\]
Neither is one or minus one: equality would require equal
squares of opposite parity or $e^2-o^2=1$, impossible for
positive distinct integers $e,o$. This also includes $o=1$
and the constant $H=3$.

\paragraph{The nonzero-value branch.}
For any effective $U\ge D$ with $U\equiv D\pmod2$ and mass $u$,
\[
 |U(0)|\le u-4,
\]
because $(U-D)/2$ is effective and its absolute value at $q=-1$
is at most its mass. In every remaining mass-twelve row of
NN-G036, $(u_1-4)(u_2-4)<64$, as follows directly from the
complete listed mass tuples. If $W(0)=8$ or $-8$, the identities
\[
 WG_1-U_1G_2=DW\delta_1,\qquad
 WG_2-U_2G_1=DW\delta_2
\]
have an invertible coefficient matrix at $D=0$ over $\mathbb Q$.
Thus $D$ divides both $G_i$ in $\mathbb Z[D]$.
Using $G_1\equiv0$ and $G_2\equiv D^2(D-1)\pmod2$, their
quotients by $D$ have even evaluation at $D=4$. Consequently
both $g_i$ are divisible by eight. This excludes any mass row
containing $84$ in this branch.
If also $\delta_1=\delta_2=\psi=D(D-3)$, divide the same
identities by $D$ and apply the invertible matrix at zero once
again. Both $G_i$ are divisible by $D^2$. Since $WS=G_1G_2$
and $W(0)\ne0$, the polynomial $S$ is divisible by $D^4$.
In particular, the mass block $(g_1,g_2)=(80,96)$ or $(96,80)$,
where $s=640$, is impossible: $4^4$ does not divide $640$.

\paragraph{The two blocks with action masses $84,88$.}
Their action masses are $84,88$ in either order, with associated
$u$-masses six and eight. The preceding paragraph leaves only
$W=DH$. If $H$ is nonconstant, it cannot divide the mass-six
character $D+2$ (the integral leading coefficient would have
to divide one), nor a mass-eight character (its value three
would have to divide eight). The common-factor argument of
(G036.R2) forces $H$ to divide both $G_i$, contradicting that
three does not divide $88$.
If $H=3$, reduction modulo three supplies the same argument
in the integral domain $\mathbb F_3[D]$. Both $U_i$ have
nonzero reduction: one is $D+2$, and the other has mass eight.
The three identities (G036.R2) then force both $G_i$ to have
all coefficients divisible by three, again contradicting
mass $88$. Thus both base rows are impossible.

\paragraph{The block with both action masses $84$.}
Here $U_1=U_2=D$ and $d_1=d_2=14$. Again only $W=DH$
remains. The common-factor argument (or reduction modulo
three when $H=3$) gives $G_i=HP_i$ integrally. Their masses
are $P_i(4)=28$, and their reductions are
$P_1\equiv0$, $P_2\equiv D^2(D-1)\pmod2$.
Cancelling $DH$ in the two cell identities gives
\[
 P_2=HP_1-D\delta_1,\qquad
 P_1=HP_2-D\delta_2.
\]
Since $H(0)^2\ne1$, evaluating at zero forces
$P_1(0)=P_2(0)=0$. In particular $P_1/D$ is an integral
polynomial with every coefficient even, but its evaluation
at four is $28/4=7$. This is a contradiction.

\paragraph{The two smallest mass-twelve rows.}
Suppose $(g_1,g_2)$ is $(80,96)$ or $(96,80)$ and both
cells are $\psi$. Let $G_-$ have mass eighty and $G_+$
have mass ninety-six; their $U$-masses are eight and twelve.
The nonzero-value branch was already excluded, so $W=DH$.
For nonconstant $H$, its value three cannot divide eighty,
so $H\nmid G_-$. The identities (G036.R2) force
$H\mid G_+$ and $H\mid U_+$. Every mass-twelve $U_+$
has value at zero equal to zero or $\pm8$. Since $H(0)$ is
odd and is neither one nor minus one, its divisibility
forces $U_+(0)=0$. Write $U_+=DJ$. The classification above
applies to $J$: it is either three or primitive irreducible
of leading coefficient two. Since $H\mid J$, the latter
case is forced and $J=H$. Thus $U_+=W$.
For $H=3$, reduce $U_+G_-=3D(G_+-D\psi)$ modulo three.
Since $G_-$ has mass eighty, its reduction is nonzero, so
all coefficients of $U_+$ are divisible by three. Its odd
coefficients at $q,q^{-1}$ are at least three each, and its
even central coefficient is at least six. These exhaust its
mass twelve, so $U_+=3D=W$ also in this case.

The larger-cell identity now gives $G_+=G_-+D\psi$.
The other identity becomes
\[
             (W-U_-)G_-=D\psi(W+U_-).
\]
The mass-eight classification gives $U_-(0)=4$ or $-4$,
while $W(0)=0$. Hence $D^2\mid G_-$.
If $G_-=G_1$, its quotient by $D^2$ has even coefficients
but mass $80/16=5$, a contradiction.
If $G_-=G_2$, use the necessary mod-four identity
$G_2-D\delta_2\equiv G_1\pmod4$ from NN-G030.
Here $\delta_2=\psi$ and $G_1=G_2+D\psi$, so it would
require $2D\psi\equiv0\pmod4$, which is false.
This excludes both rows, including every auxiliary increment.

These five exclusions use neither a degree bound nor
positivity of any quotient. They retain precisely the
closed-table character dictionary and scope of NN-G036.

\paragraph{Intermediate combined gate at 00:03:47 UTC.}
The additional four small-gluing and five mass-twelve exclusions
leave exactly $17$ necessary base rows, $26$ auxiliary increment
pairs, and $8$ first-cell mass triples. The least remaining
relaxed score is $2604$, so the intermediate closed-table
interval at this stage is $[2604,2744]$. Every strict improvement has $w=10$
or $12$; in the mass-ten case its first data must be
$(g_1,g_2,d_1,u_1)=(80,120,8,4)$, with
$d_2\in\{6,10,14,18\}$. The earlier numerical intervals in
this subsection are successive intermediate bounds.
The exact finite projection is
\texttt{verify\_second\_competitive\_refinement.py};
the preceding polynomial arguments justify every deletion
without a support-degree bound.

\paragraph{Geometric conclusion at the intermediate stage.}
Within the specified closed presentation of composing graded
$C_2$-representation objects from NN-G010, every effective
lift has at least $2604$ cancelling trivial/sign pairs and
one with $2744$ pairs exists. In the gluing-dimension-eight
stratum, that $2744$ object is the unique minimum in this
presentation. A strict improvement can occur only in the
seventeen displayed ordinary-dimension rows, whose actual
character realization remains a further condition.
The bounds concern genuine effective representation objects,
using SGA~3, I, 4.7.2--4.7.3 and the archive kernel-product
transcription at cote~2, scan~17. They are not statements
about other multiplication-table shapes or canonical GCT
multiplicities.


% NN-G036 further uniform refinement; root whole-proof PASS 2026-11-22T00:06:53Z.
\paragraph{The entire remaining mass-ten gluing branch.}
The reviewed finite projection leaves four mass-ten base rows:
\[
 \begin{aligned}
 &g_1=80,\quad g_2=120,\quad d_1=8,\quad U_1=D,\\
 &(d_2,u_2)\in\{(6,12),(10,10),(14,8),(18,6)\}.
 \end{aligned}
\]
Every allowed auxiliary increment in these rows has
\[
 a_1=8,\qquad v_1=26,\qquad v_2\in\{56,58\}.
\]
Equality with the first auxiliary's absolute trace mass
therefore fixes the entire polynomial, not only its mass:
\[
 A_1=D(D-2)(D-3)
   =2+q+q^{-1}+q^2+q^{-2}+q^3+q^{-3}.                 \tag{G036.T1}
\]
We exclude all four rows, with arbitrary support locations.

The integral $U_1=D$ argument gives
\[
 \begin{aligned}
 G_2&=DWC,\qquad G_1=DH,\qquad H=\delta_1+DC,\\
 C&\equiv D-1\pmod2,\qquad H\equiv0\pmod2,\\
 C(4)&=3,\qquad H(4)=20.
 \end{aligned}                                                   \tag{G036.T2}
\]
The mass-eight cell is necessarily
$\delta_1=\psi+2R$, where $R$ is either two or $T_k$.
In particular $\delta_1(0)$ is four or minus four.

First suppose $W$ is irreducible. If it did not divide
$H$, it would be coprime to $G_1=DH$, because $W(0)\ne0$.
The second auxiliary identity
$V_2G_1=W(DG_2-A_2)$ would force $W\mid V_2$ integrally.
Evaluation at four would then give $10\mid v_2$, contrary
to $v_2=56$ or $58$. Thus $W\mid H$.

In the exceptional case $W=(D-2)B$ of (G036.2),
$B$ is primitive irreducible, $B(4)=5$,
$B\equiv1\pmod2$, and $B(0)=-3$.
If $B\nmid H$, the same auxiliary identity gives
$B\mid V_2$, contradicting $5\nmid v_2$.
Therefore $H=2BK$ for an integral polynomial $K$;
the factor two follows by reducing $H$ modulo two.
But then $H(0)$ is divisible by six, whereas
$H(0)=\delta_1(0)=\pm4$. The exceptional case is impossible.

We are left with the irreducible case, where
\[
 H=2WK,\qquad K\in\mathbb Z[D],\qquad K(4)=1.
\]
Here the integral quotient by $W$ is even because
$\overline W=D$ and $\overline H=0$ in the domain
$\mathbb F_2[D]$. If $W(0)=6$, the identity
$H(0)=2W(0)K(0)=\pm4$ is impossible. Hence
\[
 W=D+2+2T_j,\quad j\ {\rm odd},\qquad W(0)=-2,
 \qquad K(0)\in\{1,-1\}.
\]
Every integer polynomial satisfies $K(4)\equiv K(0)\pmod4$.
Since $K(4)=1$, this forces $K(0)=1$.
Consequently $\delta_1(0)=H(0)=-4$, so
$R=T_k$ with $k$ odd. At $D=2$ the Chebyshev
evaluations are therefore
\[
 W(2)=4,\qquad \delta_1(2)=\psi(2)+2T_k(2)=-2.
                                                                  \tag{G036.T3}
\]

The first auxiliary identity, after cancelling $W$
and then $D$, is
\[
 C V_1=2DWK-A_1/D.
\]
By (G036.T1), evaluation at two gives
\[
                  C(2)V_1(2)=16K(2).               \tag{G036.T4}
\]
The integer $C(2)$ is odd by (G036.T2).
Therefore $16$ divides $V_1(2)$.
The coefficientwise lower bound for $V_1$ is the
effective character $V_{\min}=6+5(q+q^{-1})$,
of mass sixteen. Its parity and mass twenty-six imply
$V_1=V_{\min}+2P$ for an effective symmetric integral
character $P$ of mass five. At $D=2$, equivalently
$q=i$, the triangle inequality gives
$|P(i)|\le5$ and hence
\[
                         -4\le V_1(2)\le16.
\]
Its only multiples of sixteen in this interval are
zero and sixteen.
If $V_1(2)=0$, (G036.T4) gives $K(2)=0$.
The monic polynomial $D-2$ would divide $K$ integrally,
so its value two at four would divide $K(4)=1$,
which is impossible.
Thus $V_1(2)=16$, and (G036.T4) gives $K(2)=C(2)$.
But (G036.T2)--(G036.T3) also give
\[
 C(2)=\frac{2W(2)K(2)-\delta_1(2)}2
      =4K(2)+1.
\]
This would require $-3K(2)=1$, contradicting integrality.
Every mass-ten competitor is excluded.
\hfill$\square$

\paragraph{Meaning of this refinement.}
The closed-table optimum is consequently attained at
gluing dimension eight unless a strictly smaller object
occurs in the remaining dimension-twelve stratum.
The proof uses restriction of integral torus characters
at $q=-1$ and $q=i$, together with ordinary dimensions
at $q=1$; it assumes no Weil-weight estimate and no
effectivity of $C,H/W$, or $K$.
Its geometric objects and source dictionary remain the
actual composing graded sign representations of NN-G010,
SGA~3, I, 4.7.2--4.7.3, and the supplied kernel-product
transcription at archive cote~2, scan~17.

\paragraph{Projection after the mass-ten exclusion.}
Deleting its four base rows and six auxiliary increments leaves
$13$ necessary base rows, $20$ auxiliary increments and
$7$ first-cell mass triples. Every strict improvement on
$2744$ now has gluing dimension twelve. The least relaxed
score remains $2604$. The exact row projection is certified
separately by \texttt{verify\_mass10\_final\_exclusion.py}.


% Further NN-G036 refinement; root whole-proof PASS 2026-11-22T00:12:35Z.
\paragraph{A mass-six action gives integral divisibility.}
Suppose $W\in\mathbb Z[D]$ satisfies $W\equiv D\pmod2$
and $w=W(4)$ is divisible by four.
Write $W=D+2R$. Then $R(4)=(w-4)/2$ is even,
and $R(-2)\equiv R(4)\pmod2$. Therefore
\[
                         W(-2)\equiv2\pmod4.        \tag{G036.U1}
\]
In particular $W$ and $D+2$ are coprime over
$\mathbb Q[D]$.
If $(D+2)G=W F$ for integral polynomials, then
$W\mid G$ in $\mathbb Z[D]$, even when $W$ is
not primitive. To check the content explicitly,
write $W=cW_0$ with $W_0$ primitive and $c>0$.
Gauss's content identity and primitivity of $D+2$
first give $c\mid G$. Write $G=cG_0$.
Coprimality gives $W_0\mid G_0$ over $\mathbb Q[D]$,
and primitivity makes this quotient integral.

Consequently, in any of the closed tables under
consideration, $u_i=6$ forces $U_i=D+2$, and
\[
                 4\mid w,\ u_i=6
                 \quad\Longrightarrow\quad w\mid g_j
                 \qquad(i\ne j).                  \tag{G036.U2}
\]
At $w=12$ this excludes both $(g_1,g_2)=(72,112)$
rows, since $u_1=6$ but $12\nmid112$.
It also excludes the three $(80,96)$ rows with
$d_2=14,u_2=6$, and their three index-reversed rows:
in each case the opposite action has mass eighty.
Thus eight entire base rows, with all their auxiliary
increments, are impossible.

\paragraph{The forward row with action masses eighty and ninety-six.}
Consider
\[
 (w,g_1,g_2,d_1,d_2,u_1,u_2)
                    =(12,80,96,12,4,4,12).         \tag{G036.U3}
\]
Here $U_1=D$ and $\delta_2=\psi$.
If $W(0)\ne0$, the integral quotient argument for
$U_1=D$ gives
$G_2=DWC$ with $C\equiv D-1\pmod2$.
It applies to nonprimitive $W$ as well by the same
content argument just proved. Its value would be
$C(4)=96/(4\cdot12)=2$, whereas its displayed
parity requires an odd value. This is impossible.

It remains that $W=DH$, with the mass-twelve
classification proved above: $H(4)=3$,
$H\equiv1\pmod2$, and $H(0)\ne\pm1$.
Either $H=3$, or $H$ is primitive irreducible
of leading coefficient two.
Since three does not divide $g_1=80$, in the
nonconstant case $H\nmid G_1$.
The second cell identity then forces $H\mid U_2$.
As in the earlier mass-twelve proof, $H(0)$ is
odd and not $\pm1$, whereas $U_2(0)$ is zero
or $\pm8$; therefore $U_2(0)=0$.
Writing $U_2=DJ$, both $H,J$ are primitive
irreducible of leading coefficient two and mass
three, so $J=H$ and $U_2=W$.
For $H=3$, reducing the same identity modulo
three instead gives $3\mid U_2$ coefficientwise;
its effective mass twelve and parity force
$U_2=3D=W$. Thus $U_2=W$ in every case.

The two cell identities now imply
\[
 G_2=G_1+D\psi,\qquad
                    (H-1)G_1=D(H\delta_1+\psi).     \tag{G036.U4}
\]
Write $\delta_1=\psi+2R$, where $R$ is effective
of mass four. Then $\delta_1(0)\in\{-8,0,8\}$.
If $\delta_1(0)=0$, both $\delta_1$ and $\psi$
are divisible by $D$. Since $H(0)\ne1$,
(G036.U4) forces $D^2\mid G_1$.
The quotient has even coefficients but value
$80/16=5$ at four, a contradiction.

If $\delta_1(0)=\pm8$, first $D\mid G_1$.
Divide (G036.U4) by $D$ and evaluate at zero.
Coprimality of the consecutive integers
$H(0),H(0)-1$ shows that $H(0)-1$ divides eight.
The explicit root classification gives
\[
 H(0)=1+2o^2\quad\hbox{or}\quad
 H(0)=1+2(o^2-e^2),
\]
with $o$ positive odd and $e$ positive even.
In particular $H(0)\equiv3\pmod8$.
The only such odd integer with $H(0)-1\mid8$
is three. The first formula then forces $o=1$;
the second would require $o^2-e^2=1$,
impossible for positive $e$.
Consequently $H=3$ and $W=3D$.

Equation (G036.U4) becomes
$G_1=D(2\psi+3R)$.
Its even parity forces $R=2P$ integrally, where
$P$ is effective of mass two, hence either two
or some $T_k$. Thus
\[
 \delta_1=\psi+4P,\qquad
 G_1=D(2\psi+6P),\qquad
 G_2=D(3\psi+6P).
\]
The central Laurent coefficient of $2D\psi$
is four. For $P=T_k$ with $k\ge2$, that of
$6DP$ is zero, whereas $G_1$ must have central
coefficient at least sixteen.
Therefore $P=2$ or $T_1$.

Put $L=2\psi+6P$, an integral quadratic with
leading coefficient two and $L(4)=20$.
The second auxiliary identity gives
$V_2L=3(DG_2-A_2)$.
Modulo three, $L$ reduces to $2\psi\ne0$,
so $V_2=3Z$ integrally.
Since $V_2$ is effective and $V_2\equiv D\pmod2$,
$Z$ is effective and $Z\equiv D\pmod2$.
The actual second auxiliary coefficient is
\[
              A_2=D^2(3\psi+6P)-Z(2\psi+6P).       \tag{G036.U5}
\]
The first term has degree four and leading
coefficient three. A $Z$ of degree greater
than two makes the leading coefficient of
$A_2$ negative. If $\deg Z=2$, its leading
coefficient is positive even and at least two,
so the degree-four coefficient is at most
$3-2\cdot2<0$.
Thus effectivity forces $\deg Z\le1$.
But then $V_2=3Z$ has no $q^2$ coefficient,
contradicting its required coefficient at least two.
The forward row (G036.U3) is excluded.
\hfill$\square$

\paragraph{Geometric scope.}
The divisibility lemma removes ordinary-dimension
patterns for effective graded sign-representation
objects in the specified closed kernel presentation.
It does not assume positive polynomial quotients:
content and integrality alone supply the character
divisibilities. The forward-row exclusion further
uses actual coefficientwise effectivity of the
structure objects. These are the same mathematical
objects of NN-G010 and SGA~3, I, 4.7.2--4.7.3,
with composition in the supplied archive
cote~2, scan~17 transcription. No change of
multiplication-table presentation is permitted
in this conclusion.

\paragraph{The two mass-twelve rows of relaxed cost $2604$.}
Let $G_-$ have mass eighty and $G_+$ have mass ninety-six,
in either orientation of the closed table. The two rows now considered have
\[
 W(4)=12,\qquad \delta_- =\psi+2R,\quad U_-=D+2,
 \qquad \delta_+=\psi,\quad U_+(4)=12,
\]
where $\psi=D(D-3)$ and $R$ is either the constant two or $T_k$,
with $k\ge1$. Every allowed auxiliary increment is included.
In particular $\delta_-(0)=4$ or $-4$.

We first record the divisibility interface used in both orientations.
Write $W=D+2B$ with $B\in\mathbb Z[D]$. Since $W(4)=12$,
$B(4)=4$; therefore $B(-2)$ is even, and
$W(-2)\equiv2\pmod4$. Thus $W$ is coprime to $D+2$
over $\mathbb Q[D]$. From the minus-cell identity
\[
                 (D+2)G_+=W(G_--D\delta_-)
\]
it follows that $W$ divides $G_+$ in $\mathbb Z[D]$.
Here integrality follows either by Gauss's lemma and content,
or directly as follows. Write $W=cW_0$, with $W_0$ primitive.
The rational coprimality gives $G_+=W_0Q$ with $Q$ integral.
After cancellation, $c$ divides every coefficient of $(D+2)Q$;
the product-content identity and primitivity of $D+2$ imply
that $c$ divides every coefficient of $Q$.

If $W(0)\ne0$, the preceding mass-twelve classification gives
$W(0)=\pm8$ and $U_+(0)\in\{-8,0,8\}$.
The two cell identities at zero force both $G_i(0)=0$:
their coefficient determinant is
$W(0)^2-U_-(0)U_+(0)=64-2U_+(0)\ne0$.
Write $G_\pm=DP_\pm$ integrally. Dividing the identities by
$D$ and evaluating at zero yields
\[
 wP_-(0)-2P_+(0)=w\delta_-(0),\qquad
 wP_+(0)-uP_-(0)=0,
 \quad w=W(0),\ u=U_+(0).
\]
Consequently
\[
 P_-(0)=\frac{64\delta_-(0)}{64-2u}.
\]
For $u=8$ or $-8$ this is respectively
$4\delta_-(0)/3$ or $4\delta_-(0)/5$, neither an integer.
Hence $u=0$ and $P_+(0)=0$, so $D^2$ divides $G_+$.
On the other hand $G_+=WQ$ with $Q(4)=96/12=8$.
Since $W(0)\ne0$, $D^2$ divides $Q$ integrally.
Evaluation at four would make sixteen divide eight, a contradiction.

It remains to treat $W(0)=0$, so $W=DH$ with $H$ as in
the preceding mass-twelve gluing lemma. If $H$ is nonconstant,
it is primitive irreducible with leading coefficient two,
$H(4)=3$, and odd $H(0)\ne1,-1$.
It cannot divide $G_-$, since three does not divide eighty.
The plus-cell identity therefore forces $H\mid U_+$.
Divisibility at zero, together with $U_+(0)\in\{-8,0,8\}$,
forces $U_+(0)=0$. Write $U_+=DJ$.
The same classification makes $J$ either three or a primitive
irreducible polynomial of leading coefficient two and value three
at four. Since $H\mid J$, the constant case is impossible and
$J=H$. Thus $U_+=W$.
If $H=3$, reduce the plus-cell identity modulo three.
The reduction of $G_-$ is nonzero because its mass is eighty,
so all coefficients of $U_+$ are divisible by three.
Its odd coefficients at $q,q^{-1}$ are at least three each,
and its even central coefficient at least six. These exhaust
mass twelve, whence again $U_+=3D=W$.

Cancelling $W$ in the plus-cell identity gives
$G_+=G_-+D\psi$. The other cell identity becomes
\[
 (W-D-2)G_-=D\bigl(W\delta_-+(D+2)\psi\bigr).
\]
Its right side is divisible by $D^2$, whereas the first factor
on the left has value $-2$ at zero. Therefore $D^2\mid G_-$.
If $G_-=G_1$, its quotient by $D^2$ has every coefficient even
but evaluation $80/16=5$ at four, a contradiction.
If $G_-=G_2$, the necessary congruence of NN-G030 is
\[
 G_--D\delta_-\equiv G_+\pmod4.
\]
Substituting $G_+=G_-+D\psi$ gives
$2D(\psi+R)\equiv0\pmod4$, so $R\equiv\psi\pmod2$.
This is impossible. The constant two reduces to zero; if
$R=T_k$, its monic degree is $k$, so equality with the degree-two
polynomial $\psi$ would require $k=2$. But
$T_2\equiv D^2$ and $\psi\equiv D^2+D$ modulo two.
Both orientations are excluded for arbitrary grading support.
\hfill$\square$

The objects and scope are those of NN-G036: actual effective graded
sign representations composing in the symmetric closed table.
Restriction at $q=-1$, ordinary dimension at $q=1$, and integral
character divisibility are operations in the torus representation
ring of SGA~3, I, 4.7.2--4.7.3. No quotient is assumed effective,
and no statement about extra-position tables is inferred.

\paragraph{The reverse mass-twelve row with a mass-four action.}
We now exclude the remaining orientation
\[
 (w,g_1,g_2,d_1,d_2,u_1,u_2)
                  =(12,96,80,4,12,12,4).
\]
Write $G_+=G_1$, $G_-=G_2$, $U_+=U_1$, $U_-=U_2=D$,
$\delta_+=\psi=D(D-3)$, and $\delta_-=\psi+2R$.
Here $R$ is an effective symmetric character of mass four.
In particular $|\delta_-(0)|\le8$ and $|U_+(0)|\le8$.

Suppose first $W(0)\ne0$, so $W(0)=\pm8$.
The minus-cell identity
$DG_+=W(G_--D\delta_-)$ and coprimality of $W,D$
give $G_+=WQ$ with $Q$ integral, including by the content
argument if $W$ is nonprimitive.
Then $G_-=D(Q+\delta_-)$.
The plus-cell identity at zero implies $Q(0)=0$, because
its left side is $U_+(0)G_-(0)=0$ and its right side is
$W(0)^2Q(0)$.
Thus $Q=DC$ and
\[
 G_+=DWC,\qquad G_-=D\delta_-+D^2C,\qquad C(4)=2.
\]
The trace parity $G_1\equiv0$, together with $W\equiv D$
modulo two, forces $C\equiv0\pmod2$.
Dividing the plus-cell identity by $D$ now gives
\[
 U_+(\delta_-+DC)=W(WC-\psi).
\]
At zero this says
$64C(0)=U_+(0)\delta_-(0)$.
The right side has absolute value at most sixty-four.
Therefore $|C(0)|\le1$; its even parity forces $C(0)=0$.
The integral polynomial $C$ would be divisible by $D$,
contradicting $C(4)=2$.

Suppose next $W(0)=0$, so $W=DH$ with the previously
proved mass-twelve classification. Since $H(4)=3$ and
$G_-(4)=80$, the same irreducibility/content argument
as in the preceding mass-twelve rows forces $U_+=W$.
More explicitly, a nonconstant $H$ is primitive irreducible
and cannot divide $G_-$; hence it divides $U_+$.
Its odd value $H(0)\ne\pm1$ forces $U_+(0)=0$,
and the classification of $U_+/D$ gives $U_+=W$.
When $H=3$, reduction modulo three and the mass-twelve
effectivity/parity bound give the same conclusion.
Consequently
\[
 G_+=G_-+D\psi,\qquad
                  (H-1)G_-=D(H\delta_-+\psi).
\]
The necessary NN-G030 congruence is
$G_--D\delta_-\equiv G_+\pmod4$.
It forces $R\equiv\psi\pmod2$.
The effective character $R$ has mass four, while
$\psi=T_1+T_2$ already has four odd Laurent coefficients
and mass four. Hence coefficientwise effectivity and
the mass bound force $R=\psi$ exactly. In particular
$\delta_-=3\psi$ and
\[
                    (H-1)G_-=D\psi(3H+1).
\]
Put $J=(H-1)/2\in\mathbb Z[D]$.
Then $J(4)=1$. It is either the constant one (when $H=3$)
or a monic nonconstant polynomial, since the leading
coefficient of a nonconstant $H$ is two.
The preceding identity becomes
\[
                    J(G_--3D\psi)=2D\psi.
\]
Thus the monic $J$ divides $2D^2(D-3)$ in $\mathbb Z[D]$.
Its only possibilities over $\mathbb Q[D]$ are
$D^a(D-3)^b$ with $0\le a\le2$, $0\le b\le1$;
monicity fixes the scalar. Its value one at four
excludes $a>0$. Therefore $J=1$ or $J=D-3$.

If $J=1$, the identities force
\[
 W=3D,\quad G_-=5D\psi,\quad G_+=6D\psi,\quad
 S=10D^3(D-3)^2.
\]
If $J=D-3$, they force
\[
 W=D(2D-5),\quad G_-=D^2(3D-7),\quad
 G_+=2D^2(2D-5),\quad S=2D^3(3D-7).
\]
In each case $S=G_+G_-/W$ is the actual bottom
structure coefficient. Its prescribed coefficientwise
lower bound is $4D^2(D-1)^2$, whose central Laurent
coefficient is $144$.
Using $[q^0]D^h=\binom{2h}{h}$, the two forced
central coefficients of $S$ are respectively
\[
 10(252-6\cdot70+9\cdot20)=120,\qquad
                    6\cdot70-14\cdot20=140.
\]
Both are smaller than $144$, contradicting effectivity
and the given trace domination. This excludes the reverse
row and every permitted auxiliary increment, without a
support-degree restriction.
\hfill$\square$

This is an integral-character obstruction in the same
closed kernel presentation as NN-G036. The final inequality
compares actual multiplicities of the trivial torus character;
the dictionary is SGA~3, I, 4.7.2--4.7.3 and NN-G010.
No conclusion about a different multiplication table follows.

% Final NN-G036 row exclusion and exact conclusion; root whole-proof PASS 2026-11-22T00:21:30Z.
\paragraph{The last mass-twelve row.}
The only row not covered by the preceding arguments is
\[
 w=12,\quad g_1=g_2=96,\quad s=768,\quad
 \delta_1=\delta_2=\psi,\quad u_1=u_2=10.
\]
Every effective $U_i\ge D$ of mass ten with the required
parity is $D+6$ or $D+2+2T_j$. In particular
$U_i(0)\in\{6,-2\}$, so both values are nonzero.

If $W(0)\ne0$, it is $8$ or $-8$ by the mass-twelve
classification. In the pair of cell identities
\[
 WG_1-U_1G_2=DW\psi,\qquad
 WG_2-U_2G_1=DW\psi,
\]
the coefficient determinant at zero is
$64-U_1(0)U_2(0)\ne0$, because the possible products
are $36,-12,4$. Hence both $G_i$ are divisible by $D$.
After cancelling this factor, the right sides are
$W\psi$, still divisible by $D$.
Applying the same invertible matrix at zero again
shows that both $G_i$ are divisible by $D^2$.
Since $WS=G_1G_2$ and $W(0)\ne0$, it follows that
$D^4\mid S$ integrally.
The quotient has even coefficients because $S\equiv0$
modulo two, but its value at four is
$768/4^4=3$. This is impossible.

It remains that $W=DH$, where $H(4)=3$,
$H\equiv1\pmod2$, and $H$ is either three
or primitive irreducible with leading coefficient two.
If $H$ is nonconstant, it cannot divide either
$U_i$, since three does not divide ten.
The common-factor argument of (G036.R2) gives
$H\mid G_1,G_2$ integrally.
If $H=3$, reduce both cell identities modulo three.
The reductions of the $U_i$ are nonzero, since their
masses are ten; the integral domain $\mathbb F_3[D]$
therefore forces the reductions of both $G_i$ to
vanish. Thus $H\mid G_1,G_2$ also in this case.
Write $G_i=HP_i$. Then
\[
 P_i(4)=32,\qquad
 P_1\equiv0,\quad P_2\equiv D^2(D-1)\pmod2,
\]
and the divided cell identities are
\[
 DH P_1-U_1P_2=D^2\psi,\qquad
 DH P_2-U_2P_1=D^2\psi.
\]
Because each $U_i(0)\ne0$, evaluation at zero
forces both $P_i(0)=0$.
Writing $P_i=DQ_i$ and cancelling one $D$ gives
\[
 DH Q_1-U_1Q_2=D\psi,\qquad
 DH Q_2-U_2Q_1=D\psi.
\]
Evaluation at zero again forces both $Q_i(0)=0$.
Consequently $P_i=D^2R_i$ with integral $R_i$.
Their masses are $R_i(4)=32/16=2$.
But $P_2\equiv D^2(D-1)$ forces
$R_2\equiv D-1\pmod2$, whose evaluation at four
is odd. This contradicts $R_2(4)=2$.
The last row is excluded for every grading support.
\hfill$\square$

\paragraph{Exact closed-table optimum and geometric conclusion.}
The complete necessary mass gate of NN-G035,
followed by all the uniform exclusions of NN-G036,
leaves exactly the unique $2744$-pair realization
of NN-G033--NN-G034. Therefore the minimum number
of cancelling trivial/sign pairs in the specified
symmetric closed-table presentation is exactly
$2744$, and its polynomial parameter realization
is unique within that presentation.
All cell degrees and contents, all grading support
locations, and all rational auxiliary parameters
were allowed throughout, subject only to integrality
and effectivity of the actual structure constants.
The final certificate projects every surviving
ordinary-dimension row; it does not enumerate
bounded grading degrees to justify the exclusions.

In the NN-G010 category this is an exact minimum
for actual effective graded sign-representation
objects composing by the fixed kernel table.
The grading is the diagonalizable-group dictionary
of SGA~3, I, 4.7.2--4.7.3, and composition is
the supplied archive cote~2, scan~17 transcription.
The theorem concerns this presentation, not all
possible effective multiplication tables.
The unrestricted extra-position optimization still
has the actual upper bound $2744$, the finite
equal-score cutoff of NN-G029/NN-G033, and the
total-mass FPT algorithm of NN-G032; its numerical
minimum remains open.


\subsection{Current gates that are not new theorems}
\paragraph{The finite-prime model gate has been repaired.}
NN-G016 supplies the cutoff, self-adjoint lift, and Morita
identification missing from the original finite-prime argument.
Its global arithmetic edge identification remains a separately
scoped background input. The elementary support valuation
calculation NN-G005 is independent of either analytic input.

\paragraph{Rank four and nonstandard purity.}
The E4 note, section 4, leaves the canonical
Gelfand--Tsetlin scalings, label order, and minimization
data unspecified at rank four. Numerical positivity
of one $\Sigma_4$ at two rational values does not
prove Laurent coefficient positivity. NSRH F03
requires a finite effective spectrum, the opposite
path sign for unbounded powers, and uniform bounds;
the dual spectrum is needed for equality of weights.
NSRH F50 constructs the sign but its pointwise and
dual bounds are equivalent to the desired purity,
so the construction does not discharge those bounds.
NN-G008--NN-G011 reduce a chosen integral basis to
an explicit associative effective-lift problem;
they do not identify that problem with the original
canonical-basis purity assertion. The primary weight
framework is BBD 5.1.8, 5.3.2, 5.3.8 and relative
hard Lefschetz 5.4.10; the Gaussian construction's
exact locators are SGA $4\frac12$, Sommes trig.,
1.7--1.9 (especially 1.9.2), Propositions 4.2--4.3,
and Remark 4.4, as stated and rechecked in the
supplied NSRH F50 proof.

\paragraph{Positive formulas with bounded parameters.}
The GCT\_PLETH intake's G-036 (outer degree five),
G-037 (stable hooks), and G-038--G-040 are retained
with their stated parameter bounds. Their original
proof files and the named \texttt{outer\_five\_positive.py}
and \texttt{stable\_hook\_positive.py} providers are absent
from the extracted supplied sources: the only matches
for those finding IDs are the intake, consult and
implications report. No original-proof reading or new
proof certification is asserted for them.
They give useful bounded-parameter families, not a
uniform PH1 theorem for arbitrary plethysm. The
remaining uniform gate must preserve witness
length, construction time, dimension, and bitlength
under the varying parameters. For Hadamard--Howe
or Foulkes maps, a positive difference of characters
does not give the rank of the specified map; a
kernel/cokernel calculation is still required.
